\begin{center}
\medskip
\noindent\textbf{\large `ವೇದಾಂಗಗಳು - ದರ್ಶನ - ಇತಿಹಾಸ - ಪುರಾಣ'}\\
\large{ಸಂಪುಟದಲ್ಲಿ ಬಳಸಿರುವ ಸಂಸ್ಕೃತೋದ್ಧರಣಗಳ ಕನ್ನಡ}\\
\large{ಭಾವಾನುವಾದ - ಅಕಾರಾದಿಯಾಗಿ}
\end{center}

ಸೂಚನೆ: ಆಯಾ ಸಂಸ್ಕೃತೋದ್ಧರಣವಿರುವ ಪುಟದಲ್ಲಿಯೇ ಕನ್ನಡ ಭಾವಾನುವಾದವಿದೆಯೆಂಬುದನ್ನು ಸೂಚಿಸಲು ಈ  $\star$ ಗುರುತನ್ನು ಕೊಟ್ಟಿದೆ. ಒಂದೇ ವಿಷಯವನ್ನು ಪ್ರತಿಪಾದಿಸಲು ಅನೇಕ ಶ್ಲೋಕಗಳಿರುವಲ್ಲಿ ಆಯಾ ಶ್ಲೋಕಪುಂಜದ ಮೊದಲನೆಯ ಶ್ಲೋಕ ಪಂಕ್ತಿಯ ಆಕಾರಾದಿ ಅಕ್ಷರಸ್ಥಾನದಲ್ಲಿ ಆ ಶ್ಲೋಕ ಪುಂಜದ ಕನ್ನಡ ಭಾವಾನುವಾದವನ್ನು ನೀಡಿದೆ.

\medskip
\medskip
\noindent\textbf{೧. ಅಂಕೇನೋದೂಹ್ಯ ವಾಗ್ದೇವೀಂ} \hfill ಪುಟ \pageref{32}

\hfill ಹಯಗ್ರೀವ ಸಹಸ್ರನಾಮ ಶ್ಲೋಕ - ೩೧

ಹಯಗ್ರೀವನು ವಾಗ್ದೇವಿಯಾದ್ ಭಾರತಿಯನ್ನು ತನ್ನ ಮಡಿಲಲ್ಲಿ ಕುಳ್ಳಿರಿಸಿಕೊಂಡು ಆಚಾರ್ಯಭಾವವನ್ನು ಹೊಂದಿ ವೇದ, ವೇದಾಂತ ಹಾಗೂ ಶಾಸ್ತ್ರಾರ್ಥಗಳ ತಾತ್ತ್ವಿಕಾರ್ಥವನ್ನು ವ್ಯಾಖ್ಯಾನಿಸುವುದರಲ್ಲಿ ತತ್ಪರನಾಗಿದ್ದಾನೆ.

\medskip
\noindent\textbf{೨. ಅಂಕೋನ್ಮುಕ್ತ ಶಶಾಂಕಕೋಟಿಸದೃಶೀಂ} $\star$ \hfill ಪುಟ \pageref{31}

\hfill ಶಾರದಾತಿಲಕ ತಂತ್ರ ೬-೧೫

\medskip
\noindent\textbf{೩. ಅಂತರೇಣ ತಾಲುಕೇ| ಯ ಏಷಸ್ತನ ಇವಾವಲಂಬತೇ| ಸೇಂದ್ರಯೋನಿಃ|} \hfill ಪುಟ \pageref{26}

\hfill ತೈತ್ತಿರೀಯೋಪನಿಷತ್ ೧-೫

ಎರಡೂ ದವಡೆಗಳ ನಡುವೆ ಹಸುವಿನ ಮೊಲೆಯಂತೆ ಜೋಲುವ ಯಾವ ಅಡಿ ನಾಲಗೆಯುಂಟೋ ಅದೇ ಪರಮೈಶ್ವರ್ಯವುಳ್ಳ ಭಗವಂತನೊಡನೆ ಸೇರಿಕೊಳ್ಳಲು ದಾರಿಯಾಗಿದೆ. ತಲೆಗೂದಲು ಉಗಮಗೊಳ್ಳುವ ಜಾಗವದು. ಸುಷುಮ್ನೆಯು ಬ್ರಹ್ಮರಂಧ್ರ ಸೇರುವ ನಾಡೀಮಾರ್ಗವಾಗಿದೆ.

\medskip
\noindent\textbf{೪. ಅಂತರ್ಲಕ್ಷ್ಯಂ ಬಹಿರ್ದೃಷ್ಟಿ:} \hfill ಪುಟ \pageref{244}

\hfill ಶಾಂಡಿಲ್ಯೋಪನಿಷತ್ ೧-೧೪

ಶಾಂಭವೀ ಅಥವಾ ವೈಷ್ಣವೀ ಎಂಬುದು ಯೋಗಿವರೇಣ್ಯರಲ್ಲಿ ಗೋಚರಿಸುವ ಒಂದು ವಿಶೇಷವಾದ ಮುದ್ರೆ. ಮುದಂ ರಾತಿ-ಆದತ್ತೇ ಇತಿ ಮುದ್ರಾ. ಪರಮೋನ್ನತವಾದ ಆನಂದಾನುಭವದಲ್ಲಿ ಗೋಚರಿಸುವ ಸ್ಥಿತಿವಿಶೇಷ. ನಯನಗಳು ತೆರೆದಿದ್ದರೂ ದೃಷ್ಟಿಯು ಅಂತರಂಗದ ತತ್ತ್ವಭೂಮಿಕೆಯ ಲಕ್ಷ್ಯದಲ್ಲಿ ನಾಟಿ ನಿಮೇಷೋನ್ಮೇಷ ರಹಿತವಾಗಿರುವುದು. ಈ ಮುದ್ರೆಯು ವೇದ - ಶಾಸ್ತ್ರಗಳಲ್ಲಿ ಗುಪ್ತವಾಗಿಡಲ್ಪಟ್ಟಿದೆ - ಎಂದರೆ ವೇದ-ಶಾಸ್ತ್ರಗಳ ಸಾರ ಸರ್ವಸ್ವವಾಗಿದೆ.

\medskip
\noindent\textbf{೫. ಅಕರ್ದಮಮಿದಂ ತೀರ್ಥಂ } $\star$ \hfill ಪುಟ \pageref{156},\pageref{164},\pageref{175}

\hfill ರಾಮಾಯಣ ಬಾಲಕಾಂಡ ೨-೫

\medskip
\noindent\textbf{೬. ಅಕ್ಷಾರಾಣಾಮಕಾರೋಽಸ್ಮಿ } \hfill ಪುಟ \pageref{93},\pageref{95},\pageref{97}

\hfill ಭಗವದ್ಗೀತಾ ೧೦-೨೨

ಆಕಾರ ಮೊದಲುಗೊಂಡು ಕ್ಷಕಾರದವರೆಗಿನ ಎಲ್ಲಾ ಅಕ್ಷರಗಳಿಗೂ ಮೊದಲನೆಯದಾಗಿ ಸರ್ವವರ್ಣ ಪ್ರಧಾನವಾಗಿರುವ ಅಕಾರ ರೂಪನಾಗಿದ್ದೇನೆ. ಅರ್ಜುನ. 

\medskip
\noindent\textbf{೭. ಅಗಾಯತಾಂ ಮಾರ್ಗ ವಿಧಾನಸಂಪದಾ} \hfill ಪುಟ \pageref{188}

\hfill ರಾಮಾಯಣ ಬಾಲಕಾಂಡ ೪-೩೬

ಕುಶಲವರು ರಾಮಾಯಣ ಕಾವ್ಯವನ್ನು ಆತ್ಮರಂಜನಪರವಾದ ಮಾರ್ಗ ವಿಧಾನವೆಂಬ ಗಾನಪದ್ಧತಿಯಲ್ಲಿ ವೈಭವೋಪೇತವಾಗಿ ಹಾಡಿದರು.

\medskip
\noindent\textbf{೮. ಅಚ್ಛೇದ್ಯೋಽಯಂ ಅದಾಹ್ಯೋಽಯಂ } \hfill ಪುಟ \pageref{36}

\hfill ಭಗವದ್ಗೀತಾ ೨-೨೪

ಈ ದೇಹಿಯಾದ ಆತ್ಮನು ಛೇದನದಿಂದಾಗಲಿ ದಹನದಿಂದಾಗಲಿ ವಿರೂಪಗೊಳಿಸಲು ಸಾಧ್ಯನಲ್ಲ.

\medskip
\noindent\textbf{೯. ಅಜಾಯಮಾನೋ ಬಹುಧಾ} $\star$ \hfill ಪುಟ\pageref{150}

\hfill ಪುರುಷ ಸೂಕ್ತ

\medskip
\noindent\textbf{೧೦. ಅಜೀರ್ಣೇ ಭೋಜನಂ ವಿಷಂ } \hfill ಪುಟ \pageref{265}

\hfill ಅಜೀರ್ಣವಿದ್ದಾಗ ಭೋಜನವು ವಿಷಪರಿಣಾಮವನ್ನುಂಟುಮಾಡುತ್ತದೆ.

\medskip
\noindent\textbf{೧೧. ಅನುಭೂತವಿಷಯಾಸಂಪ್ರಮೋಷಃ ಸ್ಮೃತಿಃ} \hfill ಪುಟ \pageref{81}

\hfill ಪಾತಂಜಲ ಯೋಗಸೂತ್ರ ೧-೧೧

ಇಂದ್ರಿಯಗಳ ಮೂಲಕ ಅನುಭವಗೋಚರವಾದ ಯಾವುದೇ ವಿಷಯವು ಮನಸ್ಸಿನಿಂದ ಕಳೆದುಹೋಗದೇ ಉಳಿದರೆ ಅದಕ್ಕೆ ಸ್ಮೃತಿಯೆಂದು ಹೆಸರು.

\medskip
\noindent\textbf{೧೨. ಅನ್ನಂ ನ ನಿಂದ್ಯಾತ್} \hfill ಪುಟ \pageref{76}

\hfill ತೈತ್ತಿರೀಯೋಪನಿಷತ್ ೩-೭

ಅನ್ನವನ್ನು ಅನಾದರಿಸಬಾರದು. ಅದು ವ್ರತರೂಪವಾಗಿ ರಕ್ಷಿಸತಕ್ಕದ್ದು. ಅನ್ನವು ಆದರಣೀಯವಾದ್ದರಿಂದಲೇ ಯಾವುದಾದರೂ ಮಾರ್ಗವನ್ನನುಸರಿಸಿ ಬಹುವಾದ ಅನ್ನವನ್ನು ಸಂಪಾದಿಸಬೇಕು.

\medskip
\noindent\textbf{೧೩. ಅನ್ವಗಾದಿವ ಹಿ ಸ್ವರ್ಗಃ} \hfill ಪುಟ \pageref{252}

ಭುವಿಗೆ ಇಳಿದ ಪುರುಷೋತ್ತ್ಮನನ್ನು ಹಿಂಬಾಲಿಸಿ ಸ್ವರ್ಗವೂ ಭುವಿಗಿಳಿದಿದೆಯೋ ಎಂಬಂತೆ ರಾಮರಾಜ್ಯವು ಕಾಣುತ್ತಿತ್ತು.

\medskip
\noindent\textbf{೧೪. ಅಬುದ್ಧ್ವಾ ಮಾ ಲಿಖ ಕ್ವಚಿತ್ } $\star$ \hfill ಪುಟ \pageref{119}

\hfill ಮಹಾಭಾರತ ಆದಿಪರ್ವ ೧-೭೯

\medskip
\noindent\textbf{೧೫. ಅಭ್ಯಷಿಂಚನ್ ನರವ್ಯಾಘ್ರಂ} \hfill ಪುಟ \pageref{150},\pageref{208},\pageref{248}

\hfill ರಾಮಾಯಣ ಯುದ್ಧಕಾಂಡ ೧೩೧-೬೧

ಅಷ್ಟವಸ್ತುಗಳೂ ಸೇರಿ ಸಹಸ್ರಾಕ್ಷನಾದ ದೇವರಾಜನನ್ನು ಹೇಗೋ ಹಾಗೆ ವಸಿಷ್ಠ ವಾಮದೇವಾದಿ ಮುನಿಗಳೆಲ್ಲ ಒಂದು ಗೂಡಿ ಸಹಸ್ರಾಕ್ಷ ಪುರುಷೋತ್ತಮನೇ ನರೋತ್ತಮನಾಗಿಳಿದಿರುವ ರಾಮನನ್ನು ಕೇಶಾದಿ ಪಾದಾಂತವಾಗಿ ಪ್ರಸನ್ನತೆಯನ್ನೆರ್ಪಡಿಸುವಂಥವೂ, ಸುಗಂಧ (ಆತ್ಮಗಂಧ) ಯುಕ್ತವೂ ಆದ (ಸರಳವಾಗಿ ಹರಿದು ಬರುವ) ಸಲಿಲದಿಂದ ಸೇಚನಪೂರ್ವಕವಾಗಿ ಪಟ್ಟಾಭಿಷೇಕ ಮಾಡಿದರು (ಯೋಗ ಪಟ್ಟಾಭಿರಾಮನಾಗಿ ಕಂಗೊಳಿಸುವಂತೆ ಮಾಡಿದರು)

\medskip
\noindent\textbf{೧೬. ಅರ್ಥಮನುರ್ಥಂ ಭಾವಯ ನಿತ್ಯಂ} \hfill ಪುಟ \pageref{55}

\hfill ಭಜಗೋವಿಂದಂ ಸ್ತೋತ್ರ

ಹಣವು ಅನರ್ಥಕಾರಿಯಾದುದೆಂದು ಯಾವಾಗಲೂ ಭಾವಿಸು.

\medskip
\noindent\textbf{೧೭. ಅವಜಾನಂತಿ ಮಾಂ ಮೂಢಾಃ} \hfill ಪುಟ \pageref{139}

\hfill ಭಗವದ್ಗೀತಾ ೯-೧೧

ನಿಸರ್ಗದಲ್ಲಿರುವ ಸ್ಥೂಲ-ಸೂಕ್ಷ್ಮ-ಪರಭಾವಗಳ ಪೈಕಿ ಸ್ಥೂಲಭಾವವನ್ನು ಮಾತ್ರ ಅರಿಯುವಂತಹ ಮೋಹಾವೃತಜೀವಿಗಳು ಭೂತಸೃಷ್ಟಿಯ ಹಿಂಬದಿಯಲ್ಲಿರುವ, ತತ್ತ್ವ ವಿಕಾಸಕ್ಕೂ ಮೂಲವಾದ, ಮಹತ್ತತ್ತ್ವದ ಮೇಲೂ ಪ್ರಭುತ್ವವುಳ್ಳ ನನ್ನ ಪರಭಾವವನ್ನು, ಪರಾತ್ಪರ ಭಾವವನ್ನು ಅರಿಯಲಾಗದವರಾಗಿ ಮಾನುಷ ಶರೀರಿಯಾಗಿ ಕಾಣುವ ನನ್ನನ್ನು  ಮನುಷ್ಯನೆಂದೇ ಭಾವಿಸಿ ಅವಜ್ಞಾಬುದ್ಧಿಯಿಂದ ನೋಡುವರು. ಅರ್ಜುನ!

\medskip
\noindent\textbf{೧೮. ಅವಿಚಿನ್ವನ್ನಿಮಂ ವೇದೇ}  $\star$ \hfill ಪುಟ \pageref{34}

\hfill ಮಹಾಭಾರತ ಉದ್ಯೋಗಪರ್ವ ೪೩-೫೯ ರಿಂದ ೬೩

\medskip
\noindent\textbf{೧೯. ಅಶ್ವತ್ಥಃ ಸರ್ವವೃಕ್ಷಾಣಾಮ್} \hfill ಪುಟ \pageref{53},\pageref{98}

\hfill ಭಗವದ್ಗೀತಾ ೧೦-೨೬

ಅರ್ಜುನ, ವಿಶ್ವದಲ್ಲಿ ಬೆಳೆಯುವ ಎಲ್ಲ ವೃಕ್ಷಗಳ ಪೈಕಿ ತ್ರಿಮೂರ್ತ್ಯತ್ಮಕವಾದ ಅಶ್ವತ್ಥವೃಕ್ಷವಾವುದುಂಟೋ ಅದು ನನ್ನ ಉತ್ಕೃಷ್ಟ ಮೂರ್ತಿಗಳಲ್ಲೊಂದಾಗಿದೆ.

\medskip
\noindent\textbf{೨೦. ಅಸನ್ನೇವ ಸ ಭವತಿ} \hfill ಪುಟ \pageref{24}

\hfill ತ್ತೈತ್ತಿರಿಯೋಪನಿಷತ್ ೬-೧

\medskip
\noindent\textbf{೨೧. ಅಸ್ಯ ಕ್ರೂರೈಃ ನೃಶಂಸೈಶ್ಚ} \hfill ಪುಟ \pageref{192}

\hfill ರಾಮಾಯಣ ಸುಂದರಕಾಂಡ ೪೯-೧೯

ಈ ರಾವಣನ ಕ್ರೂರವೂ, ನಿರ್ದಯವೂ, ಲೋಕನಿಂದಿತವೂ ಆದಂತಹ ವರ್ತನೆಗಳಿಂದ ಬೇಸರ ಹುಟ್ಟಿದ ಕಾರಣ, ದೇವಾಸುರಮನುಷ್ಯರಿಂದ ಕೂಡಿದ ಜನರು ಹೆದರುತ್ತಿರುವರು.

ಇವನು ಕೋಪಗೊಂಡಿದ್ದೇ ಆದರೆ ಜಗತ್ತೆಲ್ಲವೂ ಸಮುದ್ರಮಯವಾಗುವಂತೆ ಮಾಡಲೂ ಶಕ್ತನಾದಾನು.

\medskip
\noindent\textbf{೨೨. ಆಸ್ತಿ ಬ್ರಹ್ಮೇತಿ ಚೇದ್ವೇದ ಸಂತಮೇನಂ ತತೋ ವಿದುರಿತಿ|} \hfill ಪುಟ \pageref{24}

\hfill ತೈತ್ತಿರೀಯೋಪನಿಷತ್ ೬-೧

ಬ್ರಹ್ಮತತ್ತ್ವದ ಅಸ್ತಿತ್ವವನ್ನು ಸಾಕ್ಷತ್ಕರಿಸಿದ್ದೇ ಆದರೆ, ಅಂಥವನನ್ನು ಸತ್ಪುರುಷನೆನ್ನುವರು.

\medskip
\noindent\textbf{೨೩. ಅಹಂ ವೇದ್ಮಿ ಮಹಾತ್ಮಾನಂ} \hfill ಪುಟ \pageref{195},\pageref{251}

\hfill ರಾಮಾಯಾಣ ಬಾಲಕಾಂಡ ೧೯-೧೪

ಮಹಾತ್ಮನೂ, ಸತ್ಯರಕ್ಷಕವಾದ ಪರಾಕ್ರಮವುಳ್ಳವನೂ ಆದ ಈ ರಾಮನನ್ನು ನಾನು (ವಿಶ್ವಾಮಿತ್ರನು) ಯಾರೆಂದು ತಾತ್ತ್ವಿಕವಾಗಿ ಅರಿತಿರುವೆನು. ಮಹಾತೇಜಸ್ವಿಯಾದ ವಸಿಷ್ಠಮುನಿಯೂ ಅರಿತಿರುವನು. ಅಲ್ಲದೇ ತಪಸ್ಸಿನಲ್ಲಿ ನೆಲೆಸಿ ತತ್ತ್ವಚಿಂತನೆ ಮಾಡಬಲ್ಲ ಇವರಾರುಂಟೋ ಅವರು ಅರಿಯಬಲ್ಲರು. 

\medskip
\noindent\textbf{೨೪. ಅಹೋ ಗೀತಸ್ಯ ಮಾಧುರ್ಯಂ} \hfill ಪುಟ \pageref{159}

\hfill ರಾಮಾಯಣ ಬಾಲಕಾಂಡ ೪-೧೭

`ಚಿರನಿರ್ವೃತ್ತ' ಇತ್ಯಾದಿ ಗ್ರಂಥದ ಭಾವಾನುವಾದವನ್ನು ಗಮನಿಸಿ

\medskip
\noindent\textbf{೨೫. ಆಕಾರ ಸದೃಶಪ್ರಜ್ಞಃ} \hfill ಪುಟ \pageref{139}

\hfill ರಘುವಂಶ ೧-೧೫

ಆಕಾರಕ್ಕೆ ಅನುರೂಪವಾದ ಪ್ರಜ್ಞಾಸಂಪತ್ತನ್ನು ಉಳ್ಳವನು (ದಿಲೀಪ ಮಹಾರಾಜ)

\medskip
\noindent\textbf{೨೬. ಆಜಗಾಮ ತತೋ ಬ್ರಹ್ಮಾ} \hfill ಪುಟ \pageref{202}

\hfill ರಾಮಾಯಣ ಬಾಲಕಾಂಡ ೨-೨೩

ಆ ಕೂಡಲೇ ಲೋಕವನ್ನೆಲ್ಲಾ ನಿರ್ಮಿಸುವವನೂ, ಲೋಕದ ಮೇಲಿನ ಪ್ರಭುತ್ವವುಳ್ಳವನೂ, ನಾಲ್ಕು ಮುಖಗಳನ್ನು ಹೊಂದಿರುವವನೂ, ಮಹಾತೇಜಸ್ವಿಯೂ ಆದ ಬ್ರಹ್ಮನು ಆ ವಾಲ್ಮೀಕಿಮುನಿಯನ್ನು ಕಾಣಲು ಸಾಕ್ಷಾತ್ತಾಗಿ ತಾನೇ ಅಲ್ಲಿಗೆ ಬಂದನು.

\medskip
\noindent\textbf{೨೭. ಆಜಾನುಬಾಹುಃ ಸುಶಿರಾಃ} \hfill ಪುಟ \pageref{247}

\hfill ರಾಮಾಯಣ ಬಾಲಕಾಂಡ ೧-೧೦

`ಇಕ್ಷ್ವಾಕುವಂಶಪ್ರಭವೋ' ಇತ್ಯಾದಿ ಗ್ರಂಥದ ಭಾವಾನುವಾದವನ್ನು ಗಮನಿಸಿ.

\medskip
\noindent\textbf{೨೮. ಆತ್ಮಕ್ರೀಡಃ ಆತ್ಮರತಿಃ} \hfill ಪುಟ \pageref{118},\pageref{143}

\hfill ಮುಂಡಕೋಪನಿಷತ್ ೩-೧-೪

ಆತ್ಮನಲ್ಲಿ ಪರ್ಯವಸಾನಗೊಳ್ಳುವಂತಹ ಕ್ರೀಡೆಯುಳ್ಳವನೂ, ಆತ್ಮನಲ್ಲಿ ಮುಗಿಯುವಂತಹ ರತಿಭಾವವುಳ್ಳವನೂ ಆಗುವ ತೆರದಲ್ಲಿ ಜೀವನದ ಎಲ್ಲಾಕ್ರಿಯೆಗಳಲ್ಲೂ ತೊಡಗುವವನು ಯಾರುಂಟೋ ಅವನೇ ಬ್ರಹ್ಮಜ್ಞರಲ್ಲಿ ಶ್ರೇಷ್ಠನಾಗುಳಿಯುವನು.

\medskip
\noindent\textbf{೨೯. ಆತ್ಮಾ ಬುದ್ಧ್ಯಾ ಸಮೇತ್ಯಾರ್ಥಾನ್} $\star$ \hfill ಪುಟ \pageref{18}

\hfill ಪಾಣಿನೀಯ ಶಿಕ್ಷಾ ೧-೫

\medskip
\noindent\textbf{೩೦. ಆತ್ಮಾ ವೈ ಪುತ್ರನಾಮಾಸಿ} \hfill ಪುಟ \pageref{228}

ಆತ್ಮಾಸ್ವರೂಪಿಯಾದ ನೀನೇ ಪುತ್ರರೂಪದಲ್ಲಿದ್ದೀಯೆ.

\medskip
\noindent\textbf{೩೧. ಆಯುಷ್ಯಂ ಪುಷ್ಟಿಜನಕಂ} \hfill ಪುಟ \pageref{210}

\hfill ರಾಮಾಯಣ ಬಾಲಕಾಂಡ ೪-೨೮

ಲವಕುಶರಿಂದ ಗಾನಮಾಡಲ್ಪಟ್ಟ ರಾಮಾಯಣವು ಆಯುಷ್ಯ ಪ್ರದವಾದದ್ದೂ ಪುಷ್ಟಿಯನ್ನುಂಟು ಮಾಡುವಂತಹುದೂ ಆಗಿತ್ತು.

\medskip
\noindent\textbf{೩೨. ಆಯುಷ್ಯಮಪರೇಪ್ರೋಚುಃ} \hfill ಪುಟ \pageref{168}

\hfill ರಾಮಾಯಣ ಬಾಲಕಾಂಡ ೪-೨೫

(ರಾಮಾಯಣಗಾನವನ್ನು ಆಲಿಸಿದ ಋಷಿಮುನಿಗಳು ಕಲಶವಲ್ಕಲ ಮೊದಲಾದ ಉಡುಗೊರೆಗಳನ್ನು ಲವಕುಶರಿಗೆ ಅನುಗ್ರಹಿಸಿದರು) ಮತ್ತೆ ಕೆಲವರು ಅವರಿಗೆ ದೀರ್ಘಾಯುಷ್ಯವನ್ನು ಹರಿಸಿದರು.

\medskip
\noindent\textbf{೩೩. ಆಹಾರಶುದ್ಧೌ ಸ್ತ್ವಶುದ್ಧಿಃ} \hfill ಪುಟ \pageref{76}

\hfill ಛಾಂದೋಗ್ಯ ೭-೨೬-೨

ಜ್ಞಾನೇಂದ್ರಿಯ - ಕರ್ಮೇಂದ್ರಿಯಗಳ ಮೂಲಕ ಜೀವಿಯು ಆಹರಣ ಮಾಡಿಕೊಳ್ಳತಕ್ಕ-ತನ್ನೆಡೆಗೆ ಬರಮಾಡಿಕೊಳ್ಳತಕ್ಕ ವಿಷಯಗಳು ನಿರ್ದುಷ್ಟವಾಗಿದ್ದಲ್ಲಿ ಸತ್ವಗುಣವು ರಜಸ್ತಮೋಗಳಿಂದ ಮೆಟ್ಟಲ್ಪಡದೆ ತನ್ನದೇ ಪ್ರಭುತ್ವಹೊಂದಿ ಶುದ್ಧವಾಗಿರುವುದು. ಆ ರೀತಿ ಶುದ್ಧವಾದ ಸತ್ವವು ದೇಹಿಯಲ್ಲೊದಗಿದಾಗ ಪರಮಾತ್ಮಾನುಭವದ ಸ್ಮೃತಿಯು ಅಚಲವಾಗಿ ನಿಲ್ಲುವುದು.

\medskip
\noindent\textbf{೩೪. ಇಂದ್ರ  ಶತ್ರುರ್ವರ್ಧಸ್ವ} $\star$ \hfill ಪುಟ \pageref{9}

\hfill ತೈತ್ತಿರೀಯ ಸಂಹಿತಾ ೨-೫-೩-೧

\medskip
\noindent\textbf{೩೫. ಇಕ್ಷ್ವಾಕುವಂಶಪ್ರಭುವೋ} \hfill ಪುಟ \pageref{154},\pageref{163},\pageref{196}

\hfill ರಾಮಾಯಣ ಬಾಲಕಂಡ ೧-೮ ರಿಂದ ೧೮

ಇಕ್ಷ್ವಾಕು ಮಹಾರಾಜನ ವಂಶದಲ್ಲಿ ಜನಿಸಿದ ಇವನು ರಾಮನೆಂದು ಪ್ರಸಿದ್ಧವಾಗಿ ಜನರೆಲ್ಲರ ಕಿವಿಗೂ ಪರಿಚಿತನು. ಮನೋನಿಯಮ, ಹಿರಿದಾದ ವೀರ್ಯ, ಶತ್ರುವನ್ನೂ ಬಿಡದೇ ಎಲ್ಲರನ್ನೂ ತನ್ನೆಡೆಗೆಳೆಯುವ ಅಂಗಶೋಭೆ, ಪ್ರತಿಜ್ಞೆಯಿಂದ ಚ್ಯುತನಾಗದಿರಲು ಬೇಕಾದ ಧೃತಿ, ದೇಹೇಂದ್ರಿಯಗಳ ಸ್ವಚ್ಛಂದತೆಗೆಡೆಗೊಡದೆ ಆತ್ಮಚ್ಛಂದವಾಗಿರುವ ಸಾಮಾರ್ಥ್ಯ, ಸೂಕ್ಷ್ಮಾರ್ಥವನ್ನೂ ನಿಶ್ಚಯಿಸಬಲ್ಲ ಕುಶಾಗ್ರಮತಿ, ಎಂಥ ಸಂಕಟದಿಂದಲೂ ಪಾರಾಗಿ ಮುಂದೆ ಸಾಗಬಲ್ಲ ನೀತಿ, ಸತ್ಯಾರ್ಥವನ್ನು ಸ್ಪಷ್ಟಪಡಿಸಬಲ್ಲ ವಾಕ್ಸಂಪತ್ತಿ, ಎಲ್ಲರೂ ಬಂದು ತನ್ನನ್ನಾಶ್ರಯಿಸಲು ಬೇಕಾದ ಧರ್ಮಸಂಪತ್ತು, ಜ್ಞಾನಸಂಪತ್ತು, ದೇಹಸಂಪತ್ತು, ಎಲ್ಲದರ ಹೊಂದಣಿಕೆ, ಶತ್ರುಗಳನ್ನು ನಿರ್ಮೂಲ ಮಾಡಬಲ್ಲ ಛಲ, ಮಹಾಪುರುಷ ಲಕ್ಷಣವಾದ ಎತ್ತಿದ ಭುಜ, ಸುಲಕ್ಷ್ಣವಾದ ಬಾಹುಗಳು, ಶಂಖಾಕೃತಿಯಾದ ಕತ್ತು, ಮಾಂಸಲವಾದ ಕಪೋಲಗಳ ಮೇಲ್ಭಾಗ, ವಿಶಾಲವಾದ ವಕ್ಷಃ ಫಲಕ, ಇವುಗಳಿಂದಾಗಿ ದೊರೆತ ಮಹಾಧನುರ್ಧಾರಿತ್ವ, ಶತ್ರುಗಳ ಸೊಕ್ಕಡಗಿಸುವ ಬಲ, ಮಾಂಸಲತೆಯಿಂದ ಗೂಢವಾಗಿರುವ ಕತ್ತಿನ ಎಲುಬುಗಳು, ಮೊಳಕಾಲಿನವರೆವಿಗೂ ನೀಡಿದ ತೋಳ್ಗಳು. ಈ ಎಲ್ಲಾ ಅಂಗಸಂಪತ್ತಿಗೆ ಹೊಂದುವ ಅಂದವಾದ ಉತ್ತಮಾಂಗ, ಒಳ್ಳೆಯ ಫಲಕದಂತಿರುವ ಹಣೆ, ಒಳ್ಳೆಯ ಗಂಭೀರವಾದ ಕಾಲ್ನಡಿಗೆ, ದೇಹದ ಉದ್ದ ದಪ್ಪಗಳು ಹೊಂದಿದ ಸುಲಕ್ಷಣವಾದ ಆಯಕಟ್ಟಾದ ದೇಹ, ಕೇಶಾದಿ ಪಾದಾಂತವಾದ ಎಲ್ಲಾ ಅಂಗಗಳ ಸುಬಿಭಕ್ತವೂ ಸರಿಸಮವೂ ಆದ ರೂಪ, ಅತಿ ಬಿಳುಪೂ ಅತಿಕಪ್ಪೂ ಆಗದೇ ಸ್ನಿಗ್ಧಶ್ಯಾಮಲವಾದ ಮೈ ಬಣ್ಣ, ದರ್ಶನ ಮಾತ್ರದಿಂದಲೇ ಎಲ್ಲರ ಮನಸ್ಸಿನ ಮೇಲೂ ಪ್ರಭಾವ ಬೀರಬಲ್ಲ ಭೀಮಕಾಂತವಾದ ಗಾಂಭೀರ್ಯ, ತುಂಬಿದ ಎದೆ, ಸುಲಕ್ಷಣವಾಗಿ ವಿಶಾಲವಾದ ಕಣ್ಗಳು, ಎಲ್ಲವೂ ಲಕ್ಷೀರೂಪವಾಗಿರಬಲ್ಲ ಸೊಬಗು, ಯಾವುದೇ ಅತಿಲಕ್ಷ್ಣವಿಲ್ಲದೇ ಅಂತರಂಗ ಬಹಿರಂಗಗಳ ಶುಭರೂಪನನ್ನೇ ಗುರುತಿಸಲು ಸಾಧನವಾದ ಹೊರಲಕ್ಷಣ, ಧರ್ಮತತ್ತ್ವದ ದರ್ಶನಾನುಭವಗಳು, ಸತ್ಯದೊಡನೇ ದೂರವಾಗದೇ ಯಾವಾಗಲೂ ಲಕ್ಷ್ಯವಿರುವುದು ಅಥವಾ, ಸತ್ಯಪ್ರತಿಜ್ಞತೆ, ಪ್ರಜಾಹಿತ ಪರಾಯಣತೆ, ಎಲ್ಲರೂ ಮುಕ್ತಕಂಠವಾಗಿ ಹೊಗಳತಕ್ಕ ಕೀರ್ತಿಶೀಲತೆ, ಮಹಾಯೋಗಿಗಳು ತಲುಪಿರುವಂತಹ ಜ್ಞಾನಸಂಪತ್ತು, ನೈರ್ಮಲ್ಯ, ತನ್ನವರನ್ನೂ ಪರರನ್ನೂ ಸ್ವಾಧೀನದಲ್ಲಿಟ್ಟುಕೊಳ್ಳುವ ನಿಶ್ಚಿತ ಪ್ರವೃತ್ತಿ, ಸದಾ ಸಮಾಹಿತಚಿತ್ತತೆ, ಸೃಷ್ಟಿಕರ್ತನಾದ ಪ್ರಜಾಪತಿಯಂತೆ ಎಲ್ಲರಿಗೂ ತಂದೆಯಾಗಿರತಕ್ಕ ನಿಷ್ಪಕ್ಷಪಾತತೆ, ಆತ್ಮಶ್ರೀ, ಎಲ್ಲರನ್ನೂ ಪೋಷಿಸಿ ಭರಿಸುವ ಔದಾರ್ಯ, ತನ್ನ ಗುರಿ ಸಾಧನೆಗೆ ಅಡ್ಡಿ ಆತಂಕಗೊಳೊದಗದಂತೆ ಮಾಡಿಕೊಳ್ಳುವ ಶಕ್ತಿ, ಜೀವಲೋಕವನ್ನೇಲ್ಲಾ ವೈಷಮ್ಯವಿಲ್ಲದೇ ರಕ್ಷಿಸುವ ಅಸಂಕುಚಿತ ಸ್ವಭಾವ, ಎಲ್ಲರೂ ಎಲ್ಲವೂ ಅವರವರ ಧರ್ಮದೊಡನೇ ಇರುವಂತೆ ರಕ್ಷಿಸುವ ಗುರಿ, ಸರ್ವಧರ್ಮರಕ್ಷಣೆಯಲ್ಲಿ ಹೊರಟು ತನ್ನ ಧರ್ಮವನ್ನು ಮರೆತು ಬಿಡದೇ ಇಟ್ಟುಕೊಳ್ಳುವುದು, ಪರರನ್ನು ರಕ್ಷಿಸುವಂತೆಯೇ ಸ್ವಜನರನ್ನೂ ರಕ್ಷಿಸಿಕೊಳ್ಳುವುದು, ವೇದ ಮತ್ತು ವೇದಾಂಗಗಳ  ತತ್ತ್ವಾರ್ಥದ ಬಗೆಗೆ ನೈಜವಾದ ಅರಿವು, ಧನುರ್ವೇದದ ನೆಲೆಯ ಅರಿವು, ಚತುರ್ದಶ ವಿದ್ಯಾಸ್ಥಾನಗಳು ಅಥವಾ ಅಷ್ಟಾದಶ ವಿದ್ಯಾಸ್ಠಾನಗಳು ಚತುಷ್ಟಷ್ಠಿಕಲೆಗಳ  ಮರ್ಮಜ್ಞಾನ, ಎಲ್ಲೂ ಪೂರ್ವಸಂಸ್ಕೃತಿಯನ್ನು ಕಳೆದುಕೊಳ್ಳದಿರುವುದು, ಸಮಯೋಚಿತವಾದ ತತ್ತ್ವಾರ್ಥಗಳ ನವನವಸ್ಫೂರ್ತಿ, ಭೂರ್ಭುವಸ್ಸುವರ್ಲೋಕಗಳಲ್ಲಿ ಯಾವಾಗಲೂ ಪ್ರೀತಿ, ತಪ್ಪು ನಡೆಗೆ ಎಡೆಗೊಡದೇ ಅಂತರ್ಬಾಹ್ಯ ಸದಾಚಾರ ಸಂಪತ್ತಿ. ಯಾವುದೇ ಕಾರಣದಿಂದ ದೈನ್ಯವನ್ನವಲಂಬಿಸದೇ ಇರುವ ಆತ್ಮವಂತಿಕೆ ನಿಜವನ್ನು ಅರಿತು ತಿಳಿಸಬಲ್ಲ ಸಾಮರ್ಥ್ಯ, ಸಮುದ್ರವು ನದಿಗಳಿಂದ ತುಂಬಿರುವಂತೆ ಜ್ಞಾನನದಿಗಳಂತಿರುವ ಸಜ್ಜನರಿಗೆ ಆತ್ಮಸಾಗರನಾಗಿರುವಿಕೆ, ದೂರದರ್ಶನ ಸಮರ್ಥವಾದ ಬುದ್ಧಿ, ಯಾವಾಗಲೂ ಎಲ್ಲರಿಗೂ ಪ್ರಿಯನಾಗಿಯೇ ಕಾಣುತ್ತಿರುವ ಮಹಾನುಭಾವತೆ, ಇಷ್ಟೆಲ್ಲಾ ಸಮಸ್ತಗುಣಪೂರ್ಣನಾದ ಕೌಸಲ್ಯಾನಂದವರ್ಧನನಾದ ಶ್ರೀರಾಮನು ಗಾಂಭೀರ್ಯದಲ್ಲಿ ಸಮುದ್ರದಂತೆಯೂ, ಧೈರ್ಯದಲ್ಲಿ ಹಿಮಾಚಲನಂತೆಯೂ, ವೀರ್ಯದಲ್ಲಿ ವಿಷ್ಣುವಿನಂತೆಯೂ, ಚಂದ್ರನಂತೆ ಸದಾಹ್ಲಾದಕರವಾದ ರೂಪವೂ ಉಳ್ಳವನಾಗಿದ್ದನು.

\medskip
\noindent\textbf{೩೬. ಇದಂ ಪವಿತ್ರಂ ಪಾಪಘ್ನಂ} \hfill ಪುಟ \pageref{181},\pageref{211}

\hfill ರಾಮಾಯಣ ಬಾಲಕಂಡ ೧-೧೮

ಈ ರಾಮಚರಿತ್ರವು ಪವಿತ್ರವಾದುದರಿಂದ ಪಾಪಗಳಿಂದ ಬಿಡುಗಡೆಯನ್ನುಂಟುಮಾಡುವಂತಹದೂ, ಪುಣ್ಯವನ್ನುಂಟುಮಾಡುವುದರಲ್ಲಿ ಮೇಲ್ಮೆ  ಪಡೆದು ವೇದಗಳಿಗಿಂತಲೂ ಶ್ರೇಷ್ಠವೆನಿಸುವಂತದೂ ಆಗಿರುವುದು. ಇಂತಹ ರಾಮಚರಿತ್ರವನ್ನು ಯಾವನು ಪಠಿಸುವನೋ ಅವನು ಎಲ್ಲಾ ಪಾಪಗಳಿಂದಲೂ ವಿಮುಕ್ತನಾಗುತ್ತಾನೆ. 

\medskip
\noindent\textbf{೩೭. ಇದಂ ಹಿ ತ್ರಿಷು ಲೋಕೇಷು} \hfill ಪುಟ \pageref{118}

\hfill ಮಹಾಭಾರತ ಆದಿಪರ್ವ ೨೬

ಈ ತಿಳಿಸಿದಂತಹ ಜ್ಞಾನವೆಂಬುದು ತ್ರಿಲೋಕವ್ಯಾಪಿಯಾಗಿ ನೆಲೆನಿಂತಿರುವುದು. 

\medskip
\noindent\textbf{೩೮. ಇಮಂ ವಿವಸ್ವತೇ ಯೋಗಂ} \hfill ಪುಟ \pageref{90}

\hfill ಭಗವದ್ಗೀತಾ ೪-೧, ೨, ೩

ಶತ್ರುತಾಪಕನಾದ ಅರ್ಜುನ! ನಿನಗೆ ತಿಳಿಸಿದ ನಾಶರಹಿತವಾದ ಈ ಯೋಗವನ್ನು ನಾನು ವಿವಸ್ವಂತನಿಗೆ ಪ್ರವಚನದ ಮೂಲಕ ಮೊದಲು ತಿಳಿಸಿದ್ದೆನಪ್ಪಾ. ವಿವಸ್ವಂತನು ಮನುವಿಗೂ, ಮನುವು ಇಕ್ಷ್ವಾಕುವಿಗೂ ಹೇಳಿದರು. ಈ ರೀತಿಯಾಗಿ ಕೈಯಿಂದ ಕೈಗೆ ಪರಂಪರೆಯಾಗಿ ದೊರಕಿದ ಈ ಯೋಗವನ್ನು ರಾಜರ್ಷಿಗಳು ಮನಗಂಡಿದ್ದರು. ಆ ಯೋಗವು ಕಾಲಕಳೆದಂತೆ ಕಳೆದಂತೆ ಮರೆಯಾಗಿ ಹೋಯಿತು. ಅದೇ ಯೋಗವನ್ನೇ ಪುರಾತನವಾಗಿ ಬಂದುದನ್ನೇ ಈಗ ನಾನು ತಿಳಿಸಿರುವೆನು.

\medskip
\noindent\textbf{೩೯. ಇಲ್ವಲಃ ಸಂಸ್ಕೃತಂ ವದನ್} $\star$ \hfill ಪುಟ \pageref{232}

\medskip
\noindent\textbf{40. ಉತ್ತಿಷ್ಠೋತ್ತಿಷ್ಠ ಗೋವಿಂದ} \hfill ಪುಟ \pageref{147}

\hfill ಸುಪ್ರಭಾತ

\medskip
\noindent\textbf{೪೧. ಉದಾತ್ತೇ ನಿಷಾದಗಾಂಧಾರೌ ಅನುದಾತ್ತೌ} \hfill ಪುಟ \pageref{13}

\hfill ಪಾಣಿನೀಯ ಶಿಕ್ಷಾ ೧೨

ಉದಾತ್ತ ಸ್ವರದಲ್ಲಿ ನಿಷಾದ - ಗಾಂಧಾರಗಳೂ, ಅನುದಾತ್ತದಲ್ಲಿ ಋಷಭ-ಧೈವತಗಳೂ, ಸ್ವರಿತದಲ್ಲಿ ಷಡ್ಜ- ಮಧ್ಯಮ - ಪಂಚಮ ಸ್ವರಗಳೂ ಹೀಗೆ ಮೂರು ಸ್ವರಗಳಲ್ಲಿ ಏಳು ಸ್ವರಗಳೂ ಅಡಕವಾಗುತ್ತವೆ.

\medskip
\noindent\textbf{೪೨. ಉದ್ಧೃತ್ಯ ಭುಜಮುಚ್ಯತೇ} \hfill ಪುಟ \pageref{130}

`ಸತ್ಯ ಸತ್ಯಂ' ಇತ್ಯಾದಿ ಗ್ರಂಥದ ಭಾವಾನುವಾದವನ್ನು ಗಮನಿಸಿ.

\medskip
\noindent\textbf{೪೩. ಉಪವಿಷ್ಟೇ ತದಾ ತಸ್ಮಿನ್} \hfill ಪುಟ \pageref{203}

\hfill ರಾಮಾಯಣ ಬಾಲಕಾಂಡ ೨-೨೭

ಲೋಕಪಿತಾಮಹನಾದ ಬ್ರಹ್ಮನೇ ಸಾಕ್ಷಾತ್ತಾಗಿ ಬಂದು ಕುಳಿತುಕೊಳ್ಳಲಾಗಿ ತತ್ಪರಾಯಣವಾದ ಮನಸ್ಸಿನಿಂದ ವಾಲ್ಮೀಕಿಯು ಧ್ಯಾನಾರೂಢನಾದನು.

\medskip
\noindent\textbf{೪೪. ಊರ್ಧ್ವಮೂಲಮದಶ್ಯಾಖಂ} $\star$ \hfill ಪುಟ \pageref{13}

\hfill ಭಗವದ್ಗೀತಾ ೧೫-೧

\medskip
\noindent\textbf{೪೫. ಋಷೇರ್ಧರ್ಮಾತ್ಮನಸ್ತಸ್ಯ} \hfill ಪುಟ \pageref{165}

\hfill ರಾಮಾಯಣ ಬಾಲಕಾಂಡ ೨-೧೩

`ತದಾ ತು ತಂ ದ್ವಿಜಂ ದೃಷ್ಟ್ವಾ' ಇತ್ಯಾದಿ ಗ್ರಂಥದ ಭಾವಾನುವಾದವನ್ನು ಗಮನಿಸಿ.

\medskip
\noindent\textbf{೪೬. ಏಕರಾತ್ರೇಣ ಕಲಲಂ} \hfill ಪುಟ \pageref{85}

\hfill ಯೋಗಾರ್ಣವ

ಗರ್ಭವು ಹುಟ್ಟಿದ ಮೊದಲ ದಿನ ಕಲ್ಲರೂಪವೂ, ಐದುದಿನದಲ್ಲಿ ಬುದ್ಬುದರೂಪವೂ ಆಗುವುದು.

\medskip
\noindent\textbf{೪೭. ಏತದಿಚ್ಛಾಮ್ಯಹಂ ಶ್ರೋತುಂ} \hfill ಪುಟ \pageref{153},\pageref{194}

\hfill ರಾಮಾಯಣ ಬಾಲಕಾಂಡ ೧-೫

ಹಿಂದೆ ಹೇಳಿದ ವಿಷಯವನ್ನು ನಿನ್ನಿಂದ ಕೇಳಿ ತಿಳಿಯಲು ನನಗೆ ಹಿರಿದಾದ ಆಸಕ್ತಿಯಿರುವುದು. ಓ ಮಹರ್ಷಿಯೇ! ಇಷ್ಟಾರು ಗುಣಗಳಿಂದ ಪರಿಪೂರ್ಣನಾದ ನರನನ್ನರಿಯಲು ಸಮರ್ಥನಾಗಿರುವೆಯಷ್ಟೆ.

\medskip
\noindent\textbf{೪೮. ಏತಾವಾನಸ್ಯ ಮಹಿಮಾ} \hfill ಪುಟ \pageref{249}

\hfill ಪುರುಷಸೂಕ್ತ

`ಸಹಸ್ರಶೀರ್ಷಾ' ಇತ್ಯಾದಿ ಗ್ರಂಥದ ಭಾವಾನುವಾದವನ್ನು ಗಮನಿಸಿ.

\medskip
\noindent\textbf{೪೯. ಓಮಿತ್ಯುಕ್ತ್ವಾ ಗಣೇಶೋಽಪಿ} \hfill ಪುಟ \pageref{119}

\hfill ಮಹಾಭಾರತ ಆದಿಪರ್ವ ೧-೧೧೫

ಹಾಗೆಯೇ ಆಗಲೆಂದೊಪ್ಪಿಕೊಂಡು ಗಣೇಶನೂ ಸಹ ಲೇಖಕನಾದನು.

\medskip
\noindent\textbf{೫೦. ಕಂಠೇಽಸ್ತಿ ಭಾರತೀಸ್ಥಾನಂ} $\star$ \hfill ಪುಟ \pageref{31}

\hfill ತಂತ್ರ ಶಾಸ್ತ್ರ

\medskip
\noindent\textbf{೫೧. ಕರಕಮಲನಿದರ್ಶಿತಾತ್ಮಮುದ್ರಃ} $\star$ \hfill ಪುಟ \pageref{48},\pageref{92},\pageref{217},\pageref{226}

\hfill ಭಗವದ್ಗೀತಾಧ್ಯಾನಶ್ಲೋಕ

\medskip
\noindent\textbf{೫೨. ಕಲ್ಯಾಣೋಲ್ಲಾಸಸೀಮಾ} \hfill ಪುಟ \pageref{206}

\hfill ವಿಶ್ವಗುಣಾದರ್ಶಚಂಪೂ

ವಿಶ್ವದ ಸಮಸ್ತ ಜೀವಿಗಳ ಮಂಗಳಕ್ಕೆ  ಅಡ್ಡಿ ಆತಂಕ ಬಾರದೇ ಸದಾ ಅಭಿವೃದ್ಧಿಶೀಲವಾಗಿ ಬೆಳೆಯುವಂತೆ ಮಾಡಿ, ಅದು ತನ್ನ ನೆಲೆದಾಣದವರೆವಿಗೂ ತಲುಪುವ ವಿಧಾನದಲ್ಲಿ ವ್ಯವಸ್ಥೆಗೊಳಿಸುವ ಮರ್ಯಾದಾಪುರುಷರೂಪಿಯೂ, ಜಲಭರಿತವಾದ ಕಾರ್ಮೋಡದೋಪಾದಿ ನೀಲವರ್ಣ ಹೊಂದಿ ಜೀವಿಗಳ ಮುಬ್ಬಗೆಯ ತಾಪಗಳನ್ನೂ ಕಳೆದು ತಂಪಾದ ಸ್ಥಿತಿಗೆ ತರುವಂತಹ ಅಮೃತ ವರ್ಷವನ್ನು ಸುರಿಸುವಂಥದಾಗಿ ನಯನ ಮನಗಳನ್ನು ತನ್ನೆಡೆಗೆ ಸೆಳೆದುಕೊಳ್ಳುವಂಥದೂ, ಸಾಕೇತವೆಂದು ಹೆಸರುವೆತ್ತ ಅಯೋಧ್ಯಾನಗರಿಯನ್ನೇ ತನ್ನ ವಿಹಾರ ಧಾಮವಾಗಿ ಹೊಂದಿ, ಭೂತಕೋಟಿಗಳ ಜನ್ಮ ವ್ಯಾಧಿ ಜರಾವಿಪತ್ತಿಮರಣಮಯವಾದ ಸಂಸಾರ ಚಕ್ರದ ಗಹನವಾದ ಗತಿಯಿಂದಾದ ಬಳಲಿಕೆಯನ್ನು ನಿವಾರಿಸುವಂತಹ ಪ್ರಣಯಕ್ಕೆ ಏಕೈಕ ಲಕ್ಷ್ಯವೂ, ತನ್ನ ಪುಂಸಾಂ ಮೋಹನವಾದ ದಿವ್ಯರೂಪ ಸೌಂದರ್ಯದಿಂದ ಕಾಮನನ್ನೂ ಲಜ್ಜೆಗೊಳಿಸುವಂಥದೂ, ತನ್ನ ವಾಮಭಾಗದಲ್ಲಿ ಬೆಳಗುವ ಜನಕಪುತ್ರಿಯ ಕರುಣಾಪೂರ್ಣವಾದ ಅಪಾಂಗವೀಕ್ಷಣವೈಭವವು ಲೋಕಕ್ಕೆ ದೊರಕುವಂತೆ ಮಾಡುವುದೂ, ದಿಗಂತವಿಶ್ರಾಂತದಿಂದ ಸರ್ವತೋಮುಖ ಹಿರಿಮೆಯುಳ್ಳದ್ದೂ, ದೇವತೆಗಳ ಸಾರ್ವಕಾಲಿಕವಾದ ಮೆಚ್ಚುಗೆಗೆ ಪಾತ್ರವೂ ಆದ `ರಾಮ' ಎಂಬ ದಿವ್ಯ ನಾಮಹೊಂದಿರುವ ಪ್ರಕಾಶಮಯವಾದ ಶಕ್ತಿಯು ನಮ್ಮೆಲ್ಲರಿಗೂ ನಿತ್ಯ ಕ್ಷೇಮವುಂಟುಮಾಡಲಿ.

\medskip
\noindent\textbf{೫೩. ಕಾಮಾರ್ಥಗುಣಸಂಯುಕ್ತಂ} \hfill ಪುಟ \pageref{212}

\hfill ರಾಮಾಯಣ ಬಾಲಕಾಂಡ ೩-೮
 
 (ವಾಲ್ಮೀಕಿಗಳು ಪ್ರಾಚೀನಾಗ್ರವಾದ ದರ್ಭೆಯಲ್ಲಿ ಆಸೀನರಾಗಿ ಧರ್ಮವೀರ್ಯದಿಂದ ರಾಮಾಯಣವನ್ನು ಅಂಗೈನೆಲ್ಲಿಯಂತೆ ಕಂಡುಕೊಂಡು ರಘುವಂಶದ ಚರಿತೆಯನ್ನು ರಚಿಸಿದರು.) ಅದು ಪುರುಷಾರ್ಥಗಳಲ್ಲಿ ಒಂದಾದ ಯಾವ ಕಾಮರೂಪವಾದ ಅರ್ಥವಿದೆಯೋ ಅದರ ಗುಣಗಳಿಂದ ಸಂಯುಕ್ತವಾದದ್ದು. ಧರ್ಮರೂಪವಾದ ಯಾವ ಅರ್ಥವಿದೆಯೋ ಆ ಗುಣಗಳಿಂದ ಅನ್ವಿತವಾದದ್ದು. ರತ್ನಗಳಿಂದ ಅಢ್ಯವಾದ ಸಮುದ್ರದಂತೆಯೂ ಸರ್ವರಿಗೂ ಶ್ರುತಿಮನೋಹರವಾಗಿಯೂ ಇರುವ ಕಾವ್ಯ.
 
\medskip
\noindent\textbf{೫೪. ಕಾಯೇನ ವಾಚಾ ಮನಸೇಂದ್ರಿಯೈರ್ವಾ} \hfill ಪುಟ \pageref{219}

ಅವನನ್ನು ಹೊತ್ತುತಂದ ಪ್ರಕೃರಿಯ ಸ್ವಭಾವಕ್ಕೆ ತಕ್ಕಂತೆ, ದೇಹ, ವಾಕ್, ಮನಸ್ಸು, ಇಂದ್ರಿಯಗಳು, ಬುದ್ಧಿ, ಆತ್ಮ ಇವುಗಳಿಂದ ಏನೇನು ಮಾಡುವೆನೋ ಅವೆಲ್ಲವನ್ನೂ ಪೂರ್ತಿಯಾಗಿ ಪರಮಪುರುಷನಾದ ನಾರಾಯಣನಿಗಾಗಿಯೇ ಎಂದು ಸಮರ್ಪಣೆ ಮಾಡುವೆನು.

\medskip
\noindent\textbf{೫೫. ಕಾಲೇ ವರ್ಷತು ಪರ್ಜನ್ಯ} \hfill ಪುಟ \pageref{149}

ಸಕಾಲದಲ್ಲಿ ಮಲೆಯು ವರ್ಷಿಸುತ್ತಿರಲಿ! ಭೂಮಿಯು ಸದಾ ಸಸ್ಯ ಸಮೃದ್ಧಿಯಿಂದ ಕೂಡಿರಲಿ! ಈ ಭರತಭೂಮಿಯು ಧರ್ಮಕ್ಷೋಭೆಯಿಲ್ಲದೆ ನೆಮ್ಮದಿಯಿಂದಿರಲಿ! ದೇಶಕ್ಕೆಲ್ಲಾ ಸತ್ಯವನ್ನು ಸಾರಿ ಕ್ಷೇಮಂಕರರಾದ ಬ್ರಹ್ಮಪರರು ಬ್ರಹ್ಮಚಿಂತನೆಗೆ ಆತಂಕ ತರುವಂತಹ ಯಾವುದೇ ಭಯಕ್ಕೆ ಎಡೆಯಿಲ್ಲದೇ ಸುಖಿಗಳಾಗಿರುವಂತಾಗಲಿ!

\medskip
\noindent\textbf{೫೬. ಕಿಮಿದಂ ವ್ಯಾಹೃತಂ ಮಯಾ} $\star$ \hfill ಪುಟ \pageref{20},\pageref{157}

\hfill ರಾಮಾಯಣ ಬಾಲಕಾಂಡ ೨-೧೬

\medskip
\noindent\textbf{೫೭. ಕುಶೀಲವೌ ಚೈವ ಮಹಾತಪಸ್ವಿನೌ} \hfill ಪುಟ \pageref{167}

\hfill ರಾಮಾಯಣ ಬಾಲಕಾಂಡ ೪-೩೫

ಕುಶಲವರೂ ಸಹ ಮಹಾತಪಸ್ವಿಗಳು.

\medskip
\noindent\textbf{೫೮. ಕುಶೀಲವೌ ತು ಧರ್ಮಜ್ಞೌ} \hfill ಪುಟ \pageref{238},\pageref{240}

\hfill ರಾಮಾಯಣ ಬಾಲಕಾಂಡ ೪-೫

`ಕೃತ್ವಾಪಿ' ಇತ್ಯಾದಿ ಗ್ರಂಥದ ಭಾವಾನುವಾದವನ್ನು ಗಮನಿಸಿ.

\medskip
\noindent\textbf{೫೯. ಕೂಜಂತಂ ರಾಮ ರಾಮೇತಿ} $\star$ \hfill ಪುಟ \pageref{136},\pageref{148},\pageref{169},\pageref{184}

\hfill ಬುಧಕೌಶಿಕಮುನಿವಿರಚಿತ ರಾಮರಕ್ಷಾ ಸ್ತೋತ್ರ

\medskip
\noindent\textbf{೬೦. ಕೂರ್ಮೋಂಗಾನೀವ ಸರ್ವಶಃ} \hfill ಪುಟ \pageref{222}

\hfill ಭಗವದ್ಗೀತಾ ೨-೫೮

ಆಮೆಯು ತನ್ನ ಅಂಗಗಳನ್ನೆಲ್ಲವನ್ನೂ ಒಳಗೆಳೆದುಕೊಳ್ಳುವಂತೆ (ಯೋಗಿಯಾದವನು ತನ್ನ ಮನಸ್ಸನ್ನು ಒಳಮುಖವಾಗಿ ಮಾಡಿಕೊಳ್ಳುತ್ತಾನೆ.)

\medskip
\noindent\textbf{೬೧. ಕೃತ್ವಾಪಿ ತನ್ಮಹಾಪ್ರಾಜ್ಞಃ} \hfill ಪುಟ \pageref{165},\pageref{167},\pageref{209},\pageref{238}

\hfill ರಾಮಾಯಣ ಬಾಲಕಾಂಡ ೪-೩ ರಿಂದ ೬

ಮಹಾಪ್ರಾಜ್ಞನಾದ ವಾಲ್ಮೀಕಿಮುನಿಯು ಮುಂದಾಗತಕ್ಕ ಕಥೆಯನ್ನೊಳಗೊಂಡ ಉತ್ತರಕಾಂಡ ಸಹಿತವಾದ ರಾಮಚರಿತ್ರೆಯನ್ನು ರಚಿಸಿದ ಕೂಡಲೇ, ಅದನ್ನು ಪ್ರಯೋಗಕ್ಕಿಡುವ ಪ್ರಭುತ್ವವನ್ನು ಪಡೆದವನಾಗಿ, ಇದನ್ನು ಯಾರು ಪ್ರಯೋಗಿಸಬಲ್ಲನೆಂದು ಚಿಂತಿಸಿದನು.

ಮಹರ್ಷಿಯು ಅದೇ ಚಿಂತೆಯಿಂದ ತದ್ಭಾವಭಾವಿತಾಂತಃಕರಣನಾಗಿರಲಾಗಿ ಆ ಕೂಡಲೇ ಮುನಿವೇಷಧಾರಿಗಳಾದ ಕುಶಲವರೀರ್ವರು ಬಂದು ಇವನ ಪಾದೋಪಸಂಗ್ರಹಣ ಮಾಡಿದರು.

ಅವರೀರ್ವರು ಧರ್ಮವನ್ನರಿತವರೂ, ರಾಜಪುತ್ರರೂ, ಕೀರ್ತಿಮಂತರೂ, ಸಹೋದರರೂ, ಗಾನಯೋಗ್ಯವಾದ ಶಾರೀರಸಂಪತ್ತುಳ್ಳವರೂ, ತನ್ನ ಆಶ್ರಮದಲ್ಲೇ ವಾಸಮಾಡುತ್ತಿರುವವರೂ ಆಗಿರುವುದನ್ನು ಕಂಡುಕೊಂಡನು.

ಇವರ ಬಗ್ಗೆ ಇದಕ್ಕೆ ಮೊದಲೇ ಗುರುವಾಗಿ ಪ್ರಭುತ್ವ ವಹಿಸಿದ್ದ ಆ ವಾಲ್ಮೀಕಿಯು, ತತ್ತ್ವಾರ್ಥ ಧಾರಣಯೋಗ್ಯವಾದ ಉತ್ತಮವಾದ ಮೇಧೆಯಿಂದ ಕೂಡಿದವರೂ ಸಮಗ್ರವಾದ ವೇದಗಳಲ್ಲಿ ಪರಿಣತರೂ ಆಗಿರುವುದನ್ನು ಕಂಡು ವೇದಾರ್ಥವನ್ನು ವಿಸ್ತರಿಸಲೋಸುಗ ಅವರೀರ್ವರನ್ನು ಸ್ವೀಕರಿಸಿದನು.

\medskip
\noindent\textbf{೬೨. ಕೃಷಿಗೋರಕ್ಷವಾಣಿಜ್ಯಂ} \hfill ಪುಟ \pageref{75},\pageref{225}

\hfill ಮಹಾಭಾರತ ಶಾಂತಿಪರ್ವ ೮೯-೭

ಕ್ಷೇತ್ರವ್ಯವಸಾಯ, ಗೋರಕ್ಷಣೆ, ವ್ಯಾಪಾರ ಇವಿಷ್ಟೂ ಜನರ ಇಹದ ಜೀವನಕ್ಕೆ ಬೇಕಾದವು. ತ್ರಯೀವಿದ್ಯೆಯು ಊರ್ಧ್ವಗತಿಯನ್ನು ಸೇರಿಸಿಕೊಂಡು ಜೀವನಸಾಧನವಾಗಿರುವುದು. ಅದು ಸಮಸ್ತಭೂತಗಳನ್ನು ಧಾರಣೆ ಮಾಡುವವಿದ್ಯೆಯಾಗಿದೆ. ಆ ತ್ರಯೀವಿದ್ಯೆಯು ಲೋಕಹಿತಕಾರಕವಾಗಿ ಕೆಲಸಮಾಡುವುದನ್ನು ತಡೆಯುವರಾರುಂಟೋ ಅವರು ದಸ್ಯುಗಳೆನಿಸುವರು. ಅವರನ್ನು ಸದೆಬಡಿಯಲು ಬ್ರಹ್ಮನು ಕ್ಷತ್ರ ಸೃಷ್ಟಿಯನ್ನು ಮಾಡಿದನು.

\medskip
\noindent\textbf{೬೩. ಕೌಸಲ್ಯಾಽ ಜನಯದ್ರಾಮಂ} \hfill ಪುಟ \pageref{252}

\hfill ರಾಮಾಯಣ ಬಾಲಕಾಂಡ ೧೮-೧೦

ಬೃಹಸ್ಪತಿಯು ಚಂದ್ರನೊಡನೆ ಸಂಗತವಾಗಿ ಉದಯಹೊಂದುತ್ತಿರಲಾಗಿ, ಜಗತ್ತಿಗೇ ನಾಯಕನೂ ಸರ್ವಲೋಕಕ್ಕೂ ನಮಸ್ಕಾರ್ಯನೂ ಆದಂತಹ ಶ್ರೀರಾಮನನ್ನು ಸಮಸ್ತಲಕ್ಷಣಸಂಪತ್ತಿನೊಡನೆ ಕೌಸಲ್ಯಾದೇವಿಯು ಪ್ರಸವಿಸಿದಳು.

\medskip
\noindent\textbf{೬೪. ಕ್ಷೀರ ಸಾಗರ ತರಂಗ} \hfill ಪುಟ \pageref{147}

\hfill ಮುಕುಂದ ಮಾಲಾ

\medskip
\noindent\textbf{೬೫. ಗಜೋಽಪಿ ಮಿಥ್ಯಾ} $\star$ \hfill ಪುಟ \pageref{54}

\medskip
\noindent\textbf{೬೬. ಗೀತೀ ಶೀಘ್ರೀ ಶಿರಃ ಕಂಪೀ} $\star$ \hfill ಪುಟ \pageref{5}

\hfill ಪಾಣಿನೀಯ ಶಿಕ್ಷಾ
 
\medskip
\noindent\textbf{೬೭. ಗುಶಬ್ದಸ್ತ್ವಂಧಕಾರಃ ಸ್ಯಾತ್} \hfill ಪುಟ \pageref{98}

\hfill ಅದ್ವಯ ತಾರಕೋಪನಿಷತ್ ೧೯

`ಗು' ಶಬ್ದವು ಅಂಧಕಾರವೆಂಬ ಅರ್ಥವನ್ನೂ `ರು' ಶಬ್ದವು ಅದರ ನಿರೋಧವನ್ನೂ ಹೇಳುತ್ತದೆ. ಅಜ್ಞಾನವೆಂಬ ಅಂಧಕಾರವು ಜೀವಿಯನ್ನಾವರಿಸದಂತೆ ತಡೆಯುವುದರಿಂದ ಜ್ಞಾನದಾತನು ಗುರುವೆಂದು ಕರೆಯಲ್ಪಡುತ್ತಾನೆ.

\medskip
\noindent\textbf{೬೮. ಗ್ರಂಥಗ್ರಂಥಿಂ ತದಾ ಚಕ್ರೇ} \hfill ಪುಟ \pageref{119}

\hfill ಮಹಾಭಾರತ ಆದಿಪರ್ವ ೧-೧೧೬

ಗಣೇಶನ ಒಪ್ಪಿಗೆಯಂತೆ ಕೈನಿಲ್ಲಿಸದೇ ಬರೆಯಲು ಅನುವಾಗಿ ಹೇಳಬೇಕೆಂಬುದನ್ನು ನಿರ್ವಹಿಸಲು ವ್ಯಾಸ ಮುನಿಯು ಗೂಢಾರ್ಥವುಳ್ಳಂತಹ, ಗಂಟಿನಂತೆ ಸುಲಭವಾಗಿ ಅರಿವಿಗೆ ಬಾರದ ಗ್ರಂಥಭಾಗವನ್ನು ಅಲ್ಲಲ್ಲೇ ಇಟ್ಟನು.

\medskip
\noindent\textbf{೬೯. ಗ್ರಂಥತಃ ಅರ್ಥತಶ್ಚ ಶಾಸ್ತ್ರಪ್ರಾಪ್ತಿಃ} \hfill ಪುಟ \pageref{105}

\hfill `ಶಿಷ್ಯೋಪಾಧ್ಯಾಯಿಕಯಾ' ಗ್ರಂಥದ ಭಾವಾನುವಾದವನ್ನು ಗಮನಿಸಿ.

\medskip
\noindent\textbf{೭೦. ಚಕ್ರವರ್ತಿತನೂಜಾಯ}  \hfill ಪುಟ \pageref{193}

ದಶರಥ ಚಕ್ರವರ್ತಿಯ ಔರಸಪುತ್ರನೂ, ಸಂಪೂರ್ಣವಾದ ಭೂಮಿಯಲ್ಲೂ ಹೆಸರು ಪಡೆದು ಸಾರ್ವಭೌಮನೂ, ಕೋಸಲದೇಶಕ್ಕೆ ರಾಜನೂ, ಸರ್ವೋತ್ಕೃಷ್ಟವಾದ ಗುಣಗಳಿಂದ ಕೂಡಿದವನೂ ಆದ ಶ್ರೀರಾಮಚಂದ್ರ ಪ್ರಭುವಿಗೆ ಮಂಗಳವಾಗಲಿ.

\medskip
\noindent\textbf{೭೧. ಚತ್ವಾರಿ ವಾಕ್ ಪರಿಮಿತಾ ಪದಾನಿ} $\star$ \hfill ಪುಟ \pageref{22},\pageref{153},\pageref{252}

\hfill ತೈತ್ತಿರೀಯ ಬ್ರಾಹ್ಮಣ ೨-೮-೮-೧೪

\medskip
\noindent\textbf{೭೨. ಚಾರುನಿಃಸ್ವನಮ್} \hfill ಪುಟ \pageref{250}

\hfill ರಾಮಾಯಣ ಬಾಲಕಾಂಡ ೨-೯

`ತಸ್ಯಾಭ್ಯಾಶೇ ತು' ಇತ್ಯಾದಿ ಗ್ರಂಥದ ಭಾವಾನುವಾದವನ್ನು ಗಮನಿಸಿ.

\medskip
\noindent\textbf{೭೩. ಚಿಂತಯಾಮಾಸ ಕೋನ್ವೇತತ್} $\star$ \hfill ಪುಟ \pageref{158}

\hfill ರಾಮಾಯಣ ಬಾಲಕಾಂಡ ೪-೩

\medskip
\noindent\textbf{೭೪. ಚಿರನಿರ್ವೃತ್ತಮಪ್ಯೇತತ್} \hfill ಪುಟ \pageref{204},\pageref{243}

\hfill ರಾಮಾಯಣ ಬಾಲಕಾಂಡ ೪-೧೮

ರಾಮಾಯನ ಗಾನದ ಇಂಪು ಅತ್ಯಾಶ್ಚರ್ಯಕರವಾಗಿದೆ. ಶ್ಲೋಕ ರಚನೆಯಂತೂ ಇನ್ನೂ ವಿಶೇಷವಾಗಿರುವುದು. ಈ ಕಥೆಯು ಬಹಳ ಹಿಂದೆಯೇ ನಡೆದುಹೋದ ಕಥೆಯಾಗಿದ್ದರೂ ಕಣ್ಣೆದುರಿಗೆ ನಡೆಯುತ್ತಿರುವಂತೆ ಗಾನ ಸಾಹಿತ್ಯಗಳು ತೋರಿಸುತ್ತಿವೆ.

\medskip
\noindent\textbf{೭೫. ಚೈತ್ರೇ ನಾವಮಿಕೇ ತಿಥೌ} \hfill ಪುಟ \pageref{147}

\hfill ರಾಮಾಯಣ ಬಾಲಕಾಂಡ ಶ್ಲೋಕ ೧೮-೮

ಚೈತ್ರಮಾಸದ ನವಮೀ ತಿಥಿಯಂದು ಶ್ರೀರಾಮನು ಆವಿರ್ಭವಿಸಿದನು.

\medskip
\noindent\textbf{೭೬. ಜನ್ಮವ್ಯಾಧಿಜರಾವಿಪತ್ತಿಮರಣೈಃ} \hfill ಪುಟ \pageref{173}

`ಯಃ ಕರ್ಣಾಂಜಲಿಸಂಪುಟೈಃ' ಇತ್ಯಾದಿ ಗ್ರಂಥದ ಭಾವಾನುವಾದವನ್ನು ಗಮನಿಸಿ.

\medskip
\noindent\textbf{೭೭. ಜಾತಿಭಿಃ ಸಪ್ತಭಿರ್ಬದ್ಧಂ} \hfill ಪುಟ \pageref{212}

\hfill ರಾಮಾಯಣ ಬಾಲಕಾಂಡ ೪-೮

`ಪಾಠ್ಯೇ ಗೇಯೇ ಚ' ಇತ್ಯಾದಿ ಗ್ರಂಥದ ಭಾವಾನುವಾದವನ್ನು ಗಮನಿಸಿ.

\medskip
\noindent\textbf{೭೮. ಜಿಹ್ವೇ ಕೀರ್ತಯ ಕೇಶವಂ} $\star$ \hfill ಪುಟ \pageref{214},\pageref{232},\pageref{271}

\hfill ಮುಕುಂದಮಾಲಾಸ್ತು

ನಾಲಗೆಯೇ (ವಾಗಿಂದ್ರಿಯವೇ)! ಕೇಶವನ ಸಂಕೀರ್ತನೆ ಮಾಡು. ಮನಸ್ಸೇ! ಮುರಾರಿಯಾದ ವಿಷ್ಣುವನ್ನು ಧ್ಯಾನಿಸು! ಹಸ್ತಯುಗಲವೇ! ಶ್ರೀಧರನನ್ನು ಚೆನ್ನಾಗಿ ಅರ್ಚಿಸು. ಕರ್ಣಯುಗಲವೇ! ನೀನು ಅಚ್ಯುತನ (ಚ್ಯುತಿ ಇಲ್ಲದಂತೆ ಮಾಡುವ) ಕಥೆಗಳನ್ನು ಆಲಿಸು. ನೇತ್ರಯುಗಲವೇ! ಕೃಷ್ಣನನ್ನು (ಮೂರ್ತಿರೂಪವನ್ನು) ನೋಡು. ಪಾದದ್ವಯವೇ! ಶ್ರೀಹರಿಯ ಆಲಯಕ್ಕೆ ಸಂಚರಿಸು. ಘ್ರಾಣೇಂದ್ರಿಯವೇ! ಮುಕ್ತಿದಾಯಕನಾದ ಜನಾರ್ದನನ ಪಾದಗಳಲ್ಲಿ ನೆಲೆಸಿದ್ದ ತುಳಸೀದಳವನ್ನು ಆಘ್ರಾಣಿಸು. (ಕರ್ಮಜ್ಞಾನೇಂದ್ರಿಯಗಳಿಗೆಲ್ಲಾ ಮೂಲಸ್ಥಾನವಾದ) ನನ್ನ ಶಿರಸ್ಸೇ! ಅಧೋಕ್ಷಜನಾದ ಸ್ವಾಮಿಯನ್ನು ಬಾಗಿ ನಮಿಸು.

\medskip
\noindent\textbf{೭೯. ಜ್ಞಾನ - ವಿಜ್ಞಾನತೃಪ್ತಾತ್ಮಾ} $\star$ \hfill ಪುಟ \pageref{66}

\hfill ಭಗವದ್ಗೀತಾ ೬-೮

\medskip
\noindent\textbf{೮೦. ಜ್ಞಾನಿನಾಮೂರ್ಧ್ವಗೋ ಭೂಯಾತ್} $\star$ \hfill ಪುಟ \pageref{13}

\hfill ಯೋಗಚೂಡಾಮಣ್ಯುಪನಿಷತ್ ೭೮

\medskip
\noindent\textbf{೮೧. ತಂತ್ರೀಲಯಸಮನ್ವಿತಂ} \hfill ಪುಟ \pageref{250}

\hfill ರಾಮಾಯಣ ಬಾಲಕಾಂಡ ೪-೮

`ಪಾಠ್ಯೇ ಗೇಯೇ ಚ' ಇತ್ಯಾದಿ ಗ್ರಂಥದ ಭಾವಾನುವಾದವನ್ನು ಗಮನಿಸಿ.

\medskip
\noindent\textbf{೮೨. ತಂ ಶೋಣಿತಪರೀತಾಂಗಂ} \hfill ಪುಟ \pageref{198}

\hfill ರಾಮಾಯಣ ಬಾಲಕಾಂಡ ೨-೧೧

ಆ ಗಂಡು ಕ್ರೌಂಚವು ಮೈಯೆಲ್ಲಾ ರಕ್ತದಿಂದ ತೋಯ್ದು ನೆಲಕ್ಕುರಳಿ ಹೊರಳಾಡುತ್ತಿದ್ದ ದೃಶ್ಯವನ್ನು ಕಂಡು ಅದರ ಸಂಗಾತಿಯಾದ ಹೆಣ್ಣು ಕ್ರೌಂಚವು ದಯನೀಯವಾದ ಸ್ವರದಿಂದ ರೋದಿಸಿತು.

\medskip
\noindent\textbf{೮೩. ತಚ್ಘುತ್ವಾ ಮುನಯಃ ಸರ್ವೇ} $\star$ \hfill ಪುಟ \pageref{158}

\hfill ರಾಮಾಯಣ ಬಾಲಕಾಂಡ ೪-೧೫

\medskip
\noindent\textbf{೮೪. ತತಃ ಕರುಣವೇದಿತ್ವಾತ್} \hfill ಪುಟ \pageref{199},\pageref{211}

\hfill ರಾಮಾಯಣ ಬಾಲಕಾಂಡ ೨-೧೪

ಕರುಣಭಾವವನ್ನು ಅರಿತ ಮಹರ್ಷಿವಾಲ್ಮೀಕಿಯು `ಇದು ಅಧರ್ಮ' ಎಂದು ನಿಶ್ಚಯಿಸಿ ಕ್ರೌಂಚದ ಆಕ್ರಂದವನ್ನು ಕೇಳಿ `ಮಾ ನಿಷಾದ' ಇತ್ಯಾದಿ ಮಾತನ್ನು ಆಡಿದನು.

\medskip
\noindent\textbf{೮೫. ತತಃ ಪಶ್ಯತಿ ಧರ್ಮಾತ್ಮಾ} \hfill ಪುಟ \pageref{204},\pageref{210},\pageref{248}

\hfill ರಾಮಾಯಣ ಬಾಲಕಾಂಡ ೩-೬

ನಂತರ ಧರ್ಮಪರಾಯಣವಾದ ಮನಸ್ಸುಳ್ಳ ವಾಲ್ಮೀಕಿಯು, ಹಿಂದೆ ಯಾವ ರಾಮವೃತ್ತಾಂತವು ನಡೆಯಿತೋ, ಅದೆಲ್ಲವನ್ನೂ ಯೋಗಾವಸ್ಥೆಯನ್ನು ತಳೆದವನಾಗಿ ಅಂಗೈ ನೆಲ್ಲಿಯಂತೆ ಕಂಡುಕೊಂಡನು.

\medskip
\noindent\textbf{೮೬. ತತ್ಸರ್ವಂ ತತ್ತ್ವತೋ ದೃಷ್ಟ್ವಾ} \hfill ಪುಟ \pageref{149}

\hfill ರಾಮಾಯಣ ಬಾಲಕಾಂಡ ೩-೭

(ಪ್ರಾಚೀನಾಗ್ರದರ್ಭೆಯಲ್ಲಿ ಆಸೀನರಾದ ವಾಲ್ಮೀಕಿಮಹರ್ಷಿಗಳು ಧರ್ಮಾತ್ಮರಾಗಿ ಯೋಗಸಮಾಧಿಯನ್ನೇರಿ) ಅಂಗೈಯಲ್ಲಿ ನೆಲ್ಲಿಯಕಾಯನ್ನು ನೋಡುವಂತೆ ಶ್ರೀರಾಮನ ಆ ಚರಿತ್ರೆಯನ್ನು ನೋಡಿ (ರಾಮಾಯಣಕಾವ್ಯವನ್ನು ರಚಿಸಿದರು.)

\medskip
\noindent\textbf{೮೭. ತತ್ಸರ್ವಂ ಧರ್ಮವೀರ್ಯೇಣ} \hfill ಪುಟ \pageref{210}

\hfill ರಾಮಾಯಣ ಬಾಲಕಾಂಡ ೩-೪

(ಪ್ರಾಚೀನಾಗ್ರದರ್ಭೆಯಲ್ಲಿ ಆಸೀನರಾದ ವಾಲ್ಮೀಕಿಗಳು) ಶ್ರೀರಾಮಚಂದ್ರನ ಆ ಚರಿತೆಯನ್ನು ಧರ್ಮವೀರ್ಯದಿಂದ - ಧರ್ಮವು ಅನುಗ್ರಹಿಸಿದ ಯೋಗದೃಷ್ಟಿಯಿಂದ ಅವಲೋಕಿಸಿದರು.

\medskip
\noindent\textbf{೮೮. ತತ್ಸವಿತುರ್ವರೇಣ್ಯಂ} \hfill ಪುಟ \pageref{186}

\hfill ಮಹಾನಾರಾಯಣೋಪನಿಷತ್ ೨೭-೧

ಯಾವನು ನಮಗೆ ಸದ್ಧರ್ಮಾದಿ ವಿಷಯಕವಾದ ಪ್ರಜ್ಞೆಯು ಹುಟ್ಟುವಂತೆ ಪ್ರೇರಿಸುವನೋ, ಅಂತಹ ಪ್ರಕಾಶಮಯನಾದ ಪ್ರಭುವಿನ - ಸವಿತೃದೇವನ, ಎಲ್ಲರಿಂದಲೂ ಪ್ರೀತ್ಯಾದರಗಳಿಂದ ಬಯಸತಕ್ಕದ್ದೂ, ಉಪಾಸಕರ ಪಾಪರಾಶಿಯನ್ನು ಹುರಿದು ಮತ್ತೆ ಮೊಳಕೆಬಾರದಂತೆ ದಗ್ಧ ಮಾಡತಕ್ಕದ್ದೂ ಆದ ತೇಜೋಮಯ ದಿವ್ಯರೂಪವನ್ನು ಧ್ಯಾನಿಸುವೆವು.

\medskip
\noindent\textbf{೮೯. ತತ್ರ್ಯೇವ ಅಂತರಧೀಯತೆ} $\star$\hfill ಪುಟ \pageref{181}

\hfill ರಾಮಾಯಣ ಬಾಲಕಾಂಡ ೨-೩೮

\medskip
\noindent\textbf{೯೦. ತದಾ ತು ತಂ ದ್ವಿಜಂ ದೃಷ್ಟ್ವಾ} \hfill ಪುಟ \pageref{199}

\hfill ರಾಮಾಯಣ ಬಾಲಕಾಂಡ ೨-೧೩

ಆಗ ಆ ಕಿರತನಿಂದ ಹೊಡೆದುರುಳಿಸಲ್ಪಟ್ಟ ಆ ಕ್ರೌಂಚ ಪಕ್ಷಿಯನ್ನು ಕಂಡು ಧರ್ಮಾತ್ಮನಾದ ವಾಲ್ಮೀಕಿ ಮುನಿಗೆ ಕಾರುಣ್ಯವು ತಲೆದೋರಿತು.

\medskip
\noindent\textbf{೯೧. ತದುಪಗತಸಮಾಸಸಂಧಿಯೋಗಂ} \hfill ಪುಟ \pageref{152},\pageref{212}

\hfill ರಾಮಾಯಣ ಬಾಲಕಾಂಡ ೨-೪೩

ವಾಲ್ಮೀಕಿಯು ನಿರ್ಮಿಸಿದ ಈ ಕಾವ್ಯವು ಅವಶ್ಯಕವಾದೆಡೆಗಳಲ್ಲಿ ಪದಗಳ ಅಚ್ಚು ಮೆಚ್ಚಾದ ಸೇರುವೆ (ಸಮಾಸ) ಯೇನು! ಸೊಬಗಾದ ವರ್ಣಗಳ ಹೊಂದಿಕೆಯೇನು! (ಸಂಧಿ) ಇವುಗಳುಳದ್ದೂ.  ಸಮತೆ ಮಾಧುರ್ಯ ಮೊದಲಾದ ಕಾವ್ಯಗುಣಗಳನ್ನು ಶಬ್ದ ಮತ್ತು ಅರ್ಥಗಳಲ್ಲಿ ಹೆಣೆದುಕೊಳ್ಳುತ್ತಿರುವ ವಾಕ್ಯ ರಚನೆಗಳುಳ್ಳದ್ದೂ, ಆಗಿ ರಘುವರನಾದ ಶ್ರೀರಾಮನ ಆದರ್ಶಜೀವನದ ಹೆಜ್ಜೆಯ ವೈಭವವನ್ನೂ, ಅದಕ್ಕೆ ವಿರುದ್ಧನಾದ ದಶಶಿರಸ್ಕ ರಾವಣನ ದಮನವನ್ನೂ ಅಂದವಾಗಿ ಚಿತ್ರಿಸುತ್ತಿರುವುದನ್ನು ಚೆನ್ನಾಗಿ ಆರಿಸಿ ಅನುಭವಿಸಿ ಪುರುಷಾರ್ಥಪೂರ್ಣರಾಗಿರಿ.

\medskip
\noindent\textbf{೯೩. ತಪಃ ಶ್ಲಾಘ್ಯೈಃ ಮಹಾತ್ಮಭಿಃ} \hfill ಪುಟ \pageref{211}

\hfill ರಾಮಾಯಣ ಬಾಲಕಾಂಡ ೪-೧೯

(ಸುಮಧುರಕಂಠಸಂಪತ್ತಿನಿಂದ ರಮ್ಯವಾಗಿ ರಾಮಚರಿತೆಯನ್ನು ಗಾನಮಾಡಿದ ಲವಕುಶರು) ತಪಸ್ಸಿನಿಂದ ಶ್ಲಾಘ್ಯರಾಗಿರುವ ಮಹರ್ಷಿಗಳಿಂದ ಪ್ರಶಂಸಿಸಲ್ಪಟ್ಟರು. 

\medskip
\noindent\textbf{೯೪. ತಪಸ್ಸ್ವಾಧ್ಯಾಯನಿರತಂ} $\star$ \hfill ಪುಟ \pageref{153},\pageref{162},\pageref{194},\pageref{232}

\hfill ರಾಮಾಯಣ ಬಾಲಕಾಂಡ ೧-೧

ತಪಸ್ವಿಯಾದ ವಾಲ್ಮೀಕಿಯು ತಪಸ್ಸು ಮತ್ತು ತತ್ತ್ವಚಿಂತನ ರೂಪವಾದ ಸ್ವಾಧ್ಯಾಯದಲ್ಲಿ ನಿರತನೂ ಮಾತಾಡುವುದರ ಮರ್ಮವರಿತವನೂ ಮುನಿಶ್ರೇಷ್ಠನೂ ಆದ ನಾರದನನ್ನು ವಿಧ್ಯುಕ್ತವಾದ ರೀತಿಯಲ್ಲಿ ಪ್ರಶ್ನಿಸಿದನು.

\medskip
\noindent\textbf{೯೫. ತಪಸ್ವಿನಾಂ ದಿವ್ಜಾತೀನಾಂ } \hfill ಪುಟ \pageref{211}

\hfill ರಾಮಾಯಣ ಬಾಲಕಾಂಡ ೪-೧೩

ತಪಸ್ವಿಗಳು, ಬ್ರಾಹ್ಮನರು ಹಾಗೂ ಸಾಧುಜನರ ಸನ್ನಿಧಿಯಲ್ಲಿ (ತತ್ತ್ವಜ್ಞರಾದ ಲವಕುಶರು ಸಮಾಹಿತಚಿತ್ತರಾಗಿ ಉಪದಿಷ್ಟವಾದ ಕ್ರಮವನ್ನನುಸರಿಸಿ ಗಾನಮಾಡಿದರು)

\medskip
\noindent\textbf{೯೬. ತನಸಸ್ತು ಪಾರೇ| ತಮಸಃ ಪರಸ್ತಾತ್|} \hfill ಪುಟ \pageref{251}

\hfill ತೈತ್ತಿರೀಯಾರಣ್ಯಕ ೩-೧೩, ೧೪

ಪ್ರಕೃತಿ ಮಂಡಲವನ್ನು ದಾಟಿ ಪರಮೋರ್ಧ್ವದಲ್ಲಿರುವಂತಹ ಆದಿತ್ಯವರ್ಣನಾದ ಪರಮಪುರುಷನನ್ನು ನಾನು ಅರಿತಿರುವೆನು.

\medskip
\noindent\textbf{೯೭. ತಮುವಾಚ ತತೋ ಬ್ರಹ್ಮಾ} \hfill ಪುಟ \pageref{203}

\hfill ರಾಮಾಯಣ ಬಾಲಕಾಂಡ ೨-೩೦

ನಂತರ ಹೀಗೆ ಶೋಲ ಪರಯಣನಾಗಿದ್ದ ವಾಲ್ಮೀಕಿ ಮುನಿಯನ್ನು ಕುರಿತು ನಗುಮುಖದಲ್ಲಿ ಬ್ರಹ್ಮನು ಹೇಳಿದನು- ``ಅಯ್ಯಾ! ಶ್ಲೋಕವನ್ನ್ನೇ ಛಂದೋಬದ್ಧವಾಗಿ ನೀನು ರಚಿಸಿರುವೆಯಪ್ಪಾ! ಆ ಬಗ್ಗೆ ನೀನು ವಿಚಿಕಿತ್ಸೆ ಮಾಡಬೇಕಾಗಿಲ್ಲ. 

\medskip
\noindent\textbf{೯೮. ತಸ್ಯ  ಚಿಂತಯಮಾನಸ್ಯ} \hfill ಪುಟ \pageref{238},\pageref{240}

`ಕೃತ್ವಾಪಿ' ಇತ್ಯಾದಿ ಭಾವಾನುವಾದವನ್ನು ಗಮನಿಸಿ. \hfill ರಾಮಾಯಣ ಬಾಲಕಾಂಡ ೪-೪

\medskip
\noindent\textbf{೯೯. ತಸ್ಯ ಧೀರಾಃ ಪರಿಜಾನಂತಿ} \hfill ಪುಟ \pageref{150},\pageref{251}

\hfill ತೈತ್ತಿರೀಯಾರಣ್ಯಕ ೩-೧೪

ಪರಮಪುರುಷನ ಆವಿರ್ಭಾವಕಾರಕವಾದ ಸನ್ನಿವೇಶದ ಒಳ ಮರ್ಮವನ್ನು, ಪಂಡಬುದ್ಧಿಯನ್ನು ಪಡೆದು ಪಂಡಿತರೆನಿಸುವ ತತ್ತ್ವದರ್ಶಿಗಳು ಮಾತ್ರವೇ ಅರಿಯಬಲ್ಲರು.

\medskip
\noindent\textbf{೧೦೦. ತಸ್ಮಾತ್ತು ಮಿಧುನಾದೇಕಂ} \hfill ಪುಟ \pageref{198}

\hfill ರಾಮಾಯಣ ಬಾಲಕಾಂಡ ೨-೧೦

ಕಾಡಿನ ಮೃಗಪಕ್ಷಿಗಳ ಮೇಲಿನ ವೈರಕ್ಕೆ ನೆಲೆದಾಣದಂತಿರುವ ಆ ನಿಷಾದನು ಆ ವಾಲ್ಮೀಕಿಯು ನೋಡುತ್ತಿರುವಂತೆಯೇ, ಆ ಕ್ರೌಂಚಗಳ ಜೋಡಿಯ ಪೈಕಿ ಒಂದಾದ ಪುರುಷ ಕ್ರೌಂಚವನ್ನು ಪಾಪಕಾರ್ಯಮಾಡುವ ಬಗ್ಗೆ ಮತ್ತೊಂದು ಆಲೋಚನೆಗೆ ಎಡೆಗೊಡದೇ ಕೊಂದೇಬಿಟ್ಟನು.

\medskip
\noindent\textbf{೧೦೧. ತಸ್ಯಾಂತೇ ಸುಷಿರಗ್ಂ ಸೂಕ್ಷ್ಮಂ} \hfill ಪುಟ \pageref{27}

\hfill ಮಹಾನಾರಾಯಣೋಪನಿಷತ್ -೧೧

ಆ ಹೃದಯಕೋಶದ ಮಧ್ಯಭಾಗದಲ್ಲಿ ಅತ್ಯಂತ ಸೂಕ್ಷ್ಮವಾಗಿ ಯೋಗಿ ಮಾತ್ರ ಸಂವೇದ್ಯವಾದ ಸುಷುಮ್ನಾರಂಧ್ರವಿರುವುದು. ಸೂಕ್ಷ್ಮಾತಮವಾದ ಆ ರಂಧ್ರದಲ್ಲಿ ಸರ್ವಮೂಲವೂ, ಸರ್ವಾತ್ಮವೂ ಆದ ಪರತತ್ತ್ವವು ಪ್ರತಿಷ್ಠಿತವಾಗಿರುವುದು.

\medskip
\noindent\textbf{೧೦೩. ತಸ್ಯಾಭ್ಯಾಶೇ ತು ಮಿಥುನಂ} $\star$ \hfill ಪುಟ \pageref{156a},\label{198a}

\hfill ರಾಮಾಯಣ ಬಾಲಕಾಂಡ ೨-೯

\medskip
\noindent\textbf{೧೦೪. ತಸ್ಯೈವಂ ಬ್ರುವತಶ್ಚಿಂತಾ} \hfill ಪುಟ \pageref{201}

\hfill ರಾಮಾಯಣ ಬಾಲಕಾಂಡ ೨-೧೬

ಆ ಕ್ರೌಂಚ ಪಕ್ಷಿಯ ವಧವನ್ನು ಕಂಡು `ಮಾ ನಿಷಾದ' ಇತ್ಯಾದಿಯಾಗಿ ಶಾಪೋಕ್ತಿಯನ್ನು ಉದ್ಗರಿಸಿದ ಆ ವಾಲ್ಮೀಕಿ ಮುನಿಗೆ ಮನದಲ್ಲಿ ಕಳವಳ ಹುಟ್ಟಿತು `ಈ ಪಕ್ಷಿಯ ಬಗೆಗೆ ಶೋಕಗೊಂಡು ವ್ಯಾಕುಲನಾದ ನನ್ನಿಂದ ಇದೇನು ವ್ಯಾಹೃತವಾಯಿತು?' (ಹೇಳಲ್ಪಟ್ಟಿತು) ಎಂದು.

\medskip
\noindent\textbf{೧೦೫. ತಸ್ಯೈಷ ಏವ ಶಾರೀರ ಆತ್ಮಾ} \hfill ಪುಟ \pageref{76a}

\hfill ತೈತ್ತಿರೀಯೋಪನಿಷತ್ ೨

ಅದಕ್ಕೂ ಈ ಶಾರಿರಪುರುಷನೇ ಆತ್ಮನಾಗಿರುವನು.

\medskip
\noindent\textbf{೧೦೬. ತೌ ಕದಾಚಿತ್ಸಮೇತಾನಾಂ} \hfill ಪುಟ \pageref{189}

\hfill ರಾಮಾಯಣ ಬಾಲಕಾಂಡ ೧೪-೮

(ಆಸೀನಾನಾಂ ಸಮೀಪಸ್ಥಾವಿದಂ ಕಾವ್ಯಮಗಾಯತಾಮ್) ಆ ಕುಶಲವರು ಯಾವಾಗಲೋ ಒಮ್ಮೆ ಭಾವಿತಾತ್ಮರಾಗಿ ಋಷಿಗಳೆಲ್ಲಾ ಒಕ್ಕೂಟವಾಗಿ ಒಂದೆಡೆ ಕುಳಿತಿರಲಾಗಿ ಅವರ ಪಕ್ಕದಲ್ಲೇ ಇದ್ದವರಾಗಿ ಈ ಕಾವ್ಯವನ್ನು ಹಾಡಿದರು.

\medskip
\noindent\textbf{೧೦೭. ತ್ರಿಪಾದೂರ್ಧ್ವ ಉದೈತ್ಪುರುಷಃ} \hfill ಪುಟ \pageref{252a}

\hfill ತೈತ್ತಿರೀಯಾರಣ್ಯಕ ೩-೧೩

ನಾರಾಯಣರೂಪಿಯಾದ ಮಹಾಪುರುಷನು ಊರ್ಧ್ವವಾದ ದಿವಿಸ್ಥಾನದಲ್ಲಿ ವಾಸುದೇವ - ಸಂಕರ್ಷಣ - ಪ್ರದ್ಯುಮ್ನರೂಪವಾದ ಉಳಿದ ಮೂರು (ಪಾದ) ಅಂಶಗಳುಳ್ಳವನಾಗಿ ಜ್ಞಾನಿಗಳ ಅಂತರಂಗಕ್ಕೆ ಗೋಚರನಾದನು. 

\medskip
\noindent\textbf{೧೦೮. ದಂಡೋ ದಮಯತಾಮ್ಸ್ಮಿ} \hfill ಪುಟ \pageref{93}

\hfill ಭಗವದ್ಗೀತಾ ೧೦-೩೮

ನಾನು (ಶ್ರೀಕೃಷ್ಣ) ದಮನಮಾಡುವವರಲ್ಲಿ ದಂಡರೂಪದಲ್ಲಿದ್ದೇನೆ. (ಅಧರ್ಮವನ್ನು) ಗೆಲ್ಲುವವರಲ್ಲಿ ನೀತಿರೂಪದಲ್ಲಿದ್ದೇನೆ. ಗೋಪ್ಯವಾದ ವಸ್ತುಗಳಲ್ಲಿ ಮೌನರೂಪದಲ್ಲಿದ್ದೇನೆ. ಜ್ಞಾನವಂತರಲ್ಲಿ ಜ್ಞಾನರೂಪದಲ್ಲಿದ್ದೇನೆ.

\medskip
\noindent\textbf{೧೦೯. ದಶಶಿರಸಶ್ಚ ವಧಂ ನಿಶಾಮಯಧ್ವಮ್} \hfill ಪುಟ \pageref{207}

\hfill ರಾಮಾಯಣ ಬಾಲಕಾಂಡ ೨-೪೩

`ತದುಪಗತ' ಇತ್ಯಾದಿ ಗ್ರಂಥದ ಭಾವಾನುವಾದವನ್ನು ಗಮನಿಸಿ.

\medskip
\noindent\textbf{೧೧೦. ದಿವ್ಯಂ ದದಾಮಿ ತೇ ಚಕ್ಷಃ} \hfill ಪುಟ \pageref{91}

\hfill ಭಗವದ್ಗೀತಾ ೧೧-೮

\medskip
\noindent\textbf{೧೧೧. ದೇವಾ ಮಾನುಷರೂಪೇಣ} \hfill ಪುಟ \pageref{146},\pageref{225},\pageref{247a}

ತತ್ತ್ವಜ್ಞರಾದ ಸಜ್ಜನರು ಯಾರೆಂದರೆ - ಇವರು ಮನುಷ್ಯದೇಹವನ್ನು ಧರಿಸಿ ಭೂಮಿಯಲ್ಲಿ ಸಂಚರಿಸುವಂತಹ ದೇವತೆಗಳೇ ಹೊರತು ಕೇವಲ ಮನುಷ್ಯರಲ್ಲ. 

\medskip
\noindent\textbf{೧೧೨. ದ್ಯುತಿಮಾನ್ ಧೃತಿಮಾನ್ ವಶೀ} \hfill ಪುಟ \pageref{249}

`ಇಕ್ಷ್ವಾಕುವಂಶಪ್ರಭವೋ' ಇತ್ಯಾದಿ ಗ್ರಂಥದ ಭಾವಾನುವಾದವನ್ನು ಗಮನಿಸಿ.

\medskip
\noindent\textbf{೧೧೩. ದ್ಯೌಃಪಿತಾ ಪೃಥಿವೀ ಮಾತಾ} $\star$ \hfill ಪುಟ \pageref{21}

\hfill ತೈತ್ತಿರೀಯ ಬ್ರಾಹ್ಮಣ ೩-೭-೫

\medskip
\noindent\textbf{೧೧೪. ಧರ್ಮಜ್ಞಸ್ಸತ್ಯಸಂಧಶ್ಚ} $\star$ \hfill ಪುಟ \pageref{155},\pageref{223}

\hfill ರಾಮಾಯಣ ಬಾಲಕಾಂಡ ೧-೧೨ ರಿಂದ ೧೯

\medskip
\noindent\textbf{೧೧೫. ಧರ್ಮನಾಶೋಽರ್ಥನಾಶಶ್ಚ} \hfill ಪುಟ \pageref{199a}

ಪುರುಷಾರ್ಥಗಳಲ್ಲಿ ಸೇರಿದಂತಹ ಧರ್ಮ ಮತ್ತು ಅರ್ಥ ಇವುಗಳ ಚ್ಯೂತಿಯು ಒದಗಿದರೆ ಕರುಣದ ಪ್ರಬೋಧವಾಗುವುದು. ಇದೇ ರೀತಿ ಬಾಂಧವರ, ಇಷ್ಟಜನರ ನಾಶದಿಂದಲೂ ಕರುಣದ ಪ್ರಬೋಧವಾಗುವುದು.

\medskip
\noindent\textbf{೧೧೬. ಧರ್ಮಸ್ಯ ಪ್ರಭುರಚ್ಯುತಃ} \hfill ಪುಟ \pageref{90a}

\hfill ಮಹಾಭಾರತ ಅನುಶಾಸನಪರ್ವ ೧೪೯-೧೩೭

ಆಚಾರದಿಂದ ಧರ್ಮವು ಆವಿರ್ಭವಿಸುವುದು. ಅಂತಹ ಧರ್ಮಕ್ಕೆ ಅಚ್ಯುತನೇ ಪ್ರಭುವು.

\medskip
\noindent\textbf{೧೧೭. ಧರ್ಮಾತ್ಮನೋ ಗುಣವತಃ} \hfill ಪುಟ \pageref{165a}

\hfill ರಾಮಾಯಣ ಬಾಲಕಾಂಡ ೨-೩೨

`ರಾಮಸ್ಯ ಚರಿತಂ' ಇತ್ಯಾದಿ ಗ್ರಂಥದ ಭಾವಾನುವಾದವನ್ನು ಗಮನಿಸಿ.

\medskip
\noindent\textbf{೧೧೮. ಧರ್ಮಾತ್ಮಾ ಸತ್ಯಸಂಧಶ್ಚ} \hfill ಪುಟ \pageref{165b},\pageref{207a},\pageref{250a}

\hfill ರಾಮಾಯಣ ಯುದ್ಧಕಾಂಡ ೯೦-೬೯

ಓ! (ನನ್ನ ಬಿಲ್ಲಿನಿಂದ ಹೊರಡುವ) ಬಾಣವೇ! ದಶರಥ ಪುತ್ರನಾದ ಈ ರಾಮನು ಧರ್ಮೈಕಪರಾಯಣವಾದ ಮನವುಳ್ಳವನಾಗಿದ್ದು (ಧರ್ಮಮಯವಾದ ಮೂರ್ತಿ ತಳೆದವನಾಗಿದ್ದು) ಜಗತ್ತನ್ನು ವಿಶ್ವಮೂಲವಾದ ಸತ್ಯದಲ್ಲಿ ಸಂಧಾನಗೊಳಿಸಿ ಸುಖೀವುಳ್ಳ ನಿರ್ಮಾಣದ ಗುರಿಯುಳ್ಳವನಾಗಿದ್ದು, ಪೌರುಷದಲ್ಲಿ ಸರಿಸಾಟಿಯಾದ ಮತ್ತೊಬ್ಬರಿಲ್ಲದ ಅಸದೃಶ ಅದ್ವಿತೀಯ ವೀರನಾಗಿರುವುದೇ ನಿಜವಾದರೆ, ರಾವಣ ಪುತ್ರನಾದ ಇಂದ್ರಜಿತ್ತನ್ನು ಧ್ವಂಸಮಾಡು.

\medskip
\noindent\textbf{೧೧೯. ಧರ್ಮೋ ರಕ್ಷತಿ ರಕ್ಷಿತಃ} \hfill ಪುಟ \pageref{225}

ನಮ್ಮಿಂದ ರಕ್ಷಿಸಲ್ಪಟ್ಟ ಧರ್ಮವು ನಮ್ಮನ್ನು ರಕ್ಷಿಸುವುದು.

\medskip
\noindent\textbf{೧೨೦. ಧರ್ಮೋ ವಿಗ್ರಹವಾನಧರ್ಮವಿರತಿಂ} \hfill ಪುಟ \pageref{138}

\hfill ದಶಾವತಾರ ಸ್ತೋತ್ರ

ಮೂರ್ತಿವೇತ್ತ ಧರ್ಮಸ್ವರೂಪನಾಗಿರುವ ಆ ಪ್ರಶಸ್ತಧನುರ್ಧಾರಿಯಾದ ಶ್ರೀರಾಮನು ನಮ್ಮಲ್ಲಿಯ ಅಧರ್ಮವೆಲ್ಲವೂ ಪೂರ್ತಿ ಅಳಿದುಹೋಗುವಂತೆ ಮಾಡಲಿ.

\medskip
\noindent\textbf{೧೨೧. ಧರ್ಮೋ ವಿಶ್ವಸ್ಯ ಜಗತಃ ಪ್ರತಿಷ್ಠಾ} \hfill ಪುಟ \pageref{182},\pageref{207b

\hfill ಮಹಾನಾರಾಯಣೋಪನಿಷತ್ ೧೨-೬

ಧರ್ಮವು ಸಮಸ್ತವಾದ ಜಗತ್ತಿಗೂ ಧೃಢವಾದ ಧಾರಕವಾಗಿದೆ.

\medskip
\noindent\textbf{೧೨೨. ಧ್ವನೇರಂತರ್ಗತಂ  ಜ್ಯೋತಿಃ} \hfill ಪುಟ \pageref{87}

\hfill ಯೋಗಶಿಖೋಪನಿಷತ್ ೬-೨೧

`ಅನಾಹತಸ್ಯ ಶಬ್ದಶ್ಯ ತಸ್ಯ ಶಬ್ದಶ್ಯ ಯೋ ಧ್ವನಿಃ'  ಎಂಬುದು ಪೂರ್ವಭಾಗ.

ಲೋಕದಲ್ಲಿ ನಾವು ಕೇಳುವ ಶಬ್ದಗಳೆಲ್ಲವೂ ಎರಡು ವಸ್ತುಗಳ ಪರಸ್ಪರ ಆಘಾತದಿಂದ ಹುಟ್ಟುತ್ತವೆ. ಆದ್ದರಿಂದ ಅವುಗಳು ಆಹತ ಶಬ್ದಗಳು. ಹೀಗಾಲ್ಲದೆ ಯೋಗಿಗಳ ಅಂತಃಶ್ರವಣಗೋಚರವಾಗುವ ಶಬ್ದಗಳೂ ಉಂಟು. ಅವು ಅನಾಹತ ಶಬ್ದಗಳು. ಪ್ರಣವವೂ ಅಂತಃಶ್ರವಣಗೋಚರವಾಗುವ ಅನಾಹತ ಶಬ್ದವೇ ಆಗಿದೆ. ಅಂತಹ ಯೋಗಿಮಾತ್ರ ವೇದ್ಯವಾದ ಪ್ರಣವಧ್ವನಿ-ನಾದಕ್ಕೂ ಒಳಗಡೆ ಸೂಕ್ಷ್ಮತಮನಾಗಿರುವ (ಜ್ಞಾನ ಚಕ್ಷುಸ್ಸಿಗೆ ಮಾತ್ರ ಎಟುಕುವ) ಪ್ರಕಾಶರೂಪಿಯಾದ ಜೀವನುಂಟು. ಅವನ ಜೊತೆಗೆ ಸೇರಿಕೊಂಡಿರುವ ಮನಸ್ಸು ಜೀವದ ಬಂಧ ಮೋಕ್ಷಗಳಿಗೆ ಕಾರಣವಾಗಿದೆ. ಅಂತಹ ಮನಸ್ಸೆಂಬ ಅಂತಃಕರಣತತ್ತವು ಜೀವಿಯನ್ನು ಸಂಸಾರಕ್ಕೆ ಪ್ರೇರಿಸದೆ ಯಾವ ಪರಮಧಾಮದಲ್ಲಿ ಲೀನಗೊಳ್ಳುವುದೋ ಅದೇ ವಿಷ್ಣುವಿನ ಪರಮಪದವಾಗಿದೆ.

\medskip
\noindent\textbf{೧೨೩. ನಕಾರಂ ಪ್ರಾಣನಾಮಾನಂ} $\star$ \hfill ಪುಟ \pageref{17}

\hfill ಶಾರದಾತಿಲಕ ೧

\medskip
\noindent\textbf{೧೨೪. ನಕ್ಷತ್ರಾಣಾಮಹಂ ಶಶೀ} \hfill ಪುಟ \pageref{98}

\hfill ಭಗವದ್ಗೀತಾ ೧೦-೨೧

ಅರ್ಜುನ! ನಾನು ನಕ್ಷತ್ರರಾಜನೆನಿಸುವ ಚಂದ್ರಸ್ವರೂಪನಾಗಿರುವೆನು.

\medskip
\noindent\textbf{೧೨೫. ನತೇ ವಾಗನೃತಾ ಕಾವ್ಯೇ} \hfill ಪುಟ \pageref{209}

\hfill ರಾಮಾಯಣ ಬಾಲಕಾಂಡ ೨-೩೫

(ಕ್ರೌಂಚವನ್ನು ಕೊಂದ ವ್ಯಾಧನನ್ನು ಶಪಿರಿ `ಕಿಮಿದಂ ವ್ಯಾಹೃತಂ ಮಯಾ' ಎಂದು ಚಿಂತಿಸುತ್ತಾ ವಾಲ್ಮೀಕಿಯು ಆಶ್ರಮಕ್ಕೆ ಹಿಂದಿರುಗಲು ಲೋಕಕರ್ತನಾದ ಬ್ರಹ್ಮನು ಅಲ್ಲಿಗೆ ಬಂದು ರಾಮಾಯಣವನ್ನು ರಚಿಸುವಂತೆ ಆದೇಶಿಸಿದನು. ಮತ್ತು) ನಿನ್ನ ಕಾವ್ಯದಲ್ಲಿ ಯಾವ ಮಾತೂ ಸುಳ್ಳಾಗುವುದಿಲ್ಲ ಎಂದು ವಾಲ್ಮೀಕಿಗೆ ವರವನ್ನಿತ್ತನು.

\medskip
\noindent\textbf{೧೨೬. ನಪುತ್ರ ಮರಣಂ ಕಿಂಚಿತ್} \hfill ಪುಟ \pageref{180}

ರಾಮನು ರಾಜ್ಯವಾಳುತ್ತಿರಲಾಗಿ ಯಾಗುದೇ ಎಡೆಯಲ್ಲೂ ಜನರು ತಮ್ಮೆದುರಿಗೆ ಮಕ್ಕಳ ಸಾವನ್ನು ಕಾಣಲಿಲ್ಲ. ಹಾಗೆಯೇ ಅಗ್ನಿಯಿಂದಾಗಲಿ, ಜ್ವರಾದಿಯೋಗಗಳಿಂದಾಗಲಿ ಭಯಕ್ಕೆ ಅವಕಾಶವಿರಲಿಲ್ಲ. ನಗರಗಳೂ, ನಾನಾನಗರಗಳಿಂಡ ಕೂಡಿದ ರಾಷ್ಟ್ರಗಳೂ ಧನಧಾನ್ಯಸಮೃದ್ಧಿಯಿಂದ ಕೂಡಿದ್ದವು. ಎಲ್ಲ ಜನರೂ ಕೃತಯುಗದಲ್ಲಿ ಹೇಗೋ ಹಾಗೆ ನಿತ್ಯತೃಪ್ತರಾಗಿ ಬಾಳುತ್ತಿದ್ದರು.

\medskip
\noindent\textbf{೧೨೭. ನ ರಾಮಸ್ತಪಸಾ ದೇವಿ} \hfill ಪುಟ \pageref{191}

\hfill ರಾಮಾಯಣ ಸುಂದರಕಾಂಡ ೨೦-೩೪

(ನ ಧನೇನ ಮಯಾ ತುಲ್ಯತ್ತೇಜಸಾ ಯಶಸಾಪಿ ವಾ||)

ಓ! ದೇವಿ ಸೀತೆಯೇ! ತಪಸ್ಸಿನಿಂದಲಾಗಲಿ, ಬಲದಿಂದಾಗಲಿ, ಪರಾಕ್ರಮದಿಂದಾಗಲಿ, ಹಣದಿಂದಾಗಲಿ, ತೇಜಸ್ಸಿನಿಂದಾಗಲಿ, ಯಶಸ್ಸಿನಿಂದಾಗಲಿ ರಾಮನು ನನಗೆ ಸರಿಸಮನಾಗಲಾರ. 

\medskip
\noindent\textbf{೧೨೮. ನ ವೇದಂ ವೇದಮಿತ್ಯಾಹುಃ}  $\star$ \hfill ಪುಟ \pageref{13}

\hfill ಯೋಗೋಪನಿಷತ್

\medskip
\noindent\textbf{೧೨೯. ನ ಶಕ್ಯಃ ಚಕ್ಷುಷಾ ದ್ರಷ್ಟುಂ} \hfill ಪುಟ \pageref{246}

ಅತ್ಯಂತ ಸೂಕ್ಷನಾಗಿ ದೇಹದ ಸರ್ವಾಂಗವನ್ನೂ ವ್ಯಾಪಿಸಿ ಬೆಳಗುತ್ತಿರುವ ಆತ್ಮ ಈ ಸರ್ವಪ್ರಾಣಿಸಾಧಾರಣವಾದ ಚರ್ಮಚಕ್ಷುಸ್ಸಿನಿಂದ ಕಾಣಲಸಾಧ್ಯನು. ಜ್ಞಾನರೂಪವೂ ಅಥವಾ ತಪೋರೂಪವೂ ಆದ ಕಣ್ಣುಳ್ಳವರಿಗೆ ಮಾತ್ರ ಗೋಚರನು. (ಜ್ಞಾನಯೋಗ - ಭಕ್ತಿಯೋಗ - ಕರ್ಮಯೋಗಗಳನ್ನು ಹೊಂದಿರತಕ್ಕ ಯೋಗಿ ಜನರಿಗೆ ಮಾತ್ರ ಗೋಚರನು)

\medskip
\noindent\textbf{೧೩೦. ನಷ್ಟೋ ಮೋಹಃ ಸ್ಮೃತಿರ್ಲಬ್ಧಾ} \hfill ಪುಟ \pageref{91},\label{104}

\hfill ಭಗವದ್ಗೀತಾ ೧೮-೭೩

ಎಲೈ ಅಚ್ಯುತನೇ ನಿನ್ನ ಪ್ರಸಾದದಿಂದ ನನಗೆ ಮೋಹವು ಕಳೆಯಿತು. ಪೂರ್ವಸ್ಮೃತಿಯು ದೊರಕಿತು. ನಾನು ಯಾರು? ನನ್ನ ಕರ್ತವ್ಯವೇನು? ಎಂಬ ಬಗ್ಗೆ ಜಾಗೃತಿಯುಂಟಾಗಿದೆ. ನನಗೆ ಸಂದೇಹವು ದೂರವಾಗಿದೆ. ನಿನ್ನ ಮಾತಿನಂತೆ ನಡೆಯುತ್ತೇನೆ.

\medskip
\noindent\textbf{೧೩೧. ನಾಽವೇದವಿನ್ಮನುತೇ ತಂ} \hfill ಪುಟ \pageref{106}

\hfill ತೈತ್ತಿರೀಯ ಬ್ರಾಹ್ಮಣ ೨-೯-೧೭

ಆ ಬ್ರಹ್ಮತತ್ತ್ವವನ್ನು ಸೃಷ್ಟಿ ಇತಿಹಾಸಮೂಲರೂಪವಾದ  ವೇದದ ಜ್ಞಾನಿಯಲ್ಲದವನು ಅರಿಯಲಾರನು.

\medskip
\noindent\textbf{೧೩೨. ನಿತ್ಯಾನಂದವಪುರ್ನಿರಂತರಗಲತ್} $\star$ \hfill ಪುಟ \pageref{31}

\hfill ಶಾರದಾತಿಲಕ ೧-೧

\medskip
\noindent\textbf{೧೩೩. ನಿಮಿತ್ತಮಾತ್ರಂ ಭವ ಸವ್ಯಸಾಚಿನ್} \hfill ಪುಟ \pageref{214a}

\hfill ಭಗವದ್ಗೀತೆ ೧೧-೩೩

"ಎಲೈ ಸವ್ಯಸಾಚಿಯೇ! ಅರ್ಜುನನೇ! (ನನ್ನಿಂದಲೇ ಇವರೆಲ್ಲರೂ ಕೊಲ್ಲಲ್ಪಟ್ಟಿದ್ದಾರೆ) ನೀನು (ಈ ಶತ್ರು ಸಂಹಾರಕಾರ್ಯದಲ್ಲಿ) ನಿಮಿತ್ತಮಾತ್ರವಾಗಿರು" ಎಂದು ಶ್ರೀಕೃಷ್ಣನು ಅರ್ಜುನನಿಗೆ ಉಪದೇಶಿಸುತ್ತಾನೆ.

\medskip
\noindent\textbf{೧೩೪. ನಿಯತಾತ್ಮಾ ಮಹಾವೀರ್ಯಃ} \hfill ಪುಟ \pageref{247b}

\hfill ರಾಮಾಯಣ ಬಾಲಕಾಂಡ ೧-೮

`ಇಕ್ಷ್ವಾಕುವಂಶಪ್ರಭವೋ' ಇತ್ಯಾದಿ ಗ್ರಂಥದ ಭಾವಾನುವಾದವನ್ನು ಗಮನಿಸಿ.

\medskip
\noindent\textbf{೧೩೫. ನಿರ್ಭರೋ ನಿರ್ಭಯೋಽಸ್ಮಿ} \hfill ಪುಟ \pageref{115}

\hfill ನ್ಯಾಸವಿಂಶತಿ ೨೨

\begin{shloka}
ಸಂಸಾರಾವರ್ತವೇಗಪ್ರಶಮನಶುಭಕೃದ್ದೇಶಿಕಪ್ರೇಕ್ಷಿತೋಽಹಂ\\
ಸಂತ್ಯಕ್ತೋಽನ್ಯೈರುಪಾಯೈರನುಚಿತಚತಿತೇಷ್ವದ್ಯ ಶಾಂತಾಭಿಸಂಧಿಃ |\\
ನಿಶ್ಯಂಕಸ್ತತ್ವದೃಷ್ವ್ಯಾ ನಿರುಪಧಿಕದಯಂ ಪ್ರಾಪ್ಯ ಸಂರಕ್ಷಕಂ ತ್ವಾಂ\\
ನ್ಯಸ್ಯ ತ್ವತ್ಪಾದಪದ್ಮೇ ವರದ ನಿಜಭರಂ ನಿರ್ಭರೋ ನಿರ್ಭಯೋಽಸ್ಮಿ \\
\end{shloka}

ಸ್ವಾಮಿನ್! ನಾನೀಗ ಸಂಸಾರವೆಂಬ ಸಾಗರದ ಸುಳಿಗಳ ವೇಗವನ್ನಡಗಿಸುವಂತ ಮಾರ್ಗದರ್ಶೀ ಗುರುವಿನಿಂದ ಕಟಾಕ್ಷಿಸಲ್ಪಟ್ಟಿರುವೆನು. ನನ್ನ ಉದ್ಧಾರಕ್ಕೆ ಗುರುಕೃಪೆಯನ್ನುಳಿದು ಮತ್ತ್ತಾವುದೇ ಉಪಾಯಗಳನ್ನೂ ಕೈಗೊಳ್ಳುತ್ತಿಲ್ಲ. ಭಗವದನುಗ್ರಹ ಪ್ರತಿಕೂಲವಾದಂತಹ ಕಾಯ - ವಾಙ್ - ಮನೋವ್ಯಾಪಾರಗಳಲ್ಲಿ ತೊಡಗಬಾರದೆಂಬ ನಿಶ್ಚಯಕ್ಕೆ ಬಂದಿದ್ದೇನೆ. ಅವ್ಯಾಜಕರುಣಾಮನಾಗಿ ರಕ್ಷ್ಣೋದ್ಯುಕ್ತನಾದ ನಿನ್ನನ್ನು ಶರಣುಹೊಂದಿ, ತತ್ತ್ವಾರ್ಥದರ್ಶನದಿಂದ ಶಂಕೆಯನ್ನೆಲ್ಲಾ ತೊರೆದಿದ್ದೇನೆ. ನಿನ್ನಡಿದಾವರೆಯಲ್ಲಿ ನನ್ನ ಆತ್ಮರಕ್ಷಣೆಯ ಹೊಣೆಯನ್ನರ್ಪಿಸಿ, ವರಪ್ರದನಾದ ನನ್ನೊಡೆಯ! ಭಾರವೆಲ್ಲಾ ಕಳೆದು ನಿರ್ಭೀತನಾಗಿ ಬಾಳುವೆನಪ್ಪಾ!

\medskip
\noindent\textbf{೧೩೬. ನೃತ್ತಾವಸಾನೇ ನಟರಾಜರಾಜಃ} $\star$ \hfill ಪುಟ \pageref{30}

\hfill ನಂದಿಕೇಶ್ವರ ಕಾರಿಕಾ ೧

\medskip
\noindent\textbf{೧೩೭. ನೇಯಂ ಸಜ್ಜನ ಸಂಗೇ ಚಿತ್ತಂ} \hfill ಪುಟ \pageref{83}

\hfill ಭಜಗೋವಿಂದ ಸ್ತೋತ್ರ

ನಮ್ಮ ಚಿತ್ತವು ಸಜ್ಜನರ ಸಹವಾಸದಲ್ಲಿ ಒಯ್ಯಲ್ಪಡಬೇಕು.

\medskip
\noindent\textbf{೧೩೮. ನ್ಯಸ್ಯತಾಂ ಕಲಶಸ್ತಾತ} \hfill ಪುಟ \pageref{198}

\hfill ರಾಮಾಯಣ ಬಾಲಕಾಂಡ ೨-೬

`ಅಪ್ಪಾ! ಕಲಶವನ್ನು ಕೆಳಕ್ಕೆ ಇಡು. ನನಗೆ (ಸ್ನಾನಾರ್ಥವಾದ) ನಾರುಮಡಿಯನ್ನು ಕೊಡು. ಇಲ್ಲಿಯೇ ಎದುರಿಗಿರುವ ಉತ್ಕೃಷ್ಟವಾದ ತಮಸಾತಿರ್ಥದಲ್ಲಿ ಅವಗಾಹನ (ಇಳಿದು ಸ್ನಾನ)ವನ್ನು ಮಾಡುವೆನು.

\medskip
\noindent\textbf{೧೩೯. ಪಂಚೇಂದ್ರಿಯಸ್ಯ ದೇಹಸ್ಯ} \hfill ಪುಟ \pageref{139}

\hfill ಚತುರ್ವಿಂಶತಿಮತ

ಐದೂ ಬಗೆಯ ದೇಹಗತವಾದ ಇಂದ್ರಿಯಗಳ, ದೇಹದ, ಬುದ್ಧಿಯ, ಮನಸ್ಸಿನ ಅಂತೆಯೇ ದ್ರವ್ಯ - ದೇಶ - ಕ್ರಿಯೆಗಳ ಶುದ್ಧಿಯೆಂಬುದು ಏನುಂಟೋ ಅದೇ ಆಚಾರವೆನಿಸುವುದು. 

\medskip
\noindent\textbf{೧೪೦. ಪದ್ಭ್ಯಾಂ ಉದಾರೇಣ ಶಿಶ್ಠ್ಯಾ} \hfill ಪುಟ \pageref{26}

ಕಾಲುಗಳಿಂದಾಗಲೀ ಹೊಟ್ಟೇಯಿಂದಾಗಲೀ, ಜನನೇಂದ್ರಿಯದಿಂದಾಗಲೀ ಹಗಲು ರಾತ್ರಿಗಳಲ್ಲಿ ಯಾವ ಪಾಪವೊದಗುವುದೋ ಅದನ್ನು ಆಯಾ ಕಾಲ ಶರೀರಿಯಾದ ಪುರುಷನೇ ಇಲ್ಲಮಾಡಲಿ.

\medskip
\noindent\textbf{೧೪೧. ಪರಂ ಕೌತೂಹಲಂ ಹಿ ಮೇ} \hfill ಪುಟ \pageref{163}

\hfill ರಾಮಾಯಣ ಬಾಲಕಾಂಡ ೧-೫

`ಏತದಿಚ್ಛಾಮ್ಯಹಂ' ಇತ್ಯಾದಿ ಗ್ರಂಥಭಾಗದ ಭಾವಾನುವಾದವನ್ನು ಗಮನಿಸಿ.

\medskip
\noindent\textbf{೧೪೨. ಪರಂ ವಿಸ್ಮಯಮಾಗತಾಃ} \hfill ಪುಟ \pageref{211a}

\hfill ರಾಮಾಯಣ ಬಾಲಕಾಂಡ ೪-೧೬

`ತಚ್ಛೃತ್ವಾ' ಇತ್ಯಾದಿ ಗ್ರಂಥದ ಭಾವಾನುವಾದವನ್ನು ಗಮನಿಸಿ.

\medskip
\noindent\textbf{೧೪೩. ಪಶುಧಾನ್ಯ ಹಿರಣ್ಯ ಸಂಪದಾ} \hfill ಪುಟ \pageref{227}

ರಾಷ್ಟ್ರವೆನಿಸಿಕೊಳ್ಳಬೇಕಾದರೆ, ಪಶು (ದನಕರುಗಳು), ದವಸ ಧಾನ್ಯಗಳು, ಸುವರ್ಣ ರಜತಾದಿ ಉತ್ತಮ ಲೋಹಗಳು, ಅವುಗಳ ಸಂಪತ್ಸಮೃದ್ಧಿಯಿಂದ ಬೆಳಗುತ್ತಿರಬೇಕು. 

ದೇಶವೆನಿಸಿಕೊಳ್ಳಬೇಕಾದರೆ, ಆಳುವ ಪ್ರಭುವಿಗೆ ದಂಡನೀತಿಯ ಮತ್ತು ಖಜಾನೆಯ ವೈಭವವನ್ನು ತಂದು ಕೊಡಬೇಕು.

ಚತುರ್ವಣ ಮತ್ತು ಚತುರಾಶ್ರಮಗಳ ಲಕ್ಷ್ಣಣದಿಂದ ಕೂಡಿದ ಜನಗಳ ವಾಸ ಸ್ಠಾನವು ಜನಪದ... ಅಥವ ಅಂತಹ ಜನರಿಗೆ ಜೀವನಕ್ಕೆ ಅಗತ್ಯವಾದ ದ್ರವ್ಯೋತ್ಪತ್ತಿಗೆ ಸ್ಠಾನವಾದದ್ದು ಜನಪದ.

\medskip
\noindent\textbf{೧೪೪. ಪಾಠ್ಯೇ ಗೇಯೇ ಚ ಮಧುರಂ} \hfill ಪುಟ \pageref{164a},\pageref{250}

\hfill ರಾಮಾಯಣ ಬಾಲಕಾಂಡ ೪-೮

ಸಾಮಾನ್ಯವಾದ ಸಾಹಿತ್ಯಮಾತ್ರೋಚ್ಚಾರಣೆಯಲ್ಲೂ, ಸ್ವರರಾಗ ಯೋಜನೆಯಲ್ಲೂ, ಎರಡೂ ಬಗೆಯಲ್ಲಿ ಇಂಪುಗೊಡುವಂಥದೂ, ದ್ರುತ - ಮಧ್ಯ - ವಿಲಂಬಿತಗಳೆಂಬ ಮೂರು ಪರಿಮಿತಿಗಳಿಂದ ಕೂಡಿ ಸೊಬಗುಗೊಂಡಿರುವುದೂ, ಗಾನಕ್ಕೆ ಸಂಬಧಿಸಿದ ಶುದ್ಧ - ವಿಕೃತ - ಸಂಕೀರ್ಣಾದಿ ಹದಿನೆಂಟು ಬಗೆಯ ಜಾತಿಗಳ ಪೈಕಿ ಪ್ರಧಾನವಾದ ಏಳು ಬಗೆಯ ಜಾತಿಗಳಿಂದ ಕಳಚಿಕೊಳ್ಳದೇ ಅವುಗಳ ಮಿತಿಗೆ ಒಳಪಟ್ಟಿರುವುದೂ, ವೀಣಾವೇಣು ಮೃದಂಗಾದಿ ವಾದ್ಯಗಳೆಲ್ಲದರ  ಹೊಂದಾಣಿಕೆಯುಳ್ಳದ್ದೂ ಆಗಿರುವಂತೆ ಹಾಡಿದರು.

\medskip
\noindent\textbf{೧೪೫. ಪಾದಬದ್ಧೋಽಕ್ಷರಮಃ} $\star$ \hfill ಪುಟ \pageref{20a},\pageref{157a},\pageref{201},\pageref{213},\pageref{242}

\hfill ರಾಮಾಯಣ ಬಾಲಕಾಂಡ ೨-೧೮

\medskip
\noindent\textbf{೧೪೬. ಪಾಪಾತ್ಮನಾ ಕೃತಂ ಕಷ್ಟಂ} \hfill ಪುಟ \pageref{203a}

ಪಾಪಸ್ವಭಾವದವನೂ ವೈರವೇ ನೆಲೆಗೊಂಡ ಬುದ್ಧಿ ಪ್ರವೃತ್ತಿಯುಳ್ಳವನೂ ಆದ ಆ ಕಿರಾತನಿಂದ ನಡೆದ ಕಾರ್ಯವು ಅತ್ಯಂತ ದಾರುಣವಾದುದು. ಏಕೆಂದರೆ ಇಂಪಾದ ಧ್ವನಿ ಮಾಡುತ್ತ ತನ್ನ ಸಂಗಾತಿಯೊಡನಿದ್ದ ಆ ಕ್ರೌಂಚಪಕ್ಷಿಯನ್ನು ನಿಷ್ಕಾರಣವಾಗಿ (ಆಹಾರಕ್ಕಾಗಿಯೂ ಅಲ್ಲದೇ) ಕೊಂದನಷ್ಟೆ.

\medskip
\noindent\textbf{೧೪೭. ಪಿಬ ರೇ ರಾಮರಸಂ} \hfill ಪುಟ \pageref{187}

ಎಲೈ! ಜೀವವೇ! ರಾಮರಸಾನುಭವವನ್ನು ಮಾಡು. 

\medskip
\noindent\textbf{೧೪೮. ಪುಣ್ಯೇ ಹಿಮವತಃ ಪಾರ್ಶ್ವೇ} \hfill ಪುಟ \pageref{118}

\hfill ಮಹಾಭಾರತ

ಪುಣ್ಯವಾದ ಹಿಮಾಲಯದ ಪರಿಸರ ಅದರ ಮಧ್ಯೆ ನೆಲೆಸಿರುವ ಬೆಟ್ಟದ ಗುಹೆಯ ಆವಾಸದಲ್ಲಿ ಸತ್ಯವತೀಸುತನಾದ ವ್ಯಾಸನು ಧರ್ಮಾತ್ಮನಾಗಿ ದೇಹವನ್ನು ಶೋಧಿಸಿಕೊಂಡು ದರ್ಭೆಯ ಹಾಸಿನ ಮೇಲೆ ಆಸೀನನಾದನು. ಶುಚಿಯಾಗಿಯೂ ನಿಯಮದಿಂದ ಕೂಡಿದವನಾಗಿಯೂ ಶಾಂತಾತ್ಮನಾಗಿಯೂ ತಪಸ್ಸಿನಲ್ಲಿ ನೆಲೆನಿಂತವನಾಗಿಯೂ ಆತನು ಆಗಿದ್ದನು. ಧರ್ಮದಿಂದ ಭಾರತ ಇತಿಹಾಸದ ಗತಿಯನ್ನು ಯಥಾರ್ಥವಾಗಿ ಅವಲೋಕಿಸಿದನು. ಯೋಗಭೂಮಿಯನ್ನು ಪ್ರವೇಶಿಸಿ ಜ್ಞಾನದೃಷ್ಟಿಯಿಂದ ಅವನು ಕೊನೆಯವರೆಗೂ ಎಲ್ಲವನ್ನು ನೋಡಿದನು. ತಪಸ್ಸು ಮತ್ತು ಬ್ರಹ್ಮಚರ್ಯೆಯ ಬಲದಿಂದ ಸನಾತನವಾದ ವೇದವನ್ನು ವಿಭಜಿಸಿ (ವೇದ ಸಾರವಾದ) ಪುಣ್ಯರೂಪವಾದ ಈ ಇತಿಹಾಸವನ್ನು ರಚಿಸಿದನು.

\medskip
\noindent\textbf{೧೪೯. ಪುಸ್ತಕಸ್ಠಾತು ಯಾ ವಿದ್ಯಾ} \hfill ಪುಟ \pageref{114}

ನಿಸರ್ಗಾವಲಂಬಿಯಾಗಿ ತಿಳುವಳಿಕೆ ಹೊಂದದೇ ಪುಸ್ತಕಾವಲಂಬಿಯಾಗಿರುವ ಅರಿವಿನಿಂದ ಮೂಢನು ಮಾತ್ರ ಮೋಸಹೋಗುವನು.

\medskip
\noindent\textbf{೧೫೦. ಪೂರ್ಣಸ್ಯ ಪೂರ್ಣಮಾದಾಯ} \hfill ಪುಟ \pageref{23},\pageref{48a},\pageref{268}

\hfill ಶಾಂತಿಪಾಠ

ನಿರುಪಾಧಿಕವಾಗಿ ಸರ್ವತ್ರ ಪೂರ್ಣನಾದ ಪರಮಾತ್ಮನ ಸ್ವಭಾವವಾದ ಪೂರ್ತಿಯೆಂಬುದನ್ನು ಪಡೆದು ನಂತರ ಮತ್ತೆ ಪೂರ್ಣವಾಗಿಯೇ ಉಳಿಯುತ್ತದೆ.

\medskip
\noindent\textbf{೧೫೧. ಪೌಲಸ್ತ್ಯವಧಮಿತ್ಯೇವ} \hfill ಪುಟ \pageref{207c}

\hfill ರಾಮಾಯಣ ಬಾಲಕಾಂಡ ೪-೭

ವಾಲ್ಮೀಕಿಯು ಸಮಗ್ರವಾದ ರಾಮಾಯಣ ಕಾವ್ಯವನ್ನು ಸೀತೆಯ ಮಹತ್ತಾದ ಚರಿತ್ರೆಯೆಂದೂ ಪೌಲಸ್ತ್ಯ - ರಾವಣವಧವೆಂದೂ ಭಾವಿಸಿಕೊಂಡು ಆ ಅಂಶವನ್ನು ಸ್ಪಷ್ಟವಾಗುವಂತೆ ಧೃಢವ್ರತಕೈಗೊಂಡು ನಿರ್ಮಾಣಮಾಡಿದನು. (ಉಳಿಯ ಬೇಕಾದುದನ್ನೂ ಅಳಿಯಬೇಕಾದುದನ್ನೂ ಖಚಿತವಾಗಿ ಸ್ಪಷ್ಟಗೊಳಿಸಿದನು).

\medskip
\noindent\textbf{೧೫೨. ಪ್ರತ್ಯಗಾನಂದಂ ಬ್ರಹ್ಮಪುರುಷಂ} \hfill ಪುಟ \pageref{106a}

\hfill ಲಘುನಾರಾಯಣೋಪನಿಷತ್

ಸೃಷ್ಟಿಯ ಹಿಂಬದಿಯಲ್ಲಿರುವ ಪ್ರತ್ಯಗ್ರೂಪಿಬ್ರಹ್ಮವು ಆನಂದೈಕ ಸ್ವರೂಪವಾದುದು. ಅದೇ ಬೀಜರೂಪವಾಗಿ ಎಲ್ಲೆಲ್ಲೂ ಪೂರ್ಣವಾಗಿರುವುದು. (ಅದೇ ಮೊದಲಿಗೆ ಅರ್ಧಮಾತ್ರೆಯಾಗಿದ್ದು) `ಅ, ಉ, ಮ' ಎಂಬ ಪ್ರಣವಾತ್ಮಕವಾಗಿ ಶಬ್ದ ಬ್ರಹ್ಮವೂ ಆಗಿರುವುದು. ಈ ಮೂರು ವರ್ಣಗಳ ಒಂದುಗೂಡಿದ - ಸಂಹಿತವಾದ ರೂಪವೇ `\includegraphics{Om-PNG.eps}' ಎಂಬುದು. ಆತ್ಮಯೋಗಾಭ್ಯಾಸ ಪ್ರವೃತ್ತನಾದವನು (ಯೋಗಿಯು) ಆ `\includegraphics{Om-PNG.eps}' ಎಂಬುದನ್ನು ನೈಜ ಸ್ವರೂಪದಲ್ಲಿ  ಉಚ್ಚಾರಣೆ ಮಾಡಿದ್ದೇ ಆದರೆ, ಜನ್ಮಜರಾಮರಣಾತ್ಮಕವಾದ ಸಂಸಾರದ ಕಟ್ಟಿನಿಂದ ಬಿಚ್ಚಿಕೊಳ್ಳುವನು. 

\medskip
\noindent\textbf{೧೫೩. ಪ್ರಮಾದಂ ವೈ ಮೃತ್ಯುಮಹಂ} \hfill ಪುಟ \pageref{126}

\hfill ಮಹಾಭಾರತ ಉದ್ಯೋಗಪರ್ವ ೪೨-೪

\begin{shloka}
`ಪ್ರಮಾದಾದ್ವೆ ಅಸುರಾಃ ಪರಾಭವನ್\\
ಅಪ್ರಮಾದಾದ್ ಬ್ರಹ್ಮಭೂತಾಃ ಸುರಾಶ್ಚ||' ಎಂಬುದು ಉತ್ತರಾರ್ಧ
\end{shloka}

ನಾನು (ಸನತ್ಸುಜಾತಮುನಿ) ಪ್ರಮಾದವನ್ನೇ ಮೃತ್ಯುವೆಂದು ಹೇಳುವೆನು. ಅಮೃತತ್ವವೆಂದರೆ ಯಾವಗಲೂ ಅಪ್ರಮತ್ತನಾಗಿರುವುದೇ ಆಗಿದೆ. (ಅಸುರರು ಪ್ರಮಾದದಿಂದಲೇ ಪರಾಭವ ಹೊಂದಿದರು. ಸುರರು ಅಪ್ರಮಾದದಿಂದಲೇ ಬ್ರಹ್ಮಭಾವ ತಲುಪಿದರು).

\medskip
\noindent\textbf{೧೫೪. ಪ್ರಯೋಗೈರ್ಯುಕ್ತಿಕಲ್ಪಿತೈಃ} \hfill ಪುಟ \pageref{136a}

`ನಾನಾದೇಶಿಕವಕ್ತ್ರಸ್ಥೈಃ' ಎನ್ನುವುದು ಇದರ ಪೂರ್ವವಾಕ್ಯ.
ಸೃಷ್ಟಿತತ್ತ್ವಜ್ಞರಾದ ಬಗೆಬಗೆಯ ಆಚಾರ್ಯರುಗಳ ಮುಖಗಳಿಂದ ಹೊರಬರತಕ್ಕ ಯುಕ್ತಿಕಲ್ಪಿತವಾದ ಪ್ರಯೋಗವಿಧಾನಗಳಿಂದಲೂ (ಸತ್ಯವನ್ನು ಸಾಧಿಸಬಹುದು.)

\medskip
\noindent\textbf{೧೫೫. ಪ್ರಶಶಂಸುಃ ಪ್ರಶಸ್ತವ್ಯೌ} \hfill ಪುಟ \pageref{168},\pageref{243a}

\hfill ರಾಮಾಯಣ ಬಾಲಕಾಂಡ ೪-೧೭

`ಸಾಧು ಸಾಧ್ವಿತಿ' ಇತ್ಯಾದಿ ಗ್ರಂಥದ ಭಾವಾನುವಾದವನ್ನು ಗಮನಿಸಿ.

\medskip
\noindent\textbf{೧೫೬. ಪ್ರಾಚೀನಾಗ್ರೇಷು ದರ್ಭೇಷು} \hfill ಪುಟ \pageref{204a},\pageref{231}

\hfill ರಾಮಾಯಣ ಬಾಲಕಾಂಡ ೩-೨

ವಾಲ್ಮೀಕಿಮಿನಿಯು ಸರಿಯಾಗಿ ಶುದ್ಧಾಚಮನ ಮಾಡಿ ಬದ್ಧಾಂಜಲಿಪುಟನಾಗಿ ಪೂರ್ವಾಗ್ರವಾಗಿ ಹಾಸಿದ ದರ್ಭೆಗಳ ಮೇಲೆ ನಿಂತು ಸತ್ಯಕ್ಕೆ ಒಯ್ಯತಕ್ಕ ಧರ್ಮವನ್ನೇ ಸಾಧನವಾಗಿ ಮಾಡಿಕೊಂಡು ರಾಮಜೀವನದ ನಡೆಯನ್ನು ಒಳಗೇ ಅನ್ವೇಷಣೆ ಮಾಡಿದನು.

\medskip
\noindent\textbf{೧೫೬. ಪ್ರಾಣಾಪಾನೌ ಸಮಾವಾಸ್ತಾಂ} \hfill ಪುಟ \pageref{151},\pageref{161},\pageref{218},\pageref{222a},\pageref{253}

\hfill ಮಹಾಭಾರತ

ಶ್ರೀರಾಮಚಂದ್ರನು ರಾಜ್ಯಪಾಲನೆ ಮಾಡುತ್ತಿರಲಾಗಿ ಸಮಸ್ತದೇಹಿಗಳ ದೇಹಗಳಲ್ಲಿ ಜೀವನಧಾರಕಗಳಾದ ಪ್ರಾಣ ಮತ್ತು ಅಪಾನಗಳೆಂಬ ಪ್ರಧಾನವಾದ ವಾಯುಗಳೆರಡೂ ಸಾಮ್ಯಹೊಂದಿ ಯಾರ ಜೀವನದಲ್ಲೂ ವೈಷಮ್ಯವೊದಗದಂತೆ ಸ್ಥಿತವಾಗಿದ್ದವು.

\medskip
\noindent\textbf{೧೫೭. ಪಾಣಾವಾಮಲಕಂ ಯಥಾ} $\star$ \hfill ಪುಟ \pageref{232}

\hfill ರಾಮಾಯಣ ಬಾಲಕಾಂಡ ೩-೬

\medskip
\noindent\textbf{೧೫೮. ಪ್ರಾತಃ ಸ್ಮರಾಮಿ ಹೃದಿ} $\star$ \hfill ಪುಟ \pageref{146a}

\hfill ಶಂಕರಾಚಾರ್ಯಕೃತ -- ಪ್ರಾತಃಸ್ಮರಣಸ್ತೋತ್ರ ೧

\medskip
\noindent\textbf{೧೫೯. ಪ್ರಾಪ್ತ ರಾಜ್ಯಸ್ಯ} \hfill ಪುಟ \pageref{188a}

\hfill ರಾಮಾಯಣ ಬಾಲಕಾಂಡ ೪-೧

ಅತೀಂದ್ರಿಯಾರ್ಥದರ್ಶಿಯಾಗಿ ತತ್ತ್ವಜ್ಞನೆನಿಸಿದ ವಾಲ್ಮೀಕಿಯು ಆತ್ಮವಂತನಾಗಿದ್ದು, ರಾಮನು ರಾಜ್ಯಾಧಿಪತ್ಯವನ್ನು ಹೊಂದಿದ ನಂತರ ಆಶ್ಚರ್ಯಕರವಾದ ಪದಗಳಿಂದ ಕೂಡಿದ ಅವನ ಚರಿತ್ರಗ್ರಂಥವನ್ನು ರಚಿಸಿದನು.

\medskip
\noindent\textbf{೧೬೦. ಬಹವೋ ದುರ್ಲಭಾಶ್ಚೈವ} $\star$\hfill ಪುಟ \pageref{163a},\pageref{194a},\pageref{219a},\pageref{247c}

\hfill ರಾಮಾಯಣ ಬಾಲಕಾಂಡ ೪-೧

ಯಾವಗುಣಗಳನ್ನು ನೀನು ಕೀರ್ತನೆ ಮಾಡಿರುವೆಯೋ ಅವು ಬಹುವೂ ಮತ್ತು ದುರ್ಲಭಗಳೂ ಆಗಿರುತ್ತವೆ. ಮುನಿಯೇ! ನಾನೂ ಸಹ ಮುನಿಯೇ ಆದ್ದರಿಂದ ಒಳಗೇ ಬೋಧದಿಂದ ಅರಿತು ಹೇಳುವೆನು. ಅಷ್ಟೂ ಗುಣಗಳುಳ್ಳ ನರನನ್ನು ಆಲಿಸು.

\medskip
\noindent\textbf{೧೬೧. ಬಹೂನಿ ಮೇ ವ್ಯತೀತಾನಿ} $\star$ \hfill ಪುಟ \pageref{91a}

\hfill ಭಗವದ್ಗೀತಾ ೪-೫

\medskip
\noindent\textbf{೧೬೨. ಬಾಷ್ಪಪರ್ಯಾಕುಲೇಕ್ಷಣಾಃ} \hfill ಪುಟ \pageref{211b}

\hfill ರಾಮಾಯಣ ಬಾಲಕಾಂಡ ೪-೧೫

`ತ್ವಚ್ಛೃತ್ವಾ' ಇತ್ಯಾದಿ ಗ್ರಂಥದ ಭಾವಾನುವಾದವನ್ನು ಗಮನಿಸಿ.

\medskip
\noindent\textbf{೧೬೩. ಬಿಂಬಾದಿವೋದ್ಧೃತೌ ಬಿಂಬೌ} \hfill ಪುಟ \pageref{167b},\pageref{228a}

ಲವಕುಶರು ರಾಮದೇಹರೂಪವಾದ ಬಿಂಬದಿಂದ ಹೊರಬಂದ ಇನ್ನೆರಡು ಪ್ರತಿಬಿಂಬಗಳೋ ಎಂಬಂತಿದ್ದರು.

\medskip
\noindent\textbf{೧೬೪. ಬುಭೂಷಯಾಸಕ್ತಮನಾ ಬಭೂವ ಹ} \hfill ಪುಟ \pageref{211c}

`ಸ ಚಾಪಿ' ಇತ್ಯಾದಿ ಗ್ರಂಥದ ಭಾವಾನುವಾದವನ್ನು ಗಮನಿಸಿ.

\medskip
\noindent\textbf{೧೬೫. ಬ್ರಹ್ಮ ಸತ್ಯಂ ಜಗನ್ಮೀಥ್ಯಾ} \hfill ಪುಟ \pageref{54a}

\hfill `ಜೀವೋ ಬ್ರಹ್ಮೈವ ನಾಪರಃ' ಎನ್ನುವುದು ಇದರ ಉತ್ತರ ವಾಕ್ಯ. ಬ್ರಹ್ಮವು - (ಪರಮಾತ್ಮ ತತ್ತ್ವ) ಪರಮಾರ್ಥವಾದುದು. ಜಗತ್ತೆಂಬುದು ಪರಮಾರ್ಥವಲ್ಲದ್ದೂ. ತಾತ್ಕಾಲಿಕ ಸತ್ಯವಷ್ಟೆ. ಜೀವತತ್ತ್ವವು ಬ್ರಹ್ಮದ ಅಂಶವೇ ಆದ್ದರಿಂದ (ಮಮೈವಾಂಶಃ ಜೀವ ಭೂತಃ) ಬ್ರಹ್ಮವೇ. ಬ್ರಹ್ಮಕ್ಕಿಂತಲೂ ಅತ್ಯಂತ ಭಿನ್ನವೇನೂ ಅಲ್ಲ.

\medskip
\noindent\textbf{೧೬೬. ಭಗವಾನ್! ಪುಂಡರೀಕಾಕ್ಷ} \hfill ಪುಟ \pageref{115a}

\hfill ವೈಷ್ಣವಪೂಜಾಕಲ್ಪ

ಕಮಲನೇತ್ರನಾದ ಭಗವಂತನೇ! ಈ ವರೆಗೆ ಮಾಡಿದ ಆಂತರವಾದ ನಿನ್ನ ಪೂಜೆಯನ್ನು ಸ್ವೀಕರಿಸು. ಸ್ವಾಮಿ! ದೇವಾಧಿದೇವನೇ! ಹೊರಗೂ ಇಂದ್ರಿಯಗಳ ಮೂಲಕವೂ ನಿನ್ನನ್ನು ಚೆನ್ನಾಗಿ ಅರ್ಚಿಸತೊಡಗುವೆನು.

\medskip
\noindent\textbf{೧೬೭. ಭಜ ಯತಿರಾಜಂ ಭಜ ಯತಿರಾಜಂ} \hfill ಪುಟ \pageref{46}

\hfill ರಾಮಾನುಜಸ್ತೋತ್ರ

ಸಂಸಾರಕ್ಕಾಗಿ ಹೆದರುವ ಜೀವಿಯೇ! ಯತಿರಾಜನಾದ ರಾಮಾನುಜನನ್ನು ಸೇವಿಸು.

ಕಾಂತಿಮತೀ ಮತ್ತು ಕೇಶವಸೋಮಯಾಜೀ ಇವರುಗಳ ಸುಕುಮಾರಪುತ್ರನೂ ಜೀವಿಗಳ ಮೇಲೆ ಔದಾರ್ಯದಿಂದ ಕೂಡಿದವನೂ, ಆದಿಶೇಷನ ಅವತಾರರೂಪಿಯೂ, ಮೂಗರೂ ಕಿವುಡರೂ ಮೊದಲಾಗಿ ದೇಹ ವೈಕಲ್ಯವುಳ್ಳ ಜೀವಿಗಳಿಗೂ ಮೋಕ್ಷಕೊಡಬಲ್ಲ ಯತಿರಾಜರಾಮಾನುಜನನ್ನು ಭಜಿಸು.

ಗುಣ ಮತ್ತು ಗುಣಿ ಇವುಗಳ ಭೇದವು ಶಾಶ್ವತವಾದುದು. ಚೇತನ - ಅಚೇತನ ಮತ್ತು ಈಶ್ವರ ಇವರುಗಳ ಭಿನ್ನತೆಯೂ ಸತ್ಯವಾದುದು. ಚೇತನಾಚೇತನಗಳೆರಡನ್ನೂ ಒಡಲಾಗಿ ಹೊಂದಿರುವ ಹರಿಯೇ ಈಶ್ವರನು. ಇದು ಸತ್ಯವಾದ ವಿಶಿಷ್ಟಾದ್ವೈತಸಿದ್ಧಾಂತವಾಗಿದೆ.

ಸತ್ಯಾರ್ಥವನ್ನು ಪರರಿಗೆ ತಿಳಿಸುವುದರಲ್ಲಿ ಶಕ್ತನಾಗಬೇಕಾದರೆ ನೈಸರ್ಗಿಕ ಪ್ರಜ್ಞಾವಂತನಾಗಬೇಕು, ಅಂತೆಯೇ ಪರರ ಹಿತದಲ್ಲಿ ಆಸಕ್ತಿಯಿರಬೇಕು ಮತ್ತು ತನ್ನ ಖ್ಯಾತಿ - ಲಾಭ - ಪೂಜೆಗಳಲ್ಲಿ ಆಸೆಯಿಲ್ಲದ ವಿರಕ್ತನಾಗಿರಬೇಕು.

\medskip
\noindent\textbf{೧೬೮. ಭಾರತಃ ಪಂಚಮೋ ವೇದಃ} \hfill ಪುಟ \pageref{122}

ಭಾರತವೆಂಬುದು ಐದನೆಯ ವೇದವಾಗಿದೆ.

\medskip
\noindent\textbf{೧೬೯. ಭೂತದಯಾಂ ವಿಸ್ತಾರಯ} \hfill ಪುಟ \pageref{115b}

\hfill ಶ್ರೀಶಂಕರರಭಗವತ್ಪಾದ ಕೃತ ಷಟ್ಪದೀ ಸ್ತೋತ್ರ

ಓ ವಿಷ್ಣುವೇ! ನನ್ನಲ್ಲಿ ಬೆಳೆದಿರುವ ಅವಿನಯವನ್ನು ಕಳೆ. ಮನಸ್ಸು ಶಮನಗೊಳ್ಳುವಂತೆ ಮಾಡು. ಶಬ್ದಾದಿ ವಿಷಯಗಳು ಸುಖಕೊಡುವುದೆಂಬ ಬಿಸಿಲು - ಗುದುರೆಯನ್ನು ದೂರಮಾಡು. ಭೂತಕೋಟಿಗಳಲ್ಲಿ ಕರುಣೆಯ ವಿಸ್ತಾರಗೊಳ್ಳುವಂತೆ ಕರುಣಿಸು. ಸಂಸಾರವೆಂಬ ಸಮುದ್ರವನ್ನು ದಾಟಿಸಿ ಅತ್ತ ಕಡೆಯ ಅಮೃತತ್ತ್ವವನ್ನು ದಯಪಾಲಿಸು.

\medskip
\noindent\textbf{೧೭೦. ಭ್ರಾತರೌ ಸ್ವರಸಂಪನ್ನೌ} \hfill ಪುಟ \pageref{250b}

\hfill ರಾಮಾಯಣ ಬಾಲಕಾಂಡ ೪-೫

`ಕೃತ್ವಾಪಿ ತನ್ಮಹಾಪ್ರಾಜ್ಞಃ' ಇತ್ಯಾದಿ ಗ್ರಂಥದ ಭಾವಾನುವಾದವನ್ನು ಗಮನಿಸಿ.

\medskip
\noindent\textbf{೧೭೧. ಮಚ್ಛಂದಾದೇವ ತೇ ಬ್ರಹ್ಮನ್} $\star$\hfill ಪುಟ \pageref{21a},\pageref{203b},\pageref{250c}

\hfill ರಾಮಾಯಣ ಬಾಲಕಾಂಡ ೨-೩೧

\medskip
\noindent\textbf{೧೭೨. ಮನನಾತ್ ತ್ರಾಣನಾಚ್ಚೈವ} $\star$\hfill ಪುಟ \pageref{10}

\hfill ಯೋಗಶಿಖೋಪನಿಷತ್

\medskip
\noindent\textbf{೧೭೩. ಮನಸಾ ಚಿಂತಯಾಮಾಸ} \hfill ಪುಟ \pageref{192}

\hfill ರಾಮಾಯಣ ಸುಂದರಕಾಂಡ ೪೯-೧೬,೧೭

ಆಂಜನೇಯನು ರಾವಣನು ದೇಹದಲ್ಲಿ ಎದ್ದು ಕಾಣುವ ಆವನ ತೇಜಸ್ಸಿನಿಂದ ಮೋಹಗೊಂಡವನಾಗಿ ಒಳಗೊಳಗೇ ಆಲೋಚಿಸಿದನು `ಏನಿದು? ಆಶ್ಚರ್ಯಕರವಾಗಿದೆ! ರಾಕ್ಷಸರಾಜನಾದ ಇವನಲ್ಲಿ ಈ ಬಗೆಯ ಸರ್ವಲಕ್ಷಣಸಂಪನ್ನತೆ? ಎಂತಹಾ ಆಶ್ಚರ್ಯಕರವಾದ ರೂಪ? ಎಂತಹ ದಿಟ್ಟತನ? ಎಂತಹಾ ಮನೋಬಲ? ಎಂತಹಾ ಕಾಂತಿ?'

\medskip
\noindent\textbf{೧೭೪. ಮನಸ್ಯನ್ಯತ್ ವಚಸ್ಯನ್ಯತ್} \hfill ಪುಟ \pageref{102}

\begin{shloka}
`ಮನಸ್ಯನ್ಯತ್ ವಚಸ್ಯನ್ಯತ್ ಕರ್ಮಣ್ಯನ್ಯದ್ದುರಾತ್ಮನಾಮ್ |\\
ಮನಸ್ಯೇಕಂ ವಚಸ್ಯೇಕಂ ಕರ್ಮಣ್ಯೇಕಂ ಮಹಾತ್ಮನಾಮ್ ||'
\end{shloka}

ದುರಾತ್ಮರಾದವರು ಮನದಲ್ಲೊಂದು ಯೋಚಿಸುತ್ತಿದ್ದರೆ, ಬಾಯಲ್ಲಿ ಮತ್ತೊಂದನ್ನಾಡುತ್ತಿದ್ದು, ದೇಹತಃ ಇವೆರಡಕ್ಕಿಂತಲು ಬೇರೆಯೇ ಮಾಡುವರು. ಮಹಾತ್ಮರು ಮನಸ್ಸು ವಾಕ್ ಕಾಯಗಳೆಂಬ ತ್ರಿಕರಣಗಳಲ್ಲೂ ಏಕರೂಪವಾದ ವೃತ್ತಿಯನ್ನವಲಂಬಿಸಿ ಏಕರೂಪರಾಗಿಯೇ ಇರುವರು.

\medskip
\noindent\textbf{೧೭೫. ಮಮ ದ್ವಿತಾ ರಾಷ್ಟ್ರಂ} $\star$ \hfill ಪುಟ \pageref{148a}

\medskip
\noindent\textbf{೧೭೬. ಮಮಾಪಿ ತದ್ಭೂತಿಕರಂ} \hfill ಪುಟ \pageref{159a},\pageref{168b},\pageref{193a},\pageref{211d}

\hfill ರಾಮಾಯಣ ಬಾಲಕಾಂಡ ೪-೩೫

ಈ ಕುಶಲವರೀರ್ವರೂ ಪಾರ್ಥಿವಲಕ್ಷಣದಿಂದ ಕೂಡಿರುವ ಮುನಿಗಳೂ, ಮಹತ್ತರವಾದ ತಫೋಬಲಹೊಂದು ಗಾಯಕರೂ ಆಗಿದ್ದಾರೆ. ಈ ಮಹಾನುಭಾವಯುಕ್ತವಾದ ಚರಿತವು ಕೇವಲ ನಿಮಗೆ ಮಾತ್ರವಲ್ಲದೇ ಕಥಾನಾಯಕನಾದ ನನಗೂ ನನ್ನ ಪೂರ್ವೋತ್ತರಗಳನ್ನು ನೆನಪಿಸಿ ಸ್ವಸ್ಥಿತಿಗೆ ಏರಿಸುವ ಮೂಲಕ ಮಂಗಲಕರವಾಗಿದೆ. ಮನಗೊಟ್ಟು ಕೇಳಿ ತಿಳಿಯಿರಿ.

\medskip
\noindent\textbf{೧೭೭. ಮಯೈವ ಪ್ರೋಚ್ಯಮಾನಸ್ಯ} \hfill ಪುಟ \pageref{119c}

\hfill ಮಹಾಭಾರತ ಆದಿಪರ್ವ ೧-೧೧೩

ಓ ಗಣೇಶನೇ! ಮನಸ್ಸಿನಿಂಡ ಸಿದ್ಧಪಡಿಸಿಕೊಂಡು ನಾನೇ ಸ್ಪಷ್ಟವಾಗಿ ಹೇಳುವ ಈ ಭಾರತಗ್ರಂಥಕ್ಕೆ ನೀನೇ ಬರಹಗಾರನಾಗು.

\medskip
\noindent\textbf{೧೭೮. ಮಹರ್ಷೇ ತ್ವಂ ಸಮರ್ಥೋಽಸಿ} \hfill ಪುಟ \pageref{246a}

\hfill ಬಾಲಕಾಂಡ ೧-೫, ೬

ವಾಲ್ಮೀಕಿಯು ನಾರದ ಮುನಿಗೆ ಹೇಳಿದ್ದೇನೆಂದರೆ `ಓ ಮುನಿಯೇ! ನಾನು ಬಯಸಿದಂತಹ ಸಮಸ್ತಗುಣಗಳಿಂದಲೂ ಪೂರ್ಣವಾದ ಈ ಬಗೆಯ ನರನೋರ್ವನನ್ನು ನೀನು ತಿಳಿಯಲು ಸಮರ್ಥನಾಗಿರುವೆ. ಏಕೆಂದರೆ ತ್ರಿಕಾಲಜ್ಞನಾದ ಮಹರ್ಷಿಯು ನೀನಲ್ಲವೇ' ಎಂದು. ಈ ಮಾತನ್ನು ಕೇಳಿಕೊಂಡು, ತ್ರಿಲೋಕವೇದಿಯಾದ ನಾರದ ಮುನಿಯು ಪ್ರಹೃಷ್ಟಾಂತಃಕರಣನಾಗಿ ಓ ಮುನಿಯೇ, ಆಲಿಸಪ್ಪಾ ನಿನ್ನ ಪ್ರಶ್ನೆಗೆ ಉತ್ತರ ವೀಯುವೆನೆಂದು ಮುಂದಿನ ಮಾತನ್ನಾಡಿದನು.

\medskip
\noindent\textbf{೧೭೯. ಮಾಕುರು ಧನಜನಯೌವನಗರ್ವಂ} \hfill ಪುಟ \pageref{60}

\hfill ಭಜಗೋವಿಂದಸ್ತೋತ್ರಂ 

ಎಲೈ ಜೀವಿಯೇ! ತಾತ್ಕಾಲಿಕವಾಗಿ ದೊರಕಿಸುವ ಹಣ, ಜನ, ಪ್ರಾಯವಯಸ್ಸು ಇವುಗಳಿಂದ ಹೆಮ್ಮೆಗೊಳ್ಳಬೇಡ.

\medskip
\noindent\textbf{೧೮೦. ಮಾ ನಿಷಾದ ಪ್ರತಿಷ್ಠಾಂ ತ್ವಮಗಮಃ} $\star$ \hfill ಪುಟ \pageref{20a},\pageref{157b},\pageref{166},\pageref{200}

\hfill ರಾಮಾಯಣ ಬಾಲಕಾಂಡ ೨-೧೫

\medskip
\noindent\textbf{೧೮೧. ಮುರಾರೆಸ್ಮೃತೀಯಃ ಪಂಥಾಃ} \hfill ಪುಟ \pageref{43}

ಮುರಾರಿಕವಿಯ ಕಾವ್ಯರಚನಾಮಾರ್ಗವು ಉಳಿದ ಕವಿಗಳ ಮಾರ್ಗಕ್ಕಿಂತ ಭಿನ್ನವಾದುದು.

\medskip
\noindent\textbf{೧೮೨. ಮೂಢ ಜಹೀಹಿ ಧನಾಗಮತೃಷ್ಣಾಂ} \hfill ಪುಟ \pageref{55a}

ಮೂಢನೇ! ಹಣವು ಯಾವಾಗಲೂ ಬರುತ್ತಿರಲೇಬೇಕೆಂಬ ಧನದಾಹವನ್ನು ದೂರಮಾಡು.

\medskip
\noindent\textbf{೧೮೩. ಯಃ ಕರ್ಣಾಂಜಲಿಸಂಪುಟೈಃ} $\star$  \hfill ಪುಟ \pageref{152},\pageref{168a},\pageref{181b},\pageref{208a},\pageref{211e},\pageref{244a}

\hfill ರಾಮಾಯಣ ಧ್ಯಾನಶ್ಲೋಕ 

\medskip
\noindent\textbf{೧೮೪. ಯಜ್ಞೋ ದಾನಂ ತಪಶ್ಚೈವ} \hfill ಪುಟ \pageref{172}

\hfill ಭಗವದ್ಗೀತಾ ೧೮-೫

ಮನಸ್ಸಿನ ಮೇಲೆ ಪ್ರಭುತ್ವವನ್ನು ಹೊಂದಿ ಮನಿಷಿಗಳಾಗಬೇಕೆಂದು ಬಯಸುವವರಿಗೆ, ಮನಸ್ಸಿನಲ್ಲಿ ಹುಟ್ಟಬಹುದಾದ ಮಾಲಿನ್ಯಗಳನ್ನು ತೊಳೆದು ಪಾವನಗೊಳಿಸುವಂತವು ಯಜ್ಞ, ದಾನ, ತಪಸ್ಸು ಇವುಗಳು.

\medskip
\noindent\textbf{೧೮೫. ಯತ್ರ ಯೋಗೇಶ್ವರಃ ಕೃಷ್ಣಃ} $\star$ \hfill ಪುಟ \pageref{92a}

\hfill ಭಗವದ್ಗೀತಾ ೧೮-೭೮

\medskip
\noindent\textbf{೧೮೬. ಯಥಾ ರಾಜಾ ತಥಾ ಪ್ರಜಾಃ} \hfill ಪುಟ \pageref{225c}

ಆಳುವ ಪ್ರಭುವು ಹೇಗಿರುವನೋ ಹಾಗೆಯೇ ಅವನ ಪ್ರಜೆಗಳೂ ಆಗುವರು.

\medskip
\noindent\textbf{೧೮೭. ಯಥಾವತ್ಪೂಜಿತಸ್ತೇನ} \hfill ಪುಟ \pageref{197}

\hfill ರಾಮಾಯಣ ಬಾಲಕಾಂಡ ೨-೨

ದೇವರ್ಷಿಯಾದ ನಾರದನು ವಾಲ್ಮೀಕಿಯಿಂದ ಯಥಾಯೋಗ್ಯ ಪೂಜಿಸಲ್ಪಟ್ಟು, ಅವನಿಂದ ಬೀಳ್ಕೊಂಡು - ಅನುಮತಿಯನ್ನು ಪಡೆದು, ಅಂತರಿಕ್ಷಕ್ಕೆ ಪ್ರಯಾಣ ಬೆಳೆಸಿದನು. ನಾರದನು ದೇವಲೋಕಕ್ಕೆ  ತೆರಳಲಾಗಿ, ವಾಲ್ಮೀಕಿಯು ಮುಹೂರ್ತಕಾಲದೊಳಗೆ ಗಂಗೆಗೆ ಅನತಿದೂರದಲ್ಲಿದ್ದ ತಮಸಾತೀರಕ್ಕೆ ತೆರಳಿದನು.

\medskip
\noindent\textbf{೧೮೮. ಯದಾ ಧರ್ಮಗ್ಲಾನಿರ್ಭವತಿ} $\star$\hfill ಪುಟ \pageref{150b}

\hfill ಶಂಕರಭಗವತ್ಪಾದರ ಕೃಷ್ಣಾಷ್ಟಕ

\medskip
\noindent\textbf{೧೮೯. ಯದಾ ಯದಾ ಹಿ ಧರ್ಮಸ್ಯ} \hfill ಪುಟ \pageref{226a}

\hfill ಭಗವದ್ಗೀತಾ ೪-೭

ಎಲೈ ಅರ್ಜುನನೇ! ಯಾವ ಯಾವ ಕಾಲಪರಿಧಿಯಲ್ಲಿ ಧರ್ಮವು ಕಾಹಯಿಸುವುದೋ, ಆ ಧರ್ಮವು ತಲೆ  ಎತ್ತುವುದೋ ಆಯಾ ಸಮಯದಲ್ಲಿ ನನ್ನನ್ನು ನಾನು (ಶ್ರೀ ಕೃಷ್ಣನು) ಸೃಷ್ಟಿಸಿಕೊಳ್ಳುತ್ತೇನೆ - ವ್ಯಕ್ತನಾಗುತ್ತೇನೆ.

\medskip
\noindent\textbf{೧೯೦. ಯದ್ಯಧರ್ಮೋ ನ ಬಲವಾನ್} \hfill ಪುಟ \pageref{192b},\pageref{215}

\hfill ರಾಮಾಯಣ ಸುಂದರಕಾಂಡ ೪೯-೧೮

ಇವನಲ್ಲಿ ಅಧರ್ಮವೆಂಬುದು ಬಲಿಷ್ಠವಾಗಿ ಇಲ್ಲದೇ ಇದ್ದ ಪಕ್ಷದಲ್ಲಿ ಇವನು ರಾಕ್ಷಸೇಶ್ವರನಾಗಿರುವುದೂ ಸರಿಯೇ, ಅಂತೆಯೇ ಇಂದ್ರನಿಂದ ಕೂಡಿದ ಸಂಪೂರ್ಣ ದೇವಲೋಕಕ್ಕೂ ಒಡೆಯನಾಗುವುದೂ ಸರಿಯೇ.

\medskip
\noindent\textbf{೧೯೧. ಯಸ್ತಾದೃಶಂ ಚಾರುರವಂ} \hfill ಪುಟ \pageref{250}

\hfill ರಾಮಾಯಣ ಬಾಲಕಾಂಡ ೨-೨೯

`ಪಾಪಾತ್ಮನಾ' ಇತ್ಯಾದಿ ಗ್ರಂಥದ ಭಾವಾನುವಾದವನ್ನು ಗಮನಿಸಿ.

\medskip
\noindent\textbf{೧೯೨. ಯಾ ನಿಶಾ ಸರ್ವಭೂತಾನಾಂ} \hfill ಪುಟ \pageref{138}

\hfill ಭಗವದ್ಗೀತಾ ೨-೭೯

ಪ್ರಾಕೃತರೂ ಪಾಮರರೂ ಆದ ಜೀವರಾಶಿಗಳೆಲ್ಲಾ ಯಾವ ತತ್ತ್ವ ಭೂಮಿಯಲ್ಲಿ ಕಣ್ಣು ತೆರೆಯಲಾಗದೇ ಅವರ ಪಾಲಿಗೆ ರಾತ್ರಿಯಾಗುವುದೋ, ಅಲ್ಲಿ ಸಂಯಮದಿಂದ ತತ್ತ್ವಜ್ಞಾನಪಡೆದು ಅಪ್ರಾಕೃತನೂ ಪಂಡಿತನೂ ಆದವನಿಗೆ ಎಚ್ಚರದ ಅರಿವಿನ ಸ್ಥಿತಿ ಇರುತ್ತದೆ. ಭೂತಕೋಟಿಗಳೆಲ್ಲಾ ಯಾವ ಭೂಮಿಯಲ್ಲಿ ಎಚ್ಚರಗೊಂಡು ಅರಿವು ಹೊಂದಿರುವುವೋ, ಆ ಪ್ರಾಕೃತಭೂಮಿಯು, ಅಂಗೈನೆಲ್ಲಿಯಂತೆ ಸ್ಫುಟವಾಗಿ ಪರತತ್ತ್ವವನ್ನು ಸಾಕ್ಷಾತ್ಕರಿಸಬಲ್ಲವನಿಗೆ ದೃಷ್ಟಿಪ್ರಸಾರಕ್ಕೆ ವಿಷಯವಾಗದೇ ರಾತ್ರಿಯಂತೆ ಕಗ್ಗತ್ತಲೆಯಾಗಿರುವುದು.

\medskip
\noindent\textbf{೧೯೨. ಯಾವತ್ ಸ್ಥಾಸ್ಯಂತಿ ಗಿರಯಃ} \hfill ಪುಟ \pageref{158b},\pageref{166},\pageref{208b},\pageref{210b}

\hfill ರಾಮಾಯಣ ಬಾಲಕಾಂಡ ೨-೩೬

ಓ ಪ್ರಾಚೇತಸನೇ! ಈ ಧರೆಯಲ್ಲಿ ಎಲ್ಲಿಯವರೆಗೆ ಗಿರಿ-ನದೀ-ಸಮುದ್ರಾದಿಗಳಿರುವುವೋ ಆ ಕಾಲದವರೆಗೆ ನೀನು ರಚಿಸಿದ ರಾಮಾಯಣ ಕಥೆಯು ದೇಹಿಗಳ - ಜೀವಗಳ ನರನಾಡಿಗಳಲ್ಲಿ ಸಂಚರಿಸುತ್ತಾ ಬೆಳಗುತ್ತದೆಯಪ್ಪಾ! (ಇದು ಬ್ರಹ್ಮ ದೇವನ ವರ)

\medskip
\noindent\textbf{೧೯೪. ಯಾವಾನರ್ಥ ಉದಪಾನೇ} $\star$ \hfill ಪುಟ \pageref{14}

\hfill ಭಗವದ್ಗೀತಾ ೨-೪೬

\medskip
\noindent\textbf{೧೯೫. ಯೋಗಮಾಸ್ಥಿತಃ} \hfill ಪುಟ \pageref{249b}

\hfill `ತತಃ ಪಶ್ಯತಿ' ಇತ್ಯಾದಿ ಗ್ರಂಥದ ಭಾವಾನುವಾದವನ್ನು ಗಮನಿಸಿ.

\medskip
\noindent\textbf{೧೯೬. ಯೋ ವೇದ ನಿಹಿತಂ ಗುಹಾಯಾಂ} \hfill ಪುಟ \pageref{108}

\hfill ತೈತ್ತಿರಿಯೋಪನಿಷತ್ ೧-೧

ಯಾವನು ಪ್ರಮವ್ಯೋಮವೆನಿಸುವ ಅತ್ಯಂತ ನಿಗೂಢ ಸ್ಠಾನದಲ್ಲಿ ನೆಲೆಸಿರುವ ಸತ್ಯಜ್ಞಾನಾನಂಟಸ್ವರೂಪವಾದ ಪುರುಷರ್ಥಗಳೆಲ್ಲವನ್ನೂ ಪರಸ್ಪರ ವಿರೋಧವಿಲ್ಲದೇ ಬ್ರಹ್ಮಾನುಭವಾತ್ಮಕವಾಗಿಯೇ ಅನುಭವಿಸುವಂತಾಗುವನು.

\medskip
\noindent\textbf{೧೯೭. ರಮಣೀಯಂ ಪ್ರನ್ನಾಂಬು} \hfill ಪುಟ \pageref{165c},\pageref{211f}

\hfill ರಾಮಾಯಣ ಬಾಲಕಾಂಡ ೨-೫

`ಅಕರ್ದಮಮಿದಂ' ಇತ್ಯಾದಿ ಗ್ರಂಥದ ಭಾವಾನುವಾದವನ್ನು ಗಮನಿಸಿ.

\medskip
\noindent\textbf{೧೯೮. ರಸೈಃ ಶೃಂಗಾರಕಾರುಣ್ಯ} \hfill ಪುಟ \pageref{188}

\hfill ರಾಮಾಯಣ ಬಾಲಕಾಂಡ ೪-೯

ಶೃಂಗಾರ, ಕರುಣಾ, ಹಾಸ್ಯ, ವೀರ, ಭಯಾನಕ, ರೌದ್ರ, ಬಿಭತ್ಸ, ಅದ್ಭುತ, ಶಾಂತಗಳೆಂಬ ನವರಸಗಳಿಂದ ತುಂಬಿ ತುಳುಕುವ ಈ ರಾಮಾಯಣಕಾವ್ಯವನ್ನು ಕುಶಲವರು ಹಾಡಿದರು.

\medskip
\noindent\textbf{೧೯೯. ರಸೋ ವೈ ಸಃ} \hfill ಪುಟ \pageref{76c}

\hfill ತೈತ್ತಿರೀಯೋಪನಿಷತ್ ೨-೭

ರಸವು ಆ ಬ್ರಹ್ಮವಸ್ತುವೇ ಆಗಿದೆ. ಆ ರಸವನ್ನು ಹೊಂದಿದಾಗಲೇ ಈ ಜೀವಿಯು ನಿರುಪಾಧಿಕಾನಂದಭಾಗಿಯಾಗುವನು.

\medskip
\noindent\textbf{೨೦೦. ರಾಮಸ್ಯ ಚರಿತಂ ಸರ್ವಂ} \hfill ಪುಟ \pageref{204c}

\hfill ರಾಮಾಯಣ ಬಾಲಕಾಂಡ ೨-೩೨

`ಶ್ಲೋಕ ಏವ' ಇತ್ಯಾದಿ ಗ್ರಂಥದ ಭಾವಾನುವಾದವನ್ನು ಗಮನಿಸಿ.

\medskip
\noindent\textbf{೨೦೧. ರಾಮೋ ರಾಜ್ಯಮುಪಾಸಿತ್ವಾ} \hfill ಪುಟ \pageref{253a}

\hfill ರಾಮಾಯಣ ಬಾಲಕಾಂಡ ೧-೯೭

ರಾಮಚಂದ್ರನು ಹಿಂದೆ ಹೇಳಿದ ಪ್ರಕಾರ ರಾಜ್ಯನಿರ್ವಹಣೆ ಮಾಡಿ ನಂತರ ಸ್ವಸ್ವರೂಪವೆನಿಸುವ ಬ್ರಹ್ಮಧಾಮಕ್ಕೆ ತೆರಳುವನು.

\medskip
\noindent\textbf{೨೦೨. ರಾಮೋ ವಿಗ್ರಹವಾನ್ ಧರ್ಮಃ} \hfill ಪುಟ \pageref{247}

\hfill ರಾಮಾಯಣ ಕಿಷ್ಕಿಂಧಾಕಾಂಡ

ರಾಮನೆಂದು ನಾವು ಯಾರನ್ನು ನೋಡುತ್ತಿರುತ್ತೇವೆಯೋ ಅವನು ಮೂರ್ತಿಯನ್ನು ತಳೆದ ಧರ್ಮವೇ ಆಗಿರುವನು. ಧರ್ಮವು ವಿಗ್ರಹವನ್ನು ಸ್ವೀಕರಿಸಿದಾಗ ಅದೇ `ರಾಮನು'

\medskip
\noindent\textbf{೨೦೩. ರಾವಣೇನ ರಣೇ ಶಕ್ತಿಃ} \hfill ಪುಟ \pageref{120}

\hfill ರಾಮಾಯಣ ಯುದ್ಧಕಾಂಡ ೧೦೧-೩೪

ಯುದ್ಧಾಂಗಣದಲ್ಲಿ ಕೋಪಗೊಂಡ ರಾವಣನಿಂದ ಮಹಾಸರ್ಪದಂತಿದ್ದ ಶಕ್ತಿಯು ನಿರ್ಭೀತನಾದ ಲಕ್ಷ್ಮಣನ ವಕ್ಷಃಸ್ಥಲದಲ್ಲಿಗೆ ಬಿಡಲ್ಪಟ್ಟು (ಪ್ರಯೋಗಿಸಲ್ಪಟ್ಟು) ಅಲ್ಲಿಯೇ ಮುಳುಗಿ ಮರೆಯಾಯಿತು.

\medskip
\noindent\textbf{೨೦೪. ಲಕ್ಷ್ಮೀನಾಥಸಮಾರಂಭಾಂ} \hfill ಪುಟ \pageref{48b}

\hfill ಶ್ರೀಪತಿಯಿಂದ ಮೊದಲಾಗಿ ನಾಥಮುನಿಗಳು ಮತ್ತು ಯಾಮುನಾಚಾರ್ಯರುಗಳನ್ನು ಮಧ್ಯದಲ್ಲುಳ್ಳದ್ದೂ ನನ್ನ ಆಚಾರ್ಯನನ್ನೆ ಕೊನೆಯಾಳುಳ್ಳದ್ದೂ ಆದ ಗುರುಪರಂಪರೆಯನ್ನು ನಮಿಸುತ್ತೇನೆ.

\medskip
\noindent\textbf{೨೦೫. ಲೋಕಸ್ಸಮಸ್ತಾಸ್ಸುಖಿನೋ ಭವಂತು} \hfill ಪುಟ \pageref{193b}

\hfill ಸಮಸ್ತ ಜನರೂ ಸುಖಿಗಳಾಗಿರಲಿ.

\medskip
\noindent\textbf{೨೦೬. ವಂದೇ ಗುರೂಣಾಂ ಚರಣಾರವಿಂದೇ} $\star$ \hfill ಪುಟ \pageref{180a}

\hfill ಯೋಗತಾರಾವಲಿ ೧

ಯಾವ ಗುರುಪಾದಾರವಿಂದಗಳು ಜನನಮರಣಚಕ್ರಕ್ಕೆ  ಸಿಕ್ಕಿದ ಜೀವಿಗಳ ಸಂಸಾರವೆಂಬ ಹಾಲಾಹಲವಿಷಸೇವನೆಯಿಂದುಂಟಾದ ಮೋಹವನ್ನು ಶಾಂತಿಗೊಳಿಸಲು ವಿಷವೈದ್ಯನಂತೆ ಇರುತ್ತವೆಯೋ, ಆತ್ಮಾನಂದದ ಅರಿವನ್ನು ತಂದುಕೊಡುವ ಅಂತಹ ಗುರುಗಳ ಅಡಿದಾವರೆಗಳನ್ನು ವಂದಿಸುತ್ತೇನೆ.

\medskip
\noindent\textbf{೨೦೭. ವಂದೇ ಪೂರಿತಚಂದ್ರಮಂಡಲಗತಂ} $\star$ \hfill ಪುಟ \pageref{32a}

\hfill ಹಯಗ್ರೀವೋಪನಿಷತ್

\medskip
\noindent\textbf{೨೦೮. ವಸವೋ ವಾಸವಂ ಯಥಾ} \hfill ಪುಟ \pageref{250e}

\hfill `ಅಭ್ಯಷಿಂಚನ್' ಇತ್ಯಾದಿ ಗ್ರಂಥದ ಭಾವಾನುವಾದವನ್ನು ಗಮನಿಸಿ.

\medskip
\noindent\textbf{೨೦೯. ವಾಗಾರ್ಥಾವಿದ ಸಂಪೃಕ್ತೌ} $\star$ \hfill ಪುಟ \pageref{21b},\pageref{31b},\pageref{223a},\pageref{???},\pageref{238}

\hfill ರಘುವಂಶ ೧-೧

\medskip
\noindent\textbf{೨೧೦. ವಾಚಂ ಚೋದಾಹರಿಷ್ಯಾಮಿ} $\star$ \hfill ಪುಟ \pageref{232}

\hfill ರಾಮಾಯಣ ಸುಂದರಕಾಂಡ ೩೦-೧೭, ೧೮

\medskip
\noindent\textbf{೨೧೧. ವಾಚೋವಿಧೇಯಂ ತತ್ಸರ್ವಂ} \hfill ಪುಟ \pageref{189a}

\hfill ರಾಮಾಯಣ ಬಾಲಕಾಂಡ ೪-೧೨

ರಾಮಾಯಣ ಕಾವ್ಯಾರ್ಥವೆಲ್ಲಾ ಮನೋವಿಧೇಯವಾಗಿರುವಂತೆಯೇ ಕಾವ್ಯ ರೂಪವಾದ ಶಬ್ದವನ್ನೆಲ್ಲಾ ವಾಚೋವಿಧೇಯ ಮಾಡಿಕೊಂಡು ಹಾಡಿದರು.

\medskip
\noindent\textbf{೨೧೨. ವಾತಾಪಿ ಗಣಪತಿಂ ಭಜೇಽಹಂ} \hfill ಪುಟ \pageref{45}

\hfill ವಾತಾಪಿ ಕ್ಷೇತ್ರದಲ್ಲಿರುವ ಗಣಪತಿಯನ್ನು ನಾನು ಭಜಿಸುತ್ತೇನೆ.

\medskip
\noindent\textbf{೨೧೩. ವಾಲ್ಮೀಕಿಗಿರಿಸಂಭೂತಾ} $\star$ \hfill ಪುಟ \pageref{151a},\pageref{166b},\pageref{174}

\medskip
\noindent\textbf{೨೧೪. ವಾಲ್ಮೀಕಿರಥ ತಂ ದೃಷ್ಟ್ವಾ} \hfill ಪುಟ \pageref{202a}

\hfill ರಾಮಾಯಣ ಬಾಲಕಾಂಡ ೨-೨೪

ಚತುರ್ಮುಖನಾದ ಆ ಬ್ರಹ್ಮನು ಅವನ ಆಶ್ರಮಕ್ಕೆ ಬಂದ ನಂತರ ವಾಲ್ಮೀಕಿಮುನಿಯು ಆ ಬ್ರಹ್ಮನನ್ನು ನೋಡಿ ಥಟಕ್ಕನೆ ಮೇಲೆದ್ದು ಬೆರಗಾಗಿ ಮೌನಹೊಂದಿ  ತಲೆಬಾಗಿ ನಮಸ್ಕರಿಸಿ ಮತ್ತೆ ಬದ್ಧಾಂಜಲಿಯಾಗಿ ಪರಮಾಶ್ಚರ್ಯ ಚಕಿತನಾಗಿ ಎದುರಿಗೆ ನಿಂತನು.

\medskip
\noindent\textbf{೨೧೫. ವಾಸಾಂತೇ ಬಾಲಮಧ್ಯೇ} \hfill ಪುಟ \pageref{29}

\hfill ದೇವೀ ಭಾಗವತ

\begin{shloka}
(ಆಧಾರೇ ಲಿಂಗ- ನಾಭೌ ಹೃದಯಸರಸಿಜೇ ತಾಲುಮೂಲೇ ಲಲಾಟೇ |\\
ದ್ವೇ ಪತ್ರೇ ಷೋಡಶಾರೇ ದ್ವಿದಶದಶದಲೇ ದ್ವಾದಶಾರ್ಧೇ ಚತುಷ್ಕೇ ||)\\
\hfill (ತಂತ್ರಶಾಸ್ತ್ರ)
\end{shloka}

ನಾಲ್ಕು ದಳಗಳುಳ್ಳ ಮೂಲಾಧಾರ ಚಕ್ರದಲ್ಲಿ ವ, ಶ, ಷ, ಸ ಗಳೆಂಬ ನಾಲ್ಕು ವರ್ಣಗಳೂ,  ಆರು ದಳಗಳುಳ್ಳ ಲಿಂಗಮೂಲಸ್ಥಾನ - ಸ್ವಾಧಿಷ್ಠಾನ ಚಕ್ರದಲ್ಲಿ ಬ, ಭ, ಮ, ಯ, ರ, ಲ ಗಳೆಂಬ ವರ್ಣಗಳಾದರೂ, ನಾಭಿಸ್ಥಾನದ ಮಣಿಪೂರ ಚಕ್ರ ಹತ್ತು ದಳಗಳುಳ್ಳದ್ದರಲ್ಲಿ - ಡ, ಢ, ಣ, ತ, ಥ, ದ ಧ, ನ, ಪ, ಫ ಗಳೆಂಬ ಹತ್ತು ವರ್ಣಗಳೂ, ಹನ್ನೆರಡು ದಳಗಳುಳ್ಳ ಹೃದಯಸ್ಥಾನದ ಅನಾಹತ ಚಕ್ರದಲ್ಲಿ ಕ, ಖ, ಗ, ಘ, ಙ, ಚ, ಛ, ಜ, ಝ, ಞ, ಟ, ಠ ಗಳೆಂಬ ಹನ್ನೆರಡು ವರ್ಣಗಳೂ, ಹದಿನಾರು ದಳಗಳುಳ್ಳ ಕಂಠಸ್ಥಾನದ ವಿಶುದ್ಧಚಕ್ರದಲ್ಲಿ ಅ, ಆ, ಇ, ಈ, ಉ, ಊ, ಋ, ೠ, ಏ, ಐ, ಓ, ಔ, ಅಂ, ಅಃ ಗಳೆಂಬ ಹದಿನಾರು ಸ್ವರಾಕ್ಷರಗಳೂ, ಎರಡು ದಳಗಳುಳ್ಳ ಭ್ರೂಮಧ್ಯಸ್ಥಾನದ ಆಜ್ಞಾಚಕ್ರದಲ್ಲಿ ಹ(ಂ) ಕ್ಷ (ಂ) ಗಳೆಂಬ ಎರಡು ವರ್ಣಗಳೂ ಆಗಿ ಪಂಚಾಶದ್ವರ್ಣಾತ್ಮಕವಾಗಿ ಐವತ್ತು ದಳಗಳಲ್ಲಿರುವ ಪರಮಾತ್ಮನನ್ನು ನಮಿಸುತ್ತೇನೆ. ಇಂತಹ ಚಕ್ರಮಾರ್ಗವನ್ನು ಎಲ್ಲಾ ಮುನಿಗಣಗಳೂ ಬ್ರಹ್ಮಮಾರ್ಗವೆಂದು ಹೇಳುವರು.

\medskip
\noindent\textbf{೨೧೬. ವಿದ್ಯಯಾ ತಾತ ಸೃಷ್ಟಾನಾಂ} \hfill ಪುಟ \pageref{133}

ಅಯ್ಯಾ! ವಿದ್ಯೆಯಿಂದ ಆವಿರ್ಭಾವಗೊಂದ ಪ್ರಜೆಗಳಿಗೆ ವಿದ್ಯೆಯೇ ಪರಮಗತಿಯಾಗಿರುವುದು.

\medskip
\noindent\textbf{೨೧೭. ವಿದ್ಯಯಾಽಮೃತಮಶ್ನುತೇ} \hfill ಪುಟ \pageref{99},\pageref{133a}

(ಅವಿದ್ಯಯಾ ಮೃತ್ಯುಂ ತೀರ್ತ್ವಾ ವಿದ್ಯಯಾಮೃತಮಶ್ನುತೇ)

\hfill ಈಶಾವಾಸ್ಯೋಪನಿಷತ್ ೧೧

ವಿದ್ಯಾವಿಲಕ್ಷಣವಾದ ಅವಿದ್ಯೆಯೆಂಬ ಆನು ಇಳಿದು ಬಂದ ಏಣಿಯ ದ್ವಾರವೇ ಮೃತ್ಯುವನ್ನು ದಾಟಿಹೋಗಿ ವಿದ್ಯೆಯ ಮೂಲಕ ಅಮೃತತ್ತ್ವವನ್ನು - ನೈಜಸ್ವರೂಪವನ್ನು ಪಡೆಯುತ್ತಾನೆ. (ಸಂಸಾರಿಜೀವನು)

\medskip
\noindent\textbf{೨೧೮. ವಿದ್ಯಾಮಯೋಽಯಂ ಪುರುಷಃ} \hfill ಪುಟ \pageref{133b}

\hfill ಮಹಾಭಾರತ ಶಾಂತಿಪರ್ವ ೨೪೧-೧೫

ಈ ಜೀವರೂಪೀ ಪುರುಷನು ಮೂಲತಃ ಕೇವಲ ವಿದ್ಯಾಮಯನಾದವನು.

\medskip
\noindent\textbf{೨೧೯. ವಿದ್ಯಾವಿನಯಸಂಪನ್ನೇ} $\star$ \hfill ಪುಟ \pageref{67}

\hfill ಭಗವದ್ಗೀತಾ ೫-೧೮

\medskip
\noindent\textbf{೨೨೦. ವಿದ್ಧಿ ಮಾಂ ಋಷಿಭಿಸ್ತುಲ್ಯಂ} \hfill ಪುಟ \pageref{151b}

\hfill ರಾಮಾಯಣ ಅಯೋಧ್ಯಕಾಂಡ ೧೯-೨೦

\begin{shloka}
ನಾಽಹಮರ್ಥಪರೋ ದೇವಿ! ಲೋಕಮಾವಸ್ತುಮುತ್ಸಹೇ |\\
ವಿದ್ಧಿಮಾಮೃಷಿಭಿಸ್ತುಲ್ಯಂ ಕೇವಲಂ ಧರ್ಮಮಾಸ್ಥಿತಮ್ ||
\end{shloka}

ದೇವಿ ಕೈಕೇಯಿ! ನಾನು (ರಾಮ) ಅರ್ಥಕಾಮಪರನಾಗಿ ಲೋಕದಲ್ಲಿ ಬಾಳಲು ಬಯಸುವುದಿಲ್ಲವಮ್ಮ. ಕೇವಲ ಧರ್ಮವನ್ನೇ ಅವಲಂಬಿಸಿ ಬಾಳುವಂತನಾಗಿದ್ದು ಋಷಿಗಳಿಗೆ ಸಮಾನನಾಗಿರತಕ್ಕವನೆಂದು ತಿಳಿದುಕೋಮ್ಮ!.

\medskip
\noindent\textbf{೨೨೧. ವೇದಃ ಪ್ರಾಚೇತಸಾದಾಸೀತ್} \hfill ಪುಟ \pageref{228},\pageref{232c}

`ವೇದವೇದ್ಯೇ' ಇತ್ಯಾದಿ ಶ್ಲೋಕಗಳ ಅರ್ಥಾನುವಾದದಲ್ಲಿ ಗಮನಿಸಿ.

\medskip
\noindent\textbf{೨೨೨. ವೇದವೇದ್ಯೇ ಪರೇ ಪುಂಸಿ} \hfill ಪುಟ \pageref{168c}

ವೇದವೇದ್ಯನಾದ ಪರಮಪುರುಷನು ದಶರಥ ಪುತ್ರನಾಗಿ ಜನಿಸಲಾಗಿ ವೇದವು ಸಾಕ್ಷಾತ್ ರಾಮಾಯನ ರೂಪವಾಗಿ ಪ್ರಾಚೇತಸನಾದ ವಾಲ್ಮೀಕಿಯಿಂದ ರೂಪುಗೊಂಡಿತು.

\medskip
\noindent\textbf{೨೨೩. ವ್ಯಾಖ್ಯಾಮುದ್ರಾಂ ಕರಸರಸಿಜೈಃ} \hfill ಪುಟ \pageref{100}

\hfill ಹಯಗ್ರೀವಸ್ತೋತ್ರ

ವಿದ್ಯಾರಾಜನಾದ ಹಯವದನನು ತನ್ನ ಕರಕಮಲಗಳಿಂದ ವ್ಯಾಖ್ಯಾನ (ಉಪದೇಶ) ಮುದ್ರೆಯನ್ನೂ ಪುಸ್ತಕವನ್ನೂ ಶಂಖಚಕ್ರಗಳನ್ನು ಧರಿಸಿಕೊಂಡು.
 
\medskip
\noindent\textbf{೨೨೪. ಶತ್ರುನಿಬರ್ಹಣಃ} \hfill ಪುಟ \pageref{250e}

\hfill ರಾಮಾಯಣ ಬಾಲಕಾಂಡ ೧-೯

`ಇಕ್ಷ್ವಾಕುವಂಶಪ್ರಭವೋ' ಇತ್ಯಾದಿ ಗ್ರಂಥದ ಭಾವಾನುವಾದವನ್ನು ಗಮನಿಸಿ.

\medskip
\noindent\textbf{೨೨೫. ಶರೀರಮಾದ್ಯಂ ಖಲು} \hfill ಪುಟ \pageref{76d}

\hfill ಕುಮಾರಸಂಭವ ೫-೩೩

ಧರ್ಮವನ್ನು ಸಂಪಾದಿಸುವ ಕಾರ್ಯದಲ್ಲಿ ದೇಹವು ಮೊದಲ ಸಾಧನ - ಉಪಕರಣವಾಗಿದೆಯಷ್ಟೆ. 

\hfill ಅಮರವಾಣಿ - ೧೧

\medskip
\noindent\textbf{೨೨೬. ಶಿವಶಕ್ತಿಮಯಂ ಸಾಕ್ಷ್ಜಾತ್} $\star$\hfill ಪುಟ \pageref{32b}

\medskip
\noindent\textbf{೨೨೭. ಶಿಷ್ಯೋಪಾಧ್ಯಾಯಿಕಯಾ ವೃತ್ಯಾ} \hfill ಪುಟ \pageref{45a}

ಅಧ್ಯಾಪಕನ ಅಧ್ಯಾಪನ ವ್ಯಾಪಾರದಿಂದಲೂ ನಿರುಕ್ತ ಶಿಷ್ಯನ ಅಧ್ಯಯನ ವ್ಯಾಪಾರದಿಂದಲೂ ಶಾಸ್ತ್ರವು ಗ್ರಂಥರೂಪದಲ್ಲೂ ಅರ್ಥರೂಪದಲ್ಲೂ ವಿಚ್ಛೇದವಾಗದೇ ಪ್ರಾಪ್ತವಾಗುವುದಕ್ಕೆ ಸಂಪ್ರದಾಯವೆನ್ನುತ್ತಾರೆ.

\medskip
\noindent\textbf{೨೨೮. ಶೀಕ್ಷಾಂ ವ್ಯಾಖ್ಯಾಸಾಮಃ | ವರ್ಣ ಸ್ವರಃ |} $\star$\hfill ಪುಟ \pageref{10a}

\hfill ತೈತ್ತಿರಿಯೋಪನಿಷತ್ ಶೀಕ್ಷಾವಲ್ಲಿ ೧

ವೇದೋಚ್ಚಾರಣೆಗೆ ಸಹಕಾರಿಯಾಗಿ ವೇದಾರ್ಥಾನುಚಿಂತನೆಗೆ ಕಾರಣವಾಗಿ ವೇದಾಂಗವಾಗಿರುವ ಶಿಕ್ಷಾಶಾಸ್ತ್ರವನ್ನು ವಿವರಿಸುವೆವು.

\medskip
\noindent\textbf{೨೨೯. ಶುಚಿರ್ವಶ್ಯಃ ಸಮಾಧಿಮಾನ್} \hfill ಪುಟ \pageref{255}

\hfill ರಾಮಾಯಣ ಬಾಲಕಾಂಡ ೧-೧೨

`ಇಕ್ಷ್ವಾಕುವಂಶಪ್ರಭವೋ' ಇತ್ಯಾದಿ ಗ್ರಂಥದ ಕನ್ನಡ ಭಾವಾನುವಾದವನ್ನು ಗಮನಿಸಿ.

\medskip
\noindent\textbf{೨೩೦. ಶೃಂಗಾರಂ ಕ್ಷಿತಿನಂದನಾವಿಹರಣೇ} \hfill ಪುಟ \pageref{182a},\pageref{208c},\pageref{219b}

ಭೂಮಿಸುತೆಯಾದ ಸೀತಾದೇವಿಯೊಡನೆ ವಿಹರಿಸುವಲ್ಲಿ ಶೃಂಗಾರ ರಸಮೂರ್ತಿಯೂ, ಶಿವಧನುಸ್ಸನ್ನು ಮುರಿಯುವ ವೇಳೆಗೆ ವೀರರಸಮೂರ್ತಿಯೂ, ಕಾಕಾಸುರನ ಬಗ್ಗೆ ಕರುಣಾರಸಮೂರ್ತಿಯೂ, ಸಮುದ್ರದಲ್ಲಿ ಸೇತುವೆಯನ್ನು ಕಟ್ಟಲೋಸುಗ ಗಿರಿ ಪರ್ವತಗಳನ್ನು ಇರಿಸುವಾಗ್ಗೆ ಅದ್ಬುತರಸಮೂರ್ತಿಯೂ, ಶೂರ್ಪನಖಿಯಲ್ಲಿ ಸಂಭಾಷಣೆ ಮಾಡುವಾಗ ಹಾಸ್ಯರಸಮೂರ್ತಿಯೂ, ಭಯತರುವಂತಹ ಯುದ್ಧಭೂಮಿಯಲ್ಲಿ ಭಯಾನಕ - ಬೀಭತ್ಸ ರಸಮೂರ್ತಿಯೂ, ರಾವಣನನ್ನು ಸಂಹರಿಸುವಾಗ್ಗೆ ರೌದ್ರರಸಮೂರ್ತಿಯೂ, ತಪೋಮಯರಾದ ಋಷಿಜನರಲ್ಲಿ ಶಾಂತರಸಮೂರ್ತಿಯೂ, ಆಗಿರುವ ಶ್ರೀರಾಮ ದಿವ್ಯಮೂರ್ತಿಯು ನಮ್ಮನ್ನು ರಕ್ಷಿಸಲಿ.

\medskip
\noindent\textbf{೨೩೧. ಶೋಚನ್ನೇವ ಮುಹುಃ ಕ್ರೌಂಚೀಂ} \hfill ಪುಟ \pageref{203c}

\hfill ರಾಮಾಯಣ ಬಾಲಕಾಂಡ ೨-೨೯

ಪುರುಷ ಕ್ರೌಂಚವನ್ನಗಲಿ ನೋಯುತ್ತಿರುವ ಹೆಣ್ಣು ಕ್ರೌಂಚದ ತೊಳಲಾಟವನ್ನೇ ನೆನೆದು ಶೋಕಪಡುತ್ತ ಅಂತರ್ಮುಖವಾದ ಮನಹೊಂದಿ, ಶೋಕದಿಂದ ಪಾರಾಗಲಾಗದ ಸ್ಥಿತಿಯಲ್ಲಿರುತ್ತಾ ಇದೇ ಶ್ಲೋಕವನ್ನೇ (ಪಾಪಾತ್ಮನಾ ಕೃತಂ) ಮತ್ತೆ ಮತ್ತೆ ಹಾಡಿದನು.

%raghu, pdf 336
\medskip
\noindent\textbf{೨೩೨. ಶ್ರುತ್ವಾ ಚೈತತ್ತ್ರಿಲೋಕಜ್ಞಃ} \hfill ಪುಟ \pageref{153},\pageref{163},\pageref{194},\pageref{246},\pageref{251}

\hfill ರಾಮಾಯಣ ಬಾಲಕಾಂಡ ೧-೬, ೭

ಮೂರು ಲೋಕದ ವೃತ್ತಾಂತವನ್ನೂ ಅರಿತ ನಾರದಮುನಿಯು ವಾಲ್ಮೀಕಿಯ ಈ ಮಾತನ್ನು ಕೇಳಿ, `ಕೇಳಪ್ಪ!' ಎಂದು ಹೇಳಿ ಸಂತುಷ್ಟಾಂತಃಕರಣನಾಗಿ ಮುಂದಿನ ಮಾತನ್ನಾಡಿದನು.

ಯಾವ ಗುಣಗಳನ್ನು ನೀನು ಕೀರ್ತನೆ ಮಾಡಿರುವೆಯೋ ಅವು ಬಹುವೂ ಮತ್ತು ದುರ್ಲಭಗಳೂ ಆಗಿರುತ್ತವೆ. ಮುನಿಯೇ! ನಾನೂ ಸಹ ಮುನಿಯೇ ಆದ್ದರಿಂದ ಒಳಗೇ ಬೋಧದಿಂದ ಅರಿತು ಹೇಳುವೆನು. ಅಷ್ಟೂಗುಣಗಳುಳ್ಳ ನರನನ್ನು ಆಲಿಸು!

\medskip
\noindent\textbf{೨೩೩. ಶ್ಲೋಕ ಏವ ತ್ವಯಾ ಬದ್ಧಃ} \hfill ಪುಟ \pageref{157},\pageref{181}

\hfill ರಾಮಾಯಣ ಬಾಲಕಾಂಡ ೨-೩೧

ಅಯ್ಯಾ! ಮುನಿಯೇ! ನಿನ್ನಿಂದ ಹೊರಬಂದುದು ಶ್ಲೋಕಬಂಧವೇ ಆಗಿದೆ. ಈ ಬಗ್ಗೆ ಸಂದೇಹಪಡಬೇಕಾಗಿಲ್ಲ. ಓ! ಬ್ರಹ್ಮರ್ಷಿಯೇ! ನಿನ್ನ ಈ ವಾಕ್ಕು  ನನ್ನ ಅಪೇಕ್ಷೇಯಂತೆಯೇ ಪ್ರವೃತ್ತವಾಗಿದೆ. (ಅಂದರೆ ವಿಧಿನಿಯಮವನ್ನನುಸರಿಸಿ ಬಂದುದೇ ಹೊರತು ಸ್ವಚ್ಛಂದವಾದುದಲ್ಲ.)

ಋಷಿಶ್ರೇಷ್ಠನೇ! ಸಮಗ್ರವಾದ ರಾಮಚರಿತ್ರೆಯನ್ನು ರಚನೆ ಮಾಡುವವನಾಗು ಧರ್ಮಾತ್ಮನೂ ಪ್ರಶಸ್ತವಾದ ಗುಣಗಳುಳ್ಳವನೂ, ಧೀಮಂತನೂ, ಮಹಾವೀರನೂ ಆದ ರಾಮನ ವೃತ್ತಾಂತವನ್ನು ನಾರದನಿಂದ ಕೇಳಿ ತಿಳಿದುಕೊಂಡಂತೆ ಕಥನ ಮಾಡುವವನಾಗು. 

\medskip
\noindent\textbf{೨೩೪. ಶ್ಲೋಕಸಾರಸಮಾಕೀರ್ಣಂ} $\star$ \hfill ಪುಟ \pageref{152}

\medskip
\noindent\textbf{೨೩೫. ಷಡ್ವಿಂಶಃ ಪರಮಾತ್ಮಾ} \hfill ಪುಟ \pageref{185}

(ತತ್ತ್ವಗಳ ಎಣಿಕೆಯಲ್ಲಿ ಪ್ರಾಕೃತ ತತ್ತ್ವಗಳು ಇಪ್ಪತ್ತನಾಲ್ಕು. ಇಪ್ಪತ್ತೈದನೆಯವನು ಜೀವನು.) ಇಪ್ಪತ್ತಾರನೆಯವನು ಈ ಜಡ ಮತ್ತು ಚೇತನರನ್ನು ಮೀರಿದ ಪರಮಾತ್ಮನಾಗಿರುವನು. 

\medskip
\noindent\textbf{೨೩೬. ಸ ಕ್ಷಯ ಏಹಿ} \hfill ಪುಟ \pageref{162}

\begin{shloka}
ಯಜ್ಞೋಭೂತ್ವಾ ಯಜ್ಞಮಾಸೀದ ಸ್ವಾಂ ಯೋನಿಮ್ |\\
ಜಾತವೇದೋ ಭುವ ಆಜಾಯಮಾನಸ್ಸಕ್ಷಯ ಏಹಿ ||
\end{shloka} 

ದೇವತಾರಾಧನರೂಪವಾದ ಯಜ್ಞವನ್ನು ನಡೆಸಲು ಯಜ್ಞರೂಪಿಯಾದ ಭಗವಂತನೇ ಬರಬೇಕು. ನಮಗೆ ಬೇಕಾದುದನ್ನು ಸ್ವೀಕರಿಸಲು ನಮ್ಮ ಕೈಯನ್ನೋ ನಮ್ಮ ಬಾಯನ್ನೋ ತಾನೇ ನೀಡಬೇಕು. `ಅಗ್ನಿರವಮೋದೇವತಾನಾಂ ವಿಷ್ಣುಃ ಪರಮಃ' ಎಂಬಂತೆ ಅವನೇ ಅವಮವಾದ ಅಗ್ನಿರೂಪ ತಾಳಿ ಕೆಳಗಿಳಿಯಬೇಕು. ಅವನೇ ಪರಮವಾದ ವಿಷ್ಣುರೂಪದೊಡನೇ ಅಗ್ನಿರೂಪದಿಂದ ತಂದ ಹವಿಸ್ಸನ್ನು ಸ್ವೀಕರಿ ಭೋಕ್ತೃವಾಗಬೇಕು. `ಹವ್ಯಕವ್ಯಭುಗೇಕಸ್ತ್ವಂ ಪಿತೃದೇವಸ್ವರೂಪಧೃತ್' ಎನ್ನುವಂತೆ ದೈವ ಮತ್ತು ಪಿತ್ರ್ಯವಾದ ಅಹುತಿಗಳನ್ನೆಲ್ಲಾ ಕೈಗೊಳ್ಳಬೇಕು.

\begin{shloka}
`ಏಕೋ ವಿಷ್ಣುರ್ಮಹದ್ಭೂತಾನ್ ಪೃಥಗ್ಭೂತಾನ್ಯನೇಕಶಃ |\\
ತ್ರೀನ್‌ಲೋಕಾನ್ವ್ಯಾಪ್ಯಭೂತಾತ್ಮಾ ಭುಂಕ್ತೇ ವಿಶ್ವಭುಗವ್ಯಯಃ ||'
\end{shloka}
ಎಂಬಂತೆ ಹವ್ಯವೂ ಹವ್ಯವಹನೂ ಹವ್ಯಭೋಜಿಯೂ ಹವ್ಯಪ್ರದನೂ ಎಲ್ಲವೂ ಒಬ್ಬನೇ ಆಗಬೇಕು. ಆದ್ದರಿಂದ ಓ ಜಾತವೇದನೇ! - ನೀನು ಯಜ್ಞಮೂರ್ತಿಯೇ ಆಗಿ ಯಜ್ಞಮಯವಾದ ನಿನ್ನ ನಿರುಪಾಧಿಕ ಮೂಲಸಾಧನವನ್ನು ಬಂದು ತಲುಪಬೇಕು. ಭುವಿಯಿಂದ ದಿವಿಗೇರುತ್ತಾ ನಿನ್ನ ಯಜ್ಞವೇದಿಕೆಯಾದ ದೈವವೂ ಧ್ರುವವೂ ಸ್ಥಾನದೊಡನೇ ಬಾಪ್ಪ (ಕ್ಷಿ-ನಿವಾಸೇ) ಕ್ಷಯವೆಂದರೆ ವಾಸಸ್ಥಾನ. ಲೋಕ ಪ್ರಯೋಗದಲ್ಲಿ ವಾಸಸ್ಥಾನವೆಂಬ ಅರ್ಥದಲ್ಲಿ ಕ್ಷಯ ಶಬ್ದ ಪ್ರಯೋಗವು ದುರ್ಲಭವಾಗಿದೆ. (ಶ್ರೀ ರಂಗಗುರುವು ದಯಪಾಲಿಸಿದ ಭಾವದ ವಿವರಣೆ)

\medskip
\noindent\textbf{೨೩೭. ಸ ಚಾಪಿ ರಾಮಃ ಪರಿಷದ್ಗತಃ} \hfill ಪುಟ \pageref{168}

\hfill ರಾಮಾಯಣ ಬಾಲಕಾಂಡ ೪-೩೬

ಆ ಶ್ರೀರಾಮಚಂದ್ರನೂ ಸಹ ಸಭೆಯಲ್ಲಿದ್ದವನಾಗಿ (ಲವಕುಶರ ರಾಮಾಯಣ ಗಾನವನ್ನು ಕೇಳಿ) ಮೈಮರೆತನು.

\medskip
\noindent\textbf{೨೩೮. ಸ ಜಗಾಮ ವಿಹಾಯಸಂ} \hfill ಪುಟ \pageref{163}

\hfill ರಾಮಾಯಣ ಬಾಲಕಾಂಡ ೨-೨

`ಯಥಾವತ್ಪೂಜಿತಃ' ಇತ್ಯಾದಿ ಗ್ರಂಥದ ಭಾವಾನುವಾದವನ್ನು ಗಮನಿಸಿ.

\medskip
\noindent\textbf{೨೩೯. ಸ ತು ಮೇಧಾವಿನೌ} \hfill ಪುಟ \pageref{239},\pageref{240}

\hfill ರಾಮಾಯಣ ಬಾಲಕಾಂಡ ೪-೬

`ಕೃತ್ವಾಪಿ' ಇತ್ಯಾದಿ ಗ್ರಂಧದ ಭಾವಾನುವಾದವನ್ನು ಗಮನಿಸಿ.

\medskip
\noindent\textbf{೨೪೦. ಸತ್ಯಂ ಸತ್ಯಂ ಪುನಸ್ಸತ್ಯಂ} \hfill ಪುಟ \pageref{130}

\hfill ವಿಷ್ಣು ಸಹಸ್ರನಾಮ ಫಲಶ್ರುತಿ

ತೋಳೆತ್ತಿ ಮತ್ತೆ ಮತ್ತೆ ಹೇಳುವ ಸತ್ಯವಿದಾಗಿದೆ.

\medskip
\noindent\textbf{೨೪೧. ನತ್ಯಮೇವೋದ್ಧರಾಮ್ಯಹಮ್} \hfill ಪುಟ \pageref{40}

ನಾನು ಸತ್ಯವನ್ನೇ ಎತ್ತಿಹಿಡಿಹುತ್ತೇನೆ. (ಪಟ್ಟಾಭಿಷಿಕ್ತ ರಾಜನೂ ಗುರುವಾಗತಕ್ಕ ಸಂನ್ಯಾಸಿಯೂ ಮಾಡುವ ಪ್ರತಿಜ್ಞೆ.)

\medskip
\noindent\textbf{೨೪೨. ಸನ್ಮನುಷ್ಯಮನೋ ಯಥಾ} $\star$\hfill ಪುಟ \pageref{197}

\hfill ರಾಮಾಯಣ ಬಾಲಕಾಂಡ ೨-೫

`ಅಕರ್ದಮಮಿದಂ' ಇತ್ಯಾದಿ ಗ್ರಂಥದ ಭಾವಾನುವಾದವನ್ನು ಗಮನಿಸಿ.

\medskip
\noindent\textbf{೨೪೪. ಸ ಪೂರ್ವೇಷಾಮಪಿ ಗುರುಃ} \hfill ಪುಟ \pageref{129}

\hfill ಪಾತಂಜಲ ಯೋಗಸೂತ್ರ ೧-೨೬

ಈಶ್ವರನು ಲೋಕಗುರುಗಳೆನಿಸಿ ಜ್ಞಾನಪ್ರದರಾದ ಬ್ರಹ್ಮಾದಿಗಳಿಗೂ ಗುರುವಾಗಿದ್ದಾನೆ. ಏಕೆಂದರೆ ಅವನು ಕಾಲಪರಿಮಿತಿಗೆ ಒಳಪಡದವನು.

\medskip
\noindent\textbf{೨೪೫. ಸಪ್ತಮೇ ಮಾಸಿ ಜೀವಸಂಯುಕ್ತೋ ಭವತಿ} $\star$\hfill ಪುಟ \pageref{85}

\hfill ಗರ್ಭೋಪನಿಷತ್

\medskip
\noindent\textbf{೨೪೬. ಸ ಪ್ರವಿಶ್ಯಶ್ರಮಪದಂ} \hfill ಪುಟ \pageref{201}

\hfill ರಾಮಾಯಣ ಬಾಲಕಾಂಡ ೨-೨೨

ಧರ್ಮತತ್ತ್ವಜ್ಞನಾದ ಆ ಮುನಿಯು ಶಿಷ್ಯನೊಡನೆ ಆಶ್ರಮಕ್ಕೆ ಪ್ರವೇಶಿಸಿ, ಅಲ್ಲಿ ಕುಳಿತು ಧ್ಯಾನವನ್ನು ಹೊಂದಿ, ನಂತರ ಇನ್ನೂ ಕೆಲವು ಬೇರೆ ಬೇರೆ ಮಾತುಗಳನ್ನೂ ಆಡಿದನು.

\medskip
\noindent\textbf{೨೪೭. ಸ ಬ್ರಹ್ಮಾ ಸ ಶಿವಃ ಸೇಂದ್ರಃ} \hfill ಪುಟ \pageref{249}

\hfill ಮಹಾನಾರಾಯಣೋಪನಿಷತ್

ಅವನೇ ಸೃಷ್ಟಿಕರ್ತನಾದ ಬ್ರಹ್ಮಾ. ಅವನೇ ಸಂಹಾರಕಾರೀ ಶಿವನು. ಅವನೇ ಇಂದ್ರಿಯಾಧೀಶ್ವರನು. ಅವನೇ ಅಕ್ಷರರೂಪಿಯಾದವನು. ಅವನೇ ಪರಮಸ್ವಾರಾಜ್ಯವಾದ ಪರಮಪದಕ್ಕೆ ಒಡೆಯನು. 

\medskip
\noindent\textbf{೨೪೮. ಸರ್ಪಾಣಾಂ ವಾಸುಕಿಶ್ಚಾಽಸ್ಮಿ} \hfill ಪುಟ \pageref{93},\pageref{95},\pageref{98}

\hfill ಭಗವದ್ಗೀತಾ ೧೦-೨೮

ಸರ್ಪಗಳಲ್ಲೆಲ್ಲಾ ಶ್ರೇಷ್ಠನಾದ ವಾಸುಕಿಯು ನಾನೇ

\medskip
\noindent\textbf{೨೪೯. ಸರ್ವಂ ಕರ್ಮಾಽಖಿಲಂ ಪಾರ್ಥ} \hfill ಪುಟ \pageref{98},\pageref{78}

\hfill ಭಗವದ್ಗೀತಾ ೪-೨೨

ಜ್ಞಾನೇಚ್ಛಾ ಕ್ರಿಯಾ ಎಂಬ ಕ್ರಮದಲ್ಲಿ ಪ್ರವೃತ್ತವಾಗುವ ಕರ್ಮವೆಲ್ಲವೂ ತಾನು ಯಾವ ಜ್ಞಾನಮೂಲದಿಂದ ವಿಸ್ತಾರಹೊಂದಿತೋ ಆ ಜ್ಞಾನದಲ್ಲಿ ಪರಿಸಮಾಪ್ತಿಯನ್ನು ಹೊಂದುವುದು.

\medskip
\noindent\textbf{೨೫೦. ಸರ್ವದಾಭಿಗತಃ ಸದ್ಭಿಃ} \hfill ಪುಟ \pageref{247}

\hfill ರಾಮಾಯಣ ಬಾಲಕಾಂಡ ೧-೧೬

`ಇಕ್ಷ್ವಾಕುವಂಶಪ್ರಭವೋ' ಇತ್ಯಾದಿ ಗ್ರಂಥದ ಭಾವಾನುವಾದವನ್ನು ಗಮನಿಸಿ.

\medskip
\noindent\textbf{೨೫೧. ಸಲಿಲೇನ ಸಹಸ್ರಾಕ್ಷಂ} \hfill ಪುಟ \pageref{248}

\hfill ರಾಮಾಯಣ ಬಾಲಕಾಂಡ ೧೩೧-೬೧

`ಅಭ್ಯಷಿಂಚನ್' ಇತ್ಯಾದಿ ಗ್ರಂಥದ ಭಾವಾನುವಾದವನ್ನು ಗಮನಿಸಿ.

\medskip
\noindent\textbf{೨೫೨. ಸಹ ವೀರ್ಯಂ ಕರವಾವಹೈ} \hfill ಪುಟ \pageref{130}

\hfill ಶಾಂತಿಪಾಠ

ನಮ್ಮೀರ್ವರ ಅಧ್ಯಯನವು ವೀರ್ಯವತ್ತರವಾಗಿ ಬೆಳೆಯಲಿ. ನಮ್ಮಿಬ್ಬರ ಅಧ್ಯಯನವು ತೇಜಃಪೂರ್ಣವಾಗಿ ಬೆಳಗಲಿ.

\medskip
\noindent\textbf{೨೫೩. ಸಹಸ್ರಶಿರ್ಷಾ ಪುರುಷಃ} \hfill ಪುಟ \pageref{108},\pageref{150},\pageref{218},\pageref{249}

\hfill ತೈತ್ತಿರೀಯಾರಣ್ಯಕ ೩-೧೩

ಪಿಂಡಾಂಡ ಬ್ರಹ್ಮಾಂಡರೂಪವಾದ ಪುರಗಳಲ್ಲಿ ತುಂಬಿಕೊಂಡು ಬಾಳುವ ಪರಮಾತ್ಮರೂಪೀ ಪುರುಷನು, ಅಸಂಖ್ಯಾತವಾದ ಶಿರಸ್ಸುಗಳನ್ನೂ ಕಣ್ಗಳನ್ನೂ ಹೊಂದಿರುವನು. ಜೀವಿಗಳೆಲ್ಲರ ತಲೆಗಳೆಲ್ಲವೂ ಅವನದೇ, ಕಣ್ಗಳೂ ಅವನದೇ, ಪಾದಗಳೂ ಅವನವೇ, ಎಲ್ಲರ ಕರ್ಮಜ್ಞಾನೇಂದ್ರಿಯಗಳೂ ಅವನು ಕೊಟ್ಟ ಧರ್ಮದಿಂದ ತಂತಮ್ಮ ಕಾರ್ಯ ನಿರ್ವಹಿಸುತ್ತಿವೆ. ಅವನು ಈ ಕಾರ್ಯಪ್ರಪಂಚವನ್ನೇಲ್ಲಾ ಜ್ಞಾಪಿಸಿಕೊಂಡು ಅಲ್ಲಿ ಮಾತ್ರ ಉಳಿಯದೇ, ದಶಾಂಗುಲ ಪ್ರಮಾಣದಲ್ಲಿರುವ ದಹರಾಕಾಶವನ್ನೂ ಮೀರಿ ನಿಂತಿರುವನು. ಭೂತ, ಭವಿಷ್ಯತ್, ವರ್ತಮಾನ ಕಾಲಗಳಲ್ಲಿಯ ಏನೇನುಂಟೋ ಅದೆಲ್ಲವೂ ಪುರುಷನೇ ಆಗಿದೆ. ಆ ಪುರುಷನು ಈ ತ್ರಿಕಾಲವನ್ನು ಮೀರಿದ ಅಮೃತತ್ವವೆಂನುಂಟೋ ಅದಕ್ಕೂ ಪ್ರಭುವಾಗಿರುವನು, ಯಾವ ಅನ್ನರಸಮಯನಾದ ಜೀವಿಯು ಕಾಲಾತೀತನಾಗಿ ಊರ್ಧ್ವಕ್ಕೆ ಸಾಗಿದಾಗ ಯಾವ ಅಮೃತತ್ವವನ್ನು ಹೊಂದುವನೋ. ಈ ಪುರುಷನ ಮಹಿಮೆಯು ಇಷ್ಟರಮಟ್ಟಿನದು. ಆದ್ದರಿಂದ ಅವನು ಮಹಾಪುರುಷನಾಗಿರುವನು. ಅಥವಾ ಅವನಿಗಿಂತಲೂ ಹಿರಿಯವನಾದ (ಉತ್ತಮ) ಪುರುಷನಾಗಿರುವನು. ಸಮಸ್ತ ಪ್ರಕೃತಿ - ಪ್ರಾಕೃತ ಪ್ರಪಂಚವು ಇವನ ಒಂದು ಕಾಲುಭಾಗವಷ್ಟೇ. ಅಂದರೆ ಒಂದು ಅಂಶಮಾತ್ರವಾಗಿರುವುದು. ಈ ಚತುಷ್ಪಾದ್ ಬ್ರಹ್ಮರೂಪೀ ಪುರುಷನ ಉಳಿದ ಮೂರು ಪಾದವು ದಿವಿಯಲ್ಲಿ ಅಮೃತಸ್ವರೂಪವಾಗಿರುವುದು. ಈ ಪುರುಷನು ಊರ್ಧ್ವವಾದ ದಿವಿಯಲ್ಲಿ ಪಾದತ್ರಯವಂತನಾಗಿ ಜ್ಞಾನಿಗಳಿಗೆ ಗೋಚರಿಸಿದನು. ಒಂದು ಪಾದವು ಈ ಪ್ರಾಕೃತಭೂಮಿಯಲ್ಲಿ ಮತ್ತೆ ಮತೇ ಹುಟ್ಟಿ ಸಾಯುವ ಪ್ರಪಂಚಕ್ಕೆ ಆಶ್ರಯವಾಗಿ ನಿಂತಿತು. ಆದ್ದರಿಂದಲೇ ಆಹಾರ ಜೀವಿ ಮತ್ತು ಅನಾಹಾರಜೀವಿಗಳಾದ ಜಂಗಮಸ್ಥಾವರಾತ್ಮಕವಾಗತಕ್ಕ ಜೀವಿಗಳ ಕಡೆಗೆ ಸರ್ವತೋಮುಖವಾಗಿ ಪಾದನ್ಯಾಸ ಮಾಡಿಕೊಂಡನು. ಪ್ರಾಕೃತಪ್ರಪಂಚದ ಕಡೆಗೆ ಪಾದಬೆಳೆಸಿದ ಪುರುಷತತ್ತ್ವದಿಂದ `ವಿರಾಟ್' ಎಂಬ ಹೆಸರು ಪಡೆಯತಕ್ಕ ಅಂಡವೊಂದು (ಬ್ರಹ್ಮಾಂಡ) ಆವಿರ್ಭವಿಸಿತು.

\medskip
\noindent\textbf{೨೫೪. ಸ ಶಿಷ್ಯ ಹಸ್ತಾದಾದಾಯ} \hfill ಪುಟ \pageref{198}

\hfill ರಾಮಾಯಣ ಬಾಲಕಾಂಡ ೨-೮

ನಿಯತೇಂದ್ರಿಯನಾದ ಆ ವಾಲ್ಮೀಕಿಯು ಶಿಷ್ಯನಾದ ಭರದ್ವಾಜನ ಕೈಯಿಂದ ನಾರುಮಡಿಯನ್ನು ತೆಗೆದುಕೊಂಡು, ವಿಸ್ತಾರವಾದ ಆ ಸುತ್ತಮುತ್ತಲ ವನದ ಮೇಲೆ ಕಣ್ಣು ಹಾಯಿಸುತ್ತಾ ಸಂಚರಿಸಿದನು.

\medskip
\noindent\textbf{೨೫೫. ಸಹಿತೌ ಮಧುರಂ ರಕ್ತಂ} \hfill ಪುಟ \pageref{189}

\hfill ರಾಮಾಯಣ ಬಾಲಕಾಂಡ ೪-೧೯

ಸೋದರನಾದ ಕುಶಲವರಿರ್ವರೂ ಒಂದುಗೂಡಿ ಮಧುರವಾಗಿಯೂ ಕಾವ್ಯತತ್ತ್ವಾರ್ಥಗಳಿಂದ ರಂಜಿತವಾಗಿಯೂ ಇರುವಂತೆ (ಸ್ವರಸಂಪತ್ತಿನಿಂದ ಮೆರೆಯಿಸಿ) ಈ ರಾಮಾಯಣ ಕಾವ್ಯವನ್ನು ಹಾಡಿದರು.

\medskip
\noindent\textbf{೨೫೬. ಸಾಮಾಂ ಪಾತು ಸರಸ್ವತೀ} \hfill ಪುಟ \pageref{100}

\hfill ಜೀವನದಲ್ಲಿಯ ಜಾಡ್ಯವನ್ನು ಪೂರ್ತಿಯಾಗಿ ನಾಶಗೊಳಿಸುವ, ಜ್ಞಾನಾದಿಗುಣ ಪೂರ್ಣೆಯಾದ ಭಗವತಿಯಾದ ಸರಸ್ವತಿಯು ನನ್ನನ್ನು ಕಾಪಾಡಲಿ.

\medskip
\noindent\textbf{೨೫೭. ಸಾಧು ಸಾಧ್ವಿತಿ ತಾವೂಚುಃ} \hfill ಪುಟ \pageref{168},\pageref{211},\pageref{242}

\hfill ರಾಮಾಯಣ ಬಾಲಕಾಂಡ ೪-೧೬

(ಕುಶಲವರಿಂದ ಹಾಡಲ್ಪಟ್ಟ ರಾಮಾಯಣ ಕಾವ್ಯವನ್ನು ಶ್ರವಣಮಾಡಿ ಆನಂದ ಭಾಷ್ಪ ತುಂಬಿದ ಕಣ್ಣುಗಳುಳ್ಳ ಮುನಿಗಳೆಲ್ಲರೂ) ಅತ್ಯಾಶ್ಚರ್ಯ ಹೊಂದಿ, ಅವರಿಬ್ಬರನ್ನು `ಸಾಧು ಸಾಧು' ಎಂದು ಉದ್ಗಾರ ಮಾಡಿದರು. ಸ್ವಭಾವವಾಗಿಯೇ ಧರ್ಮವಾತ್ಸಲ್ಯವುಳ್ಳ ಆ ಮುನಿಗಳು ಪ್ರೀತಿಹೊಂದಿದ ಅಂತಃಕರಣವುಳ್ಳವರಾಗಿ ಹಾಡುತ್ತಿರುವ ಕುಶೀಲವರಾದ ಕುಶಲವರನ್ನು ಪ್ರಶಂಸಾಯೋಗ್ಯರಾದ ಕಾರಣ ಪ್ರಶಂಸೆಮಾಡಿದರು.

\medskip
\noindent\textbf{೨೫೮. ಸೋಽಭಿಷೇಕಂ ಕತಃ ಕೃತ್ವಾ} \hfill ಪುಟ \pageref{201}

\hfill ರಾಮಾಯಣ ಬಾಲಕಾಂಡ ೨-೨೦

ಆ ವಾಲ್ಮೀಕಿಮುನಿಯು ನಂತರ ಆ ತಮಸಾತಿರ್ಥದಲ್ಲಿ ವಿಧ್ಯುಕ್ತವಾದ ರೀತಿಯಲ್ಲಿ ಸ್ನಾನ ಮುಗಿಸಿ ತಾನು ಆಡಿದ ಮಾತಿನ ಅಭಿಪ್ರಾಯವನ್ನೇ ಮೆಲುಕುಹಾಕುತ್ತಾ ಆಶ್ರಮಕ್ಕೆ ಹಿಂತಿರುಗಿದನು.

\medskip
\noindent\textbf{೨೫೯. ಸೋಽಹಮಾಜನ್ಮಶುದ್ಧಾನಾಂ} $\star$\hfill ಪುಟ \pageref{148}

\hfill ರಘುವಂಶ ೧-೫

\medskip
\noindent\textbf{೨೬೦. ಸ್ವಂಚಿತಾಯತನಿಃಸ್ವನಮ್} \hfill ಪುಟ \pageref{250}

\hfill ರಾಮಾಯಣ ಬಾಲಕಾಂಡ ೪-೩೩

ರಾಮಚಂದ್ರ ಪ್ರಭುವಿನಿಂದ ಪ್ರೇರಿತರಾದ, ಸೋದರರಾದ ಆ ಕುಶಲವರೂ ಸಹ, ಮಧುರ (ಶ್ರುತಿಮನೋಹರ)ವೂ ಯಥಾಯೋಗ್ಯವಾದ ರಾಗಪುಷ್ಟಿಯುಳ್ಳದ್ದೂ, ಯೋಗ್ಯವಾದ ಗತಿ (ಆರೋಹ - ಅವರೋಹ) ಮತ್ತು ವಿಷಯಕ್ಕೆ ತಕ್ಕ ದೀರ್ಘತೆ ಇವುಗಳಿಂದ ಕೂಡಿದ ಆಲಾಪನೆಗಳುಳ್ಳದ್ದೂ, ವೀಣಾತಂತ್ರಿಯಲ್ಲಿ ನುಡಿಸಿದಾಗ ಇರತಕ್ಕಂತಹ ಲಯಸಂಪತ್ತಿನಿಂದ ಕೂಡಿರುವುದೂ, ಅರ್ಥಸ್ಫೂರ್ತಿಗೆ ಅತ್ಯಂತಾನುಗುಣವೂ, ಆಗಿರುವ ತೆರದಲ್ಲಿ ಬಹಳ ಸುಷ್ಠುವಾಗಿ ಹಾಡಿದರು.

\medskip
\noindent\textbf{೨೬೧. ಸ್ವತೋ ರಂಜಯತೀತಿ ಸ್ವರಃ} \hfill ಪುಟ \pageref{11},\pageref{17}

\hfill ಪರಾವಲಂಬನೆಯಿಲ್ಲದೇ ತಾನೂ ಪ್ರಕಾಶಿಸಿ ಮತ್ತೊಂದನ್ನೂ ಪ್ರಕಾಶಗೊಳಿಸುತ್ತೆಯಾದ್ದರಿಂದ `ಸ್ವರ' ವೆಂದು ಹೆಸರು.

\medskip
\noindent\textbf{೨೬೨. ಸ್ವಸ್ತಿ ಪ್ರಜಾಭ್ಯಃ ಪರಿಪಾಲಯಂತಾಂ} \hfill ಪುಟ \pageref{149}

\hfill ವಿಶ್ವದ ಎಲ್ಲಾ ಪ್ರಜೆಗಳಿಗೂ ಕ್ಷೇಮವಿರಲಿ! ಭೂಮಿಯನ್ನಾಳಹೊರಡುವ ಪ್ರಭುಗಳು ನ್ಯಾಯದಿಂದ ಹೊರತಾಗದ ಹಾದಿಯನ್ನವಲಂಬಿಸಿ ಭೂಮಿಯನ್ನು ಪಾಲಿಸಲಿ! ಯಜ್ಞಕ್ಕೆ ಪೋಷಣೆ ಮಾಡುವ ಗೋವುಗಳಿಗೂ ಯಜ್ಞಮಾಡಬಲ್ಲ ಬ್ರಾಹ್ಮಣರಿಗೂ ಸದಾ ಶುಭವಿರಲಿ! ಸಮಸ್ತಲೋಕಗಳೂ ನೈಜವಾದ ನೆಮ್ಮದಿಯಿಂದಿರಲಿ!

\medskip
\noindent\textbf{೨೬೩. ಸ್ಥಾನಭ್ರಷ್ಟಾ ನ ಶೋಭಂತೇ} \hfill ಪುಟ \pageref{48}

\hfill (ಮನುಷ್ಯನ) ಹಲ್ಲು ಕೂದಲು ಉಗುರು ಇವು ಅದಿರುವ ಜಾಗವನ್ನು ತೊರೆದರೆ ಅದಕ್ಕೆ ಸೊಬಗಿಲ್ಲ. ಹುಲಿಯುಗುರು, ಚಮರಿಯಕೇಶ, ಆನೆಯದಂತ ಇವು ಆಯಾಜಾಗ ಬಿಟ್ಟರೇ ಸೊಬಗು.

\medskip
\noindent\textbf{೨೬೪. ಸ್ವಾದುಮೇಧ್ಯರಸೋಪೇತಂ} \hfill ಪುಟ \pageref{119}

\hfill ಮಹಾಭಾರತ ಆದಿಪರ್ವ ೧-೧೨೦

\begin{shloka}
(ತಸ್ಯ ವೃಕ್ಷಸ್ಯ ವಕ್ಷ್ಯಾಮಿ ಶಾಖಾಪುಷ್ಪಫಲೋದಯಮ್|\\
ಪರ್ಜನ್ಯ ಇವ ಭೂತಾನಾಮಕ್ಷಯೋ ಭಾರತದ್ರುಮಃ|)
\end{shloka}

ಭೂತಕೋಟಿಗಳಿಗೆ ಮಳೆಯು ಹೇಗೋ ಹಾಗೆ ಸಮಸ್ತಜೀವಿಗಳಿಗೂ ಸಹ ಅಕ್ಷಯವಾಗಿ ಉಪಜೀವ್ಯವಾಗಿರುವುದು ಭಾರತದ್ರುಮವು. ಆ ಭಾರತವೃಕ್ಷದ ಶಾಖೆಯಲ್ಲಿ ಉದಯಿಸುವ ಪುಷ್ಪ ಮತ್ತು ಫಲಸ್ತೋಮವು ಅಸ್ವದನೀಯವೂ ಪವಿತ್ರವೂ ಆದ ರಸದಿಂದ ಪೂರ್ಣವಾಗಿರುವುದು. ಆದರೆ ದೇವತೆಗಳಿಗೂ ಇದು ಕಿತ್ತು ತಿನ್ನಲು ಸುಲಭವಾಗಿರದು.

\medskip
\noindent\textbf{೨೬೫. ಸ್ವಾಮಿನೇ ಮಮ ಮಂಗಲಂ} \hfill ಪುಟ \pageref{234}

\hfill ಶ್ರೀ ರಾಮಮಂಗಲಾಶಾಸನ

ನನ್ನ ಸ್ವಾಮಿಗೆ ಮಂಗಲವಾಗಲಿ

\medskip
\noindent\textbf{೨೬೬. ಹೃದಯಂ ತದ್ವಿಜಾನೀಯಾತ್} \hfill ಪುಟ \pageref{27},\pageref{249}

\hfill ಮಹಾನಾರಾಯಣ ಸಹಸ್ರಶಿರ್ಷಾನುವಾಕ ದ್ರಾವಿಡ ಪಾಠ

ಇದರ ಹಿಂದಿನ ವಾಕ್ಯ `ಅಧೋ ನಿಷ್ಟ್ಯಾ ವಿತಸ್ತ್ಯಾಂತು (ನ್ತೇ) ನಾಭ್ಯಾಮುಪರಿತಿಷ್ಠತಿ' ಎಂಬುದು. (ಕಮಲ ಮುಕುಲಾಕಾರವಾದ ಹೃದಯಸ್ಥಾನವಾವುದನ್ನು ಹೇಳಿದ್ದಾಯಿತೋ ಅದು) ನಿಷ್ಟ್ಯಾ (ನಿಷ್ಟಿ) ಎಂಬ ಕತ್ತಿನ ಗ್ರಂಥಿಯ ದೆಸೆಯಿಂದ ಹನ್ನೆರಡು ಬೆರಳು ಕೆಳಗಡೆಗೂ, ನಾಭಿ (ಹೊಕ್ಕಳು) ಪ್ರದೇಶದಿಂದ ಹನ್ನೆರಡು ಬೆರಳು ಮೇಲ್ಗಡೆಗೂ ಇರುವುದು. ಅದನ್ನೇ ಹೃದಯವೆಂದು ಸಾಧಕನು ಒಳಗೆ ಸಾಗುವಾಗ ನಿಶ್ಚಯಿಸಿಕೊಳ್ಳಬೇಕು. ಅದೇ ವಿಶ್ವವೆಂಬ ಸರ್ವವ್ಯಾಪೀ ಪರಮಾತ್ಮನ (ಜೀವಿಯೊಡಗೂಡಿದ) ಸರ್ವಾಧಿಕತ್ತ್ವವನ್ನು ವ್ಯಕ್ತಗೊಳಿಸುವ ಸಾಧನವಾಗಿದೆ.

\medskip
\noindent\textbf{೨೬೭. ಹೃದಯೆಷ್ವೇವ ತಿಷ್ಠಂತಿ} \hfill ಪುಟ \pageref{162}

\hfill ರಾಮಾಯಣ ಅಯೋಧ್ಯಾಕಾಂಡ ೪೫-೨೫

ನಮ್ಮ ಪರಮಧನವಾದ ವೇದಗಳು ಹೃದಯದಲ್ಲಿಯೇ ನೆಲೆಸಿವೆ.

\medskip
\noindent\textbf{೨೬೮. ಹ್ಲಾದಯತ್ಸರ್ವಗಾತ್ರಾಣಿ} \hfill ಪುಟ \pageref{189}

\hfill ರಾಮಾಯಣ ಬಾಲಕಾಂಡ ೪-೩೪

ಕುಶಲವರು ಹಾಡಿದ ರಾಮಾಯಣಗಾನವು ಸಭ್ಯರೆಲ್ಲರ ಸರ್ವಾಂಗಗಳನ್ನೂ ಮನಸ್ಸು ಮತ್ತು ಹೃದಯಗಳನ್ನೂ ಆಹ್ಲಾದಪೂರ್ಣವಾಗಿ ಮಾಡುತ್ತಾ (ರಾಜರ) ಜನರು ನೆರೆದಿದ್ದ ಸಭೆಯಲ್ಲಿ ಶ್ರೋತ್ರಕ್ಕು ಅದಕ್ಕಾಶ್ರಯನಾದ ದೇಹಿಗೂ (ಯಾರ್ಯಾರಿಗೂ ಕಿವಿಯುಂಟೋ ಅವರೆಲ್ಲರಿಗೂ) ಸುಖಕರವಾಗಿ ಬೆಳಗಿತು.

